\documentclass[11pt]{article}
\usepackage[T1]{fontenc}
\usepackage{lmodern,microtype}
\usepackage[margin=1in]{geometry}
\usepackage{amsmath,amssymb,amsthm,mathtools}
\usepackage{booktabs,array,longtable}
\usepackage{xcolor}
\usepackage[colorlinks=true,linkcolor=blue!45!black,citecolor=blue!45!black,urlcolor=blue!45!black]{hyperref}
\usepackage{bookmark}
\allowdisplaybreaks[2]
\setlength{\emergencystretch}{2em}
\numberwithin{equation}{section}
\newtheorem{theorem}{Theorem}[section]
\newtheorem{proposition}[theorem]{Proposition}
\newtheorem{lemma}[theorem]{Lemma}
\newtheorem{corollary}[theorem]{Corollary}
\theoremstyle{definition}
\newtheorem{definition}[theorem]{Definition}
\theoremstyle{remark}
\newtheorem{remark}[theorem]{Remark}
\newcommand{\e}{\mathrm e}
\newcommand{\ii}{\mathrm i}
\newcommand{\dd}{\,\mathrm d}
\newcommand{\R}{\mathbb R}
\newcommand{\C}{\mathbb C}
\newcommand{\N}{\mathbb N}
\newcommand{\Cu}{C_{\mathrm u}}
\DeclareMathOperator{\Log}{Log}
\DeclareMathOperator{\Res}{Res}
\newcommand{\cC}{\mathcal C}
\newcommand{\repo}[1]{\texttt{\detokenize{#1}}}
\title{Exterior moment hierarchies and coupled gap dynamics\\in the Newman heat flow}
\author{Edward Baker\thanks{Draft prepared for Edward Baker with OpenAI Codex LLM assistance in derivation, exposition and internal checking. The current revision's actual model and effort were unavailable in the recorded configuration; helper settings were inherited. The internal assessments are not specialist peer review.}}
\date{18 November 2026\\\small Analytical working draft}
\hypersetup{pdftitle={Exterior moment hierarchies and coupled gap dynamics in the Newman heat flow},pdfauthor={Edward Baker}}
\begin{document}
\maketitle
\begin{abstract}
We develop the dynamics of a counted zero pair and its complete exterior
divisor for the Newman heat field. Symmetric pair coordinates give an
infinite hierarchy of rational moments, regular through mergers and
algebraically equivalent to split-pole moments. Indexing those moments by
adjacent pair produces a coupled zero lattice: shared-root identities,
a common meromorphic resolvent and an exact block flux express the
compatibility between neighboring rows. These identities distinguish a
common-divisor reconstruction from independent moment truncations, whose
small residuals need not have small propagated absolute allowances.
Complex powers and damped Cauchy integrals organize the same complete
interaction, with explicit divisor and contour errors.
For backward gap comparison, we first identify the limitation of initial
interaction alone: its canonical high-height allowance is asymptotic to
$g^2/4$, compared with $g^2/6$ from two neighbors. Retaining the joint
variables $V=1-2DM$ and $K=2D^2\mathcal F$ gives an upper barrier for $V$
and a lower bound for the squared half-gap. On each fixed positive-time
slab, a growing canonical neighbor head yields relative loss
$O(\Delta/\log X)$; the source theorem already protects that exterior.
A separate analytic-residual argument recovers the published high-zero
velocity method and permits a matched direct-source comparison.
Trial-coordinate contour identities yield a local acquisition and
stopping theorem without assuming exterior reality. Forward block
dynamics give conditional force contraction and a minimum-gap tail
comparison; retaining physical forcing gives an explicit backward
Riccati comparison. An exact derivative of the initial-interaction
horizon identifies what an unchanged forward--backward return loses.
We then formulate recursive two-sided refinement through common-source
trajectory enclosures and a continuous forward-error/backward-adjoint
identity. Sound refinement is distinguished from an unproved strict
improvement of canonical bounds. The results are analytical; inherited
certificates and fixed scalar arithmetic appear only as supporting
examples. No stable finite closure, new exterior coverage or improved
de Bruijn--Newman bound is asserted.
\end{abstract}
\clearpage
\tableofcontents
\clearpage
\section{The heat field and a counted pair}
\label{sec:framework}

The local interaction of a pair of zeros is determined by the entire
exterior divisor, including nonreal zeros and their multiplicities.
The useful coordinates for this interaction are the center $c$ and the
signed discriminant $D$, rather than individually labelled target roots.
Our starting point is the pair law established in the base manuscript
and its cluster-transport development \cite{BakerBase,BakerGT41}.
The present paper extends its rational moments to every integer order,
relates them to split-pole moments, and then treats a complex exponent.
The Gaussian-resolvent predecessor \cite{BakerGR02,BakerGR05,BakerGR07}
supplies the first higher moment law and the quantitative comparison
used in Section~\ref{sec:bounds}.

The distinction between an identity and an estimate is essential to the
development. An exact representation of the exterior interaction is
valuable for organizing its evolution. A better gap estimate additionally
requires a quantitatively useful signed estimate for the relevant
coupled quantity, on a time interval whose boundary count and separation
have been established independently.
The later all-real comparison uses an additional source of sign: each
exterior denominator of an adjacent real pair is positive. This gives
an explicit local backward allowance from neighboring gaps, while its
statistical application still requires information about rare exceptions.

The development begins with exact structure. Sections~\ref{sec:integer}
and \ref{sec:split} derive the integer hierarchy and its split-pole
coordinates. Section~\ref{sec:lattice} places those rows on the same zero
lattice and proves their overlap and block-flux identities. This gives the
compatibility question in Section~\ref{sec:finite} a concrete spatial
meaning. Complex exponents, contours and signed forcing then lead to the
joint dynamics in Sections~\ref{sec:bounds}--\ref{sec:joint}.

The second stage concerns quantitative use. The backward and regional
laws of Section~\ref{sec:statistics} are followed by a joint-in-time
barrier in Section~\ref{sec:jointtime}. Section~\ref{sec:canonicaltime}
compares the initial-data and joint consumers on the same canonical
positive-time source, retaining all height and time errors. The compact
contour method in Section~\ref{sec:compact} separates source acquisition
from a noncircular stopping argument. Section~\ref{sec:forward-statistics}
then asks what finite time-zero geometry can acquire in the forward
direction. Section~\ref{sec:retained-forcing} retains the physical forcing
in the backward comparison and computes the exact cost of an unchanged
forward--backward excursion. These results motivate the common-source
two-sided refinement in Section~\ref{sec:two-sided}: a discrete enclosure
construction and a continuous residual identity, with their unacquired
strict-gain requirement stated explicitly. Supporting certificate-based
examples are deferred to Appendix~\ref{app:compact-example}; none of the
general theorems depends on those numerical values.

Strictly separating contours suffice for the local algebra, lattice and
compact stopping results. The quantitative moment-order, complex-power
and truncation estimates retain the stronger collar fixed below. Stating
these domains separately prevents an adjacent-pair or negative-time
application from silently importing a stronger geometric hypothesis.

\subsection{Canonical normalization and hypotheses}

For $u\ge0$, set
\begin{equation}
 r_k(u)=\pi k^2\e^{4u},\qquad
 \Phi(u)=\sum_{k\ge1}\e^u(2r_k(u)^2-3r_k(u))\e^{-r_k(u)},\qquad
 H(t,z)=\int_0^\infty\e^{tu^2}\Phi(u)\cos(zu)\dd u.
 \label{eq:source}
\end{equation}
This is the unchanged complete theta source, with the normalization of
\cite{BakerBase}; in particular $H(0,z)=\xi((1+\ii z)/2)/8$.
Its positive half-line source gives $H(t,0)>0$ for real $t$.
On compact sets in $(t,z)$ the integrand, including every fixed derivative,
is dominated by a constant times a polynomial in $u$ times
\[
 \exp\{9u+\Re(t)u^2+|\Im z|u-\pi\e^{4u}\}.
\]
Consequently $H$ is jointly entire, real on the real plane, even in $z$,
and satisfies the physical-time equation
\begin{equation}
 H_t=-H_{zz}.
 \label{eq:heat}
\end{equation}
The sign of time in this equation is fixed throughout the paper.
Maximizing $ru-(\pi/4)\e^{4u}$ after reserving an integrable part
of the source bound gives
$\log\max_{|z|\le r}|H(t,z)|=O(1+r\log(2+r))$
uniformly for $t$ in any fixed compact real interval. The bounded factor
$\Re(t)u^2$ is absorbed by the reserved double-exponential decay; at the
maximizing scale it contributes only $O((\log(2+r))^2)$. Thus the order
is at most one on every such interval. This does not extend any later
positive-time localization or numerical error constant to negative time.

For the quantitative strong-collar estimates,
fix a compact real interval $I=[t_0,t_1]\subset[0,1/2]$ on which an independently
established count-two tube isolates a conjugation-invariant pair. Freeze
the height and collar parameters
\begin{equation}
 \begin{gathered}
 T\ge64,\quad L=\log T,\quad \ell=L^{-1},\quad d\ge0,
 \quad \max(2,4\sqrt d)\le a\le L,\quad r_0=a\ell,\\
 c(t)\in[T,2T],\quad D(t)\in\R,\quad |D(t)|\le d\ell^2,\qquad
 z_\pm=c\pm s,\quad s^2=D,\\
 |z_\pm-c|\le r_0/4,\qquad
 w_j=\rho_j-c,\quad |w_j|\ge r_0
 \quad\text{for every exterior occurrence }\rho_j.
 \end{gathered}
 \label{eq:domain}
\end{equation}
An occurrence means a zero counted with its integer multiplicity.
The two target occurrences coincide when $D=0$; $D<0$ describes
a conjugate nonreal pair. The fixed scale $\ell$ is not differentiated.

We use the complete strip, unit-count and safe-line bounds inherited
from the canonical analysis in \cite{BakerBase,BakerGT41}. Specifically,
the base manuscript gives the projected count in ``Uniform counts,
including time zero'' and the safe-line estimate in ``A convex balance
with conjugate dissipation.'' With one
constant $\Cu\ge1$, enlarged to cover these inputs,
\begin{equation}
 \begin{gathered}
 |\Im\rho|\le1,\qquad
 \#\{\rho:v\le\Re\rho\le v+1\}\le\Cu\log(2+|v|),\\
 \left|\frac{H_z}{H}(t,v\pm4\ii)\right|
       \le\Cu\log(2+|v|),\qquad v\in\R,\quad t\in I.
 \end{gathered}
 \label{eq:count-safe}
\end{equation}
These are stated inputs, not conclusions of the moment identities.
In particular, those identities do not create or prolong the tube.

Write
\begin{equation}
 Q(t,w)=w^2-D(t),\quad Q_*(t,z)=(z-c(t))^2-D(t),\quad
 Q_j=Q(t,w_j),\quad \mu_j=\frac{\ell^2}{Q_j},\quad q_0=\frac{16}{15}.
 \label{eq:quadratic}
\end{equation}
The collar implies
\begin{equation}
 |Q_j|\ge\frac{15}{16}|w_j|^2,\qquad
 |\mu_j|\le\gamma:=\frac{q_0}{a^2}<1.
 \label{eq:collar-kernel}
\end{equation}
No right-half-plane condition is imposed on $Q_j$.

\subsection{Symmetric coordinates and the pair equation}

At a reference time $t_*$ choose the positive circle
$\Gamma=\{|z-c(t_*)|=r_0/2\}$, fixed while differentiating in $t$.
It has positive distance from all zeros. Its count remains two for
nearby complex times. If $B=H_z/H$, the argument principle gives
\begin{equation}
 J_k=\frac1{2\pi\ii}\oint_\Gamma z^kB(t,z)\dd z,
 \qquad c=J_1/2,\qquad D=J_2/2-c^2.
 \label{eq:coordinates}
\end{equation}
These formulas define holomorphic coordinates even at a double zero.
The factorization $H=Q_*A$ has $A$ jointly holomorphic and nonzero
inside $\Gamma$. This follows either from division by the counted
quadratic factor or from the Cauchy extension of the analytic boundary
quotient. Put $g=A_z/A$ there. Inserting this factorization into
\eqref{eq:heat}, with $w=z-c$, gives
\begin{equation}
 \frac{A_t}{A}=-(g_z+g^2)
       +\frac{2wc'+D'-2-4wg(t,c+w)}{w^2-D}.
 \label{eq:deflated}
\end{equation}
The numerator must vanish to the target multiplicities. Interpolating
$h_0(w)=wg(t,c+w)$ at $\pm s$, with Hermite interpolation at $s=0$,
therefore yields
\begin{equation}
 \begin{split}
 c'&=g(t,c-s)+g(t,c+s),\\
 D'&=2+4Dg[t;c-s,c+s]=2-4DM,\qquad
 M:=-g[t;c-s,c+s],\\
 \frac{A_t}{A}&=-(g_z+g^2)-4h_0[-s,s,w].
 \end{split}
 \label{eq:pairlaw}
\end{equation}
Divided differences are understood by analytic continuation at repeated
nodes. At a merger $M=-g_z(t,c)$ and $D'=2$. Thus a collision creates
no singularity in the signed discriminant equation.

\subsection{Rational traces and the complete divisor}

\begin{proposition}[Exterior traces]\label{prop:trace}
For $n\ge1$ and $k\ge2$, respectively, the moments
\begin{equation}
 S_n=\sum_j Q_j^{-n},\qquad T_k=\sum_j w_jQ_j^{-k}
 \label{eq:moments}
\end{equation}
are absolutely convergent and real on $I$. They have the local
holomorphic contour expressions
\begin{equation}
 S_n=-\frac1{2\pi\ii}\oint_\Gamma Q_*^{-n}B\dd z,
 \qquad
 T_k=-\frac1{2\pi\ii}\oint_\Gamma(z-c)Q_*^{-k}B\dd z.
 \label{eq:traces}
\end{equation}
In particular $M=S_1$. Their time derivatives exist through arbitrary
exterior mergers and a double target.
\end{proposition}

\begin{proof}
The unit count implies $\sum_\rho|\rho|^{-2}<\infty$.
The order bound and evenness give the canonical product
\begin{equation}
 H(t,z)=H(t,0)\prod_{\{\rho,-\rho\}}
              (1-z^2/\rho^2),
 \label{eq:product}
\end{equation}
where opposite groups, their conjugates and multiplicities are retained.
Indeed the genus-one factors of an opposite pair multiply to the displayed
factor, and the only possible linear exponential prefactor is constant
by evenness. The grouped logarithmic derivative converges locally away
from the divisor. The finite collar and the unit count give absolute
convergence of the moments; the slowest tails are $|\rho|^{-2}$
and $|\rho|^{-3}$.

For either kernel $R(w)=Q(w)^{-n}$ or $wQ(w)^{-k}$, the residue theorem
for the rational function $R(z-c)/(z-\rho_j)$ gives
\[
 \frac1{2\pi\ii}\oint_\Gamma
                 \frac{R(z-c)}{z-\rho_j}\dd z=-R(w_j).
\]
There is no residue at infinity. The genus-one additive constants
integrate to zero because the sum of the residues of $R$ is zero.
The same is true of the target contribution $R\,\partial_z Q_*/Q_*$:
it has poles only at the targets and decays faster than $1/z$.
Its combined target residue vanishes, even though the two residues
need not vanish individually. Uniform convergence on $\Gamma$ now
proves \eqref{eq:traces}. Taking $n=1$ also gives
$S_1=-g[t;c-s,c+s]$. The fixed-circle formulas and
\eqref{eq:coordinates} prove local analyticity without labelling
individual exterior roots. Reality follows from conjugation.
\end{proof}

\begin{lemma}[Passage from polynomial identities at a time slice]
\label{lem:slice-passage}
Suppose an identity for a fixed finite collection of the moments
\eqref{eq:moments}, $D$, and their first time derivatives is polynomial
in absolutely convergent moments, after all nonabsolute terms have
cancelled. If it holds for every finite polynomial heat field with an
isolated target pair, it holds for $H$ under \eqref{eq:domain}--
\eqref{eq:count-safe}, including multiple zeros.
\end{lemma}

\begin{proof}
First, for a finite polynomial at $t_*$, perturb its roots to distinct
ones preserving the count on $\Gamma$, and initialize its own polynomial
heat evolution there. The derivative
\begin{equation}
 B_t=-\frac{H_{zzz}}H+\frac{H_zH_{zz}}{H^2}
 \label{eq:log-time}
\end{equation}
and the fixed-circle coordinate and trace formulas depend continuously
on the polynomial and its first three spatial derivatives. The contour
denominators stay bounded away from zero. Thus an identity proved for
simple roots passes to all finite multiplicities.

Now truncate \eqref{eq:product} at increasing radii, retaining all ties,
opposite groups and conjugates. Include the targets for all sufficiently
large cutoffs. The product tail on a fixed compact set is
$O(\log N/N)$ by the unit count. The polynomials and every fixed
spatial derivative converge locally uniformly to $H(t_*,\cdot)$.
Give each truncated polynomial its own heat evolution initialized at
$t_*$. Formula \eqref{eq:log-time} and the trace formulas give convergence
of the values and first time derivatives of the stated observables at
that time. The remaining exterior sums converge absolutely. Passing
the identity to the limit proves it at $t_*$; repeat at every time
of $I$, with inward derivatives at its endpoints. This does not require
a common time interval for the truncated heat evolutions, and gives
no estimate uniform as the moment order tends to infinity.
\end{proof}

\section{A collision-regular integer hierarchy}
\label{sec:integer}

On the tube \eqref{eq:domain}, set \(w_j=\rho_j-c\), \(Q_j=w_j^2-D\),
and \(m_j=Q_j^{-1}\) for each exterior occurrence. The rational
moments are
\[
 S_n=\sum_j m_j^n\quad(n\geq 1),\qquad
 T_n=\sum_j w_jm_j^n\quad(n\geq 2),\qquad M=S_1 .
\]
They are real by conjugation. The collar gives
\(\lvert Q_j\rvert\geq(15/16)\lvert w_j\rvert^2\). This bound and the
unit count in \eqref{eq:count-safe} give absolute convergence in the
indicated index ranges. In particular, the first odd moment is outside
these ranges. The hierarchy extends the interaction law
of \cite{BakerGT41} and the next even moment calculation of
\cite{BakerGR02}, using the rational divided differences of
\cite{BakerCM09}. Primes denote physical time derivatives.

\begin{theorem}[Integer moment evolution]\label{thm:hierarchy}
On the complete canonical tube, the regular moments satisfy
\begin{align}
 S_n'={}&2n\sum_{k=1}^{n}S_kS_{n+1-k}
       +2n\sum_{k=2}^{n}\bigl(T_kT_{n+2-k}+DS_kS_{n+2-k}\bigr)
       \notag\\[-1pt]
       &\quad-4n(n+2)S_{n+1}-4n(n+3)DS_{n+2},
       \qquad n\geq1, \label{eq:S-hierarchy}\\
 T_n'={}&2(2n-1)\sum_{k=2}^{n}T_kS_{n+1-k}
       +4nD\sum_{k=2}^{n}T_kS_{n+2-k}
       \notag\\[-1pt]
       &\quad-4(n^2+n-1)T_{n+1}-4n(n+3)DT_{n+2},
       \qquad n\geq2. \label{eq:T-hierarchy}
\end{align}
Empty sums are zero. These equations and \(D'=2-4DM\) hold through
target and exterior multiplicities, with inward derivatives at the
interval endpoints. Each equation involves finitely many terms of the
infinite relative moment family. The absolute midpoint \(c\) uses
its counted contour law.
\end{theorem}

\begin{proof}
First consider a polynomial satisfying the physical heat equation at
a slice with simple roots. Factoring out a simple root gives
\(\rho'=2\sum_{\eta\ne\rho}(\rho-\eta)^{-1}\). At this finite stage
write \(P=\sum_j w_jm_j=T_1\). Adding the target velocities gives
\(c'=-2P\). Multiplying their velocity difference by \(s\), where
\(s^2=D\), gives \(D'\). The mutual target force contributes \(2\)
to this derivative, and each exterior occurrence contributes
\(-4Dm_j\). Thus \(D'=2-4DM\), as in \eqref{eq:pairlaw}.
The contribution of the targets to the velocity at an exterior root
is \(4w_jm_j\). Hence
\begin{equation}
 w_j'=4w_jm_j+2P+2\sum_{k\ne j}\frac{1}{w_j-w_k},
 \label{int:relative}
\end{equation}
where the last sum ranges over exterior occurrences.

For the symmetrization, let \(A=(u^2-D)^{-1}\), \(B=(v^2-D)^{-1}\),
and \(L_h=\sum_{i=0}^{h-1}A^{h-1-i}B^i\) for \(h\geq1\), with
\(L_0=0\). The two divided differences are
\begin{align}
 \mathcal E_h(u,v)&=\frac{A^h-B^h}{u-v}=-(u+v)AB L_h, \notag\\
 \mathcal O_h(u,v)&=\frac{uA^h-vB^h}{u-v}
              =-AB L_{h-1}-(D+uv)AB L_h. \label{int:divdiff}
\end{align}
Indeed \(A-B=-(u-v)(u+v)AB\), and
\(uA-vB=-(u-v)(D+uv)AB\). The identity
\(x^h=xL_h-AB L_{h-1}\) for \(x=A,B\) follows by cancellation of
successive terms. It gives the second line of \eqref{int:divdiff}
from the first-power difference. Both rational expressions on the
right are regular at \(u=v\) away from zeros of the denominators.

The product \(AB L_h\) is \(\sum_{k=1}^h A^k B^{h+1-k}\).
The symmetric kernels in \eqref{int:divdiff} therefore give
\begin{align}
 -\sum_{i<j}\mathcal O_{n+1}(w_i,w_j)
  ={}&\frac12\sum_{k=1}^{n}
             \bigl(S_kS_{n+1-k}-S_{n+1}\bigr)\notag\\
   &+\frac D2\sum_{k=1}^{n+1}
             \bigl(S_kS_{n+2-k}-S_{n+2}\bigr)\notag\\
   &+\frac12\sum_{k=1}^{n+1}
             \bigl(T_kT_{n+2-k}-S_{n+1}-DS_{n+2}\bigr),
     \label{int:odd-pairs}\\
 -\sum_{i<j}\mathcal E_h(w_i,w_j)
  ={}&\sum_{k=1}^{h}\bigl(T_kS_{h+1-k}-T_{h+1}\bigr).
     \label{int:even-pairs}
\end{align}
The factor \(1/2\) in \eqref{int:odd-pairs} converts the ordered
off-diagonal sum to the unordered sum. The three diagonal subtractions
come from \(A^{n+1}\), \(DA^{n+2}\), and
\(w^2 A^{n+2}=A^{n+1}+DA^{n+2}\), respectively.
The two contributions from \(u+v\) in the even difference give the
same ordered sum after \(k\) is reversed. Each diagonal contributes
\(T_{h+1}\), and the two terms cancel the ordering factor \(1/2\).

For the even weight \(f(w,D)=(w^2-D)^{-n}\),
\(f_D=n(w^2-D)^{-n-1}\) and \(f_w=-2nw(w^2-D)^{-n-1}\).
Differentiating along \eqref{int:relative} gives
\begin{align}
 S_n'={}&nD'S_{n+1}-8n(S_{n+1}+DS_{n+2})
          -4nPT_{n+1}-4n\sum_{i<j}\mathcal O_{n+1}(w_i,w_j).
 \label{int:even-before}
\end{align}
The \(k=1,n+1\) terms of the \(T\)-convolution in
\eqref{int:odd-pairs} contribute \(4nPT_{n+1}\) in
\eqref{int:even-before}. The same two endpoints of the \(D\)-weighted
convolution contribute \(4nDMS_{n+1}\), which cancels the interaction
term of \(nD'S_{n+1}\). The remaining coefficient of \(S_{n+1}\)
is \(2n-8n-2n^2-2n(n+1)=-4n(n+2)\). The coefficient of
\(DS_{n+2}\) is \(-8n-4n(n+1)=-4n(n+3)\). This yields
\eqref{eq:S-hierarchy}.

For the odd weight \(f(w,D)=w(w^2-D)^{-n}\) one has
\[
 f_D=nw(w^2-D)^{-n-1},\qquad
 f_w=(1-2n)(w^2-D)^{-n}-2nD(w^2-D)^{-n-1}.
\]
Consequently
\begin{align}
 T_n'={}&nD'T_{n+1}+4(1-2n)T_{n+1}-8nDT_{n+2}\notag\\
       &+2P\bigl((1-2n)S_n-2nDS_{n+1}\bigr)\notag\\
       &+2(1-2n)\sum_{i<j}\mathcal E_n(w_i,w_j)
          -4nD\sum_{i<j}\mathcal E_{n+1}(w_i,w_j).
       \label{int:odd-before}
\end{align}
By \eqref{int:even-pairs}, the \(k=1\) terms cancel both terms
proportional to \(P\). The \(k=n+1\) term of the \(D\)-weighted
sum yields \(4nDMT_{n+1}\), canceling the interaction part of
\(nD'T_{n+1}\). The other diagonal and direct terms give
\(2n+4(1-2n)+2n(1-2n)=-4(n^2+n-1)\) for \(T_{n+1}\) and
\(-8n-4n(n+1)=-4n(n+3)\) for \(DT_{n+2}\).
The two surviving convolution ranges are \(2\leq k\leq n\).
This proves \eqref{eq:T-hierarchy} at the finite simple slice.
The arguments have eliminated every use of \(P\).

For the complete passage, fix a reference slice \(\tau\) and a
stationary positive physical circle \(\Gamma\) enclosing the two
targets and excluding every exterior occurrence. Use
\(u=z-c(t)\), \(Q_*(t,z)=u^2-D(t)\).
For a weight \(f=Q_*^{-n}\) or \(f=uQ_*^{-n}\),
Proposition~\ref{prop:trace} writes the exterior sum as
\(-\langle f(H_z/H-Q_{*,z}/Q_*)\rangle_\Gamma\), where
\(\langle g\rangle_\Gamma=(2\pi\ii)^{-1}\oint_\Gamma g(z)\,\dd z\).
Both poles remain inside the fixed circle, including when they
coalesce. The combined target term integrates to zero since
\(fQ_{*,z}/Q_*=O(z^{-3})\) at infinity and all its poles are inside.
It vanishes as an identity in time. Differentiation of the sum
retains the full motion of the kernel, using \(u_t=-c'\) and
\(Q_{*,t}=-2uc'-D'\), together with
\[
 \partial_t(H_z/H)=-H_{zzz}/H+H_zH_{zz}/H^2 .
\]
The contour expression is locally holomorphic in time.

For a finite multiple-root slice, perturb the root factors at that
slice and initialize a separate polynomial heat solution for each
perturbation. The strict clearance keeps the target count. The
spatial derivatives through order three converge on \(\Gamma\).
The displayed heat derivative and the counted target moments then
give convergence of the first contour jets, including \(c',D'\).
The right sides of \eqref{eq:S-hierarchy} and \eqref{eq:T-hierarchy}
are continuous finite sums of occurrence values, so the identities
persist at polynomial multiples.

For the canonical field, apply Lemma~\ref{lem:slice-passage} at
\(\tau\) with modulus cutoff \(N\), keeping the entire opposite and
conjugate groups and every tied occurrence. The truncated slice
contains the exact targets for large \(N\). Its heat derivative at
\(\tau\) is determined by its first three spatial derivatives on
\(\Gamma\). The omitted inverse-square product tail
\(O(\log N/N)\) gives uniform compact convergence of those
derivatives by Cauchy's formula. Together with the strict clearance
on the fixed circles, this gives convergence of the time jets of
the rational moments and of \(c,D\) at the reference slice.

For \(N>4T+8\), an omitted root has \(\lvert w_j\rvert\geq
\lvert\rho_j\rvert/2\). A dyadic shell of radii comparable to
\(2^kN\) contains \(O(\Cu\,2^kN\log(2^kN))\) roots by the strip
and unit count. Summation against the decay powers \(2n\) and
\(2n-1\) gives the respective moment errors
\(O(N^{1-2n}\log N)\) and \(O(N^{2-2n}\log N)\).
Thus all products in each fixed equation converge. A finite
cover by cleared circles gives the jet passage uniformly on the
compact established tube at fixed \(T\) and fixed orders. The argument
proves the canonical equations using these reference-slice jets alone.
\end{proof}

In particular the first forcing moments and the first odd derivative are
\begin{equation}
 \begin{aligned}
 S_1'&=2M^2-12S_2-16DS_3,\\
 S_2'&=8MS_2-32S_3+4DS_2^2-40DS_4+4T_2^2,\\
 T_2'&=(6M+8DS_2)T_2-20T_3-40DT_4 .
 \end{aligned}
 \label{eq:low-orders}
\end{equation}
The spatial term \(4T_2^2\) in \(S_2'\) has its inherited
positive sign \cite{BakerGR02,BakerGR05}. The equations describe
componentwise evolution for the given field. Arbitrary-sequence
well-posedness, finite truncation, and a signed forcing estimate
are separate questions.

\section{Separate powers of the target poles}
\label{sec:split}

Choose an analytic square root \(s\) of \(D\) on a local chart where
\(D\ne0\). Write \(z_\pm=c\pm s\), \(d_0=2s\), \(x_j=w_j-s\), and
\(y_j=w_j+s=x_j+d_0\). The split-pole moments are
\begin{equation}
 P_{p,q}=\sum_j x_j^{-p}y_j^{-q},
 \qquad p,q\in\mathbb Z_{\geq0},\quad p+q\geq2.
 \label{split:def}
\end{equation}
Both \(\lvert x_j\rvert\) and \(\lvert y_j\rvert\) are at least
\(3\lvert w_j\rvert/4\) by \eqref{eq:domain}. Thus all of
\eqref{split:def}, including its pure powers, converge absolutely.

The two moment families have an exact algebraic relation.
For \(h=\max(p,q)\) one has
\begin{equation}
 \frac1{(w-s)^p(w+s)^q}
    =\frac{(w-s)^{h-p}(w+s)^{h-q}}{(w^2-D)^h}.
 \label{split:reduction}
\end{equation}
The numerator has degree at most \(2h-2\) in \(w\).
Write an even power as \(w^{2k}=(Q+D)^k\) and an odd power as
\(w^{2k+1}=w(Q+D)^k\), with \(Q=w^2-D\).
Then \(k\leq h-1\) in the even terms and \(k\leq h-2\) in the
odd terms. Expansion of the binomial gives only \(S_n\) with
\(n\geq1\) and \(T_n\) with \(n\geq2\), with polynomial
coefficients in \(s\) after \(D=s^2\). Conversely,
\begin{equation}
 S_n=P_{n,n}\quad(n\geq1),\qquad
 T_n=\tfrac12(P_{n-1,n}+P_{n,n-1})\quad(n\geq2).
 \label{split:inverse}
\end{equation}
In particular, since $w=(w-s)+s=(w+s)-s$,
\begin{equation}
 T_2=P_{1,2}+sP_{2,2}=P_{2,1}-sP_{2,2}.
 \label{split:T2}
\end{equation}
This uses \(xy=Q\) and \(x+y=2w\). Both changes of expressions
use the absolutely convergent index ranges.

\begin{proposition}[Split-pole recurrence]\label{split:prop}
For \(m,n\geq1\) write
\begin{equation}
 C_{m,n}=\sum_{i=2}^{m}P_{i,0}P_{m+1-i,n}
            +\sum_{j=1}^{n-1}P_{m,j}P_{0,n+1-j}.
 \label{split:convolution}
\end{equation}
On the axes use \(C_{m,0}=\sum_{i=2}^{m-1}P_{i,0}P_{m+1-i,0}\)
and \(C_{0,n}=\sum_{j=2}^{n-1}P_{0,j}P_{0,n+1-j}\) for the
respective index at least two. Let \(A_d=2/d_0-d_0M\).
For \(p,q\geq1\), the complete hierarchy \cite{BakerCM08} is
\begin{align}
 P_{p,q}'={}&p\bigl(C_{p+1,q}-(p+3)P_{p+2,q}-(q+2)P_{p+1,q+1}
                                  +A_dP_{p+1,q}\bigr)\notag\\
   &+q\bigl(C_{p,q+1}-(p+2)P_{p+1,q+1}-(q+3)P_{p,q+2}
                                  -A_dP_{p,q+1}\bigr).
 \label{split:mixed}
\end{align}
The pure powers satisfy
\begin{align}
 P_{p,0}'&=p\bigl(C_{p+1,0}-(p+3)P_{p+2,0}-2P_{p+1,1}
                                   +(2/d_0)P_{p+1,0}\bigr),
       \qquad p\geq2,\notag\\
 P_{0,q}'&=q\bigl(C_{0,q+1}-2P_{1,q+1}-(q+3)P_{0,q+2}
                                   -(2/d_0)P_{0,q+1}\bigr),
       \qquad q\geq2 .
 \label{split:axes}
\end{align}
The separation obeys \(d_0'=4/d_0-2d_0M\). The equations
are valid on the separated charts, including exterior multiplicities.
\end{proposition}

\begin{proof}
At a finite simple-root slice define \(P_{1,0}\) and \(P_{0,1}\)
as finite sums. The target velocities are
\[
 z_+'=2/d_0-2P_{1,0},\qquad z_-'=-2/d_0-2P_{0,1}.
\]
The identity \(x_j^{-1}-y_j^{-1}=d_0/Q_j\) gives
\(P_{1,0}-P_{0,1}=d_0M\) and the stated equation for \(d_0'\).
For the complete divisor this difference is the absolutely
convergent series \(\sum_j d_0/Q_j\), or the limit of differences
with a common cutoff.

At an exterior root the contribution of the two targets to
\(\rho_j'\) is \(2(x_j^{-1}+y_j^{-1})\).
Subtracting \(z_+'\) for \(x_j'\) and \(z_-'\) for \(y_j'\), and
differentiating the weight \(f_{p,q}(w)=(w-s)^{-p}(w+s)^{-q}\),
gives the following two contributions to \(P_{p,q}'\):
\begin{align*}
 &p\bigl((2/d_0-2P_{1,0})P_{p+1,q}
                    -2P_{p+2,q}-2P_{p+1,q+1}+U_{p+1,q}\bigr),\\
 &q\bigl((-2/d_0-2P_{0,1})P_{p,q+1}
                    -2P_{p+1,q+1}-2P_{p,q+2}+U_{p,q+1}\bigr).
\end{align*}
Omit a block with zero prefactor. Here the exterior pair term is
\[
 U_{m,n}=-2\sum_{j<k}
        \frac{f_{m,n}(w_j)-f_{m,n}(w_k)}{w_j-w_k}.
\]
Splitting the difference between the two factors gives
\begin{equation}
 -\frac{f_{m,n}(u)-f_{m,n}(v)}{u-v}
   =\sum_{i=1}^{m} f_{i,0}(u)f_{m+1-i,n}(v)
     +\sum_{j=1}^{n} f_{m,j}(u)f_{0,n+1-j}(v).
 \label{split:divdiff}
\end{equation}
For the first sum subtract \((u-s)^{-m}\) and \((v-s)^{-m}\)
using the difference of reciprocal powers, with
\((v+s)^{-n}\) held fixed. For the second use the corresponding
reciprocal powers at \(-s\) with \((u-s)^{-m}\) fixed.
The full kernel is symmetric in \(u,v\). Therefore \(U_{m,n}\)
is the sum of the right side over ordered off-diagonal pairs.
Subtracting the diagonal gives \(-mP_{m+1,n}\) and
\(-nP_{m,n+1}\) in the two sums.

When \(m,n\geq1\) the term \(i=1\) and the term \(j=n\) yield
\((P_{1,0}+P_{0,1})P_{m,n}\). Their combination with the two
moving-target blocks is \(-d_0MP_{p+1,q}\) and
\(+d_0MP_{p,q+1}\), respectively. The other products form
\eqref{split:convolution}. Adding the two target terms in each
block to the diagonal subtractions gives the four coefficients
in \eqref{split:mixed}. On the \(m\)-axis the two endpoints
\(i=1,m\) give \(2P_{1,0}P_{m,0}\), and on the \(n\)-axis the
endpoints \(j=1,n\) give \(2P_{0,1}P_{0,n}\). They cancel the
first-power sums in the corresponding blocks. The diagonals
here are \(-mP_{m+1,0}\) and \(-nP_{0,n+1}\). This proves
\eqref{split:axes}. The cancellation precedes all limiting steps.

The fixed physical circle \(\Gamma\) of Proposition~\ref{prop:trace}
also gives the contour formula for the split weight
\(k(z)=(z-z_+)^{-p}(z-z_-)^{-q}\).
The integral of \(k\) vanishes by its decay \(O(z^{-p-q})\).
The residue of \(k(z)/(z-\rho_j)\) at an exterior pole is
\(k(\rho_j)\), so the integral on \(\Gamma\) gives its negative.
In the convergent product logarithmic derivative each constant
regularizer multiplies the zero integral of \(k\). Hence
\begin{equation}
 P_{p,q}=-\frac1{2\pi\ii}\oint_\Gamma
                  k(z)\left(\frac{H_z}{H}-\frac{Q_{*,z}}{Q_*}\right)
                   \,\dd z.
 \label{split:trace}
\end{equation}
The \(Q_{*,z}/Q_*\) term accounts for both target poles.
Its weighted integral vanishes since all poles of \(kQ_{*,z}/Q_*\)
are enclosed and the decay is \(O(z^{-p-q-1})\).
For \(p=2,q=0\) the individual target residues are
\(-d_0^{-2}\) at \(z_+\) and \(+d_0^{-2}\) at \(z_-\).
Separating the two pole contours would retain these two terms.
For \(p=q\) the target term is the spatial derivative of
\(-Q_*^{-p}/p\).

Differentiate \eqref{split:trace} on the fixed circle, including
\(k_t=k(p z_+'/(z-z_+)+q z_-'/ (z-z_-))\).
The first contour jets converge by the argument of
Lemma~\ref{lem:slice-passage}. At a finite multiple-root slice
the root factors can be perturbed without changing the separated
counts, then evolved from that slice. The derivatives of \(z_\pm\)
come from the jets of \(c,D\) and \(s'=D'/(2s)\), with \(s\) bounded
away from zero on each compact subchart. This also applies to
the opposite-paired canonical product truncations there.
The omitted occurrences in \eqref{split:def} contribute
\(O(N^{1-p-q}\log N)\) at fixed orders, by the strip and unit count
as in the proof of Theorem~\ref{thm:hierarchy}. Every term of the
already canceled equations converges. This proves the complete
laws. At an interval endpoint contained in the chart the derivatives
are inward.
\end{proof}

For example, on the diagonal the identities
\[
 P_{n+1,n}-P_{n,n+1}=d_0 S_{n+1},\qquad
 P_{n+2,n}+P_{n,n+2}=2S_{n+1}+d_0^2S_{n+2}
\]
give, directly from \eqref{split:mixed},
\begin{align*}
 S_n'={}&n(C_{n+1,n}+C_{n,n+1})-4n(n+2)S_{n+1}\\
        &\quad-n(n+3)d_0^2S_{n+2}-nd_0^2MS_{n+1}.
\end{align*}
At \(n=1\) the convolution is \(M(2M+d_0^2 S_2)\).
At \(n=2\) its five products are
\[
 S_2(P_{2,0}+P_{0,2})+P_{3,0}P_{1,2}+P_{2,1}P_{0,3}
                       +P_{3,1}P_{0,2}+P_{2,0}P_{1,3}.
\]
Pair these products using \(x+y=2w\), \(y-x=d_0\), and \(xy=Q\).
For the two crossed pairs use
\begin{align*}
 P_{3,0}+P_{0,3}&=2T_2+2d_0^2T_3,&
 P_{3,0}-P_{0,3}&=d_0(3S_2+d_0^2S_3),\\
 P_{3,1}+P_{1,3}&=2S_2+d_0^2S_3,&
 P_{3,1}-P_{1,3}&=2d_0 T_3,\\
 P_{2,0}+P_{0,2}&=2M+d_0^2 S_2,&
 P_{2,0}-P_{0,2}&=2d_0 T_2,\\
 P_{2,1}+P_{1,2}&=2T_2,&
 P_{2,1}-P_{1,2}&=d_0 S_2 .
\end{align*}
The terms \(T_2T_3\) and \(d_0^4S_2S_3\) cancel, leaving
\(4MS_2+2T_2^2+d_0^2MS_3+(d_0^2/2)S_2^2\).
The two values of \(n\) reproduce \eqref{eq:low-orders} for \(S_1,S_2\).

\subsection{Coalescing labels and an exact quartic}

Changing \(s\) to \(-s\) swaps the two indices, and the right
sides of the recurrence transform by the same exchange. At real
times, conjugation fixes the moments for \(D>0\) and swaps the
two indices for \(D<0\). At a merger the two raw labels have the
same limit. In fact, \eqref{split:reduction} and the analytic
rational moments give a decomposition
\[
 P_{p,q}(t)=A(t)+s(t)B(t)
\]
with \(A,B\) locally analytic in \(t\). They are obtained by
separating the even and odd polynomial powers of \(s\) and
substituting \(s^2=D(t)\). At a target merger \(t_*\), expansion of
the summand in \(s\) at the nonzero exterior displacements gives
\[
 A(t_*)=\sum_j w_j(t_*)^{-p-q},\qquad
 B(t_*)=(p-q)\sum_j w_j(t_*)^{-p-q-1}.
\]
Both sums converge absolutely. The symmetric combination and
the divided antisymmetric combination recover \(A\) and \(B\).
An individual coordinate has derivative asymptotic to \(B(t_*)/s\)
if \(B(t_*)\ne0\), since \eqref{eq:pairlaw} gives \(D'(t_*)=2\).
The condition on \(B(t_*)\) remains open at a canonical merger.

The singularity occurs in a finite even quartic heat field.
For \(b>0\) consider
\begin{equation}
 F(t,z)=(z^2-b^2)^2-t(12z^2-4b^2)+12t^2 .
 \label{split:quartic}
\end{equation}
It satisfies \(F_t=-F_{zz}\) by differentiation. A fixed small
circle around \(+b\) counts the two targets near \(t=0\). The
exterior pair is their reflection, with center \(-c\) and the same
\(D\). Setting \(u=z-b\) at \(t=0\), one has
\[
 \frac{F_t}{F}
  =\frac{-8b^2-24bu-12u^2}{u^2(2b+u)^2}
  =-2u^{-2}-\frac4b u^{-1}+O(1).
\]
Comparison with the derivative of the target quadratic gives
\(D'(0)=2\), \(c'(0)=2/b\). The mean exterior displacement
\(\bar w=-2c\) thus has value \(u_0=-2b\) and derivative \(8/u_0\).
For a weight \(f(w,D)\) analytic near \((u_0,0)\), the trace over
the exterior pair is
\[
 f(\bar w+s,D)+f(\bar w-s,D)
    =2f(\bar w,D)+Df_{ww}(\bar w,D)+O(D^2).
\]
The derivative at the double is
\((16/u_0) f_w+4f_D+2f_{ww}\), evaluated at \((u_0,0)\).
For \(f=Q^{-n}\), the three contributions have coefficients
\(-32n,4n,4n(2n+1)\) of \(u_0^{-2n-2}\). For \(f=wQ^{-n}\)
they have coefficients \(16(1-2n),4n,-4n(1-2n)\) of
\(u_0^{-2n-1}\). In particular
\[
 S_n'(0)=8n(n-3)u_0^{-2n-2},\qquad
 T_n'(0)=8(n^2-4n+2)u_0^{-2n-1}.
\]
Independently, inserting \(S_n(0)=2u_0^{-2n}\),
\(T_n(0)=2u_0^{1-2n}\), and \(D(0)=0\) in the regular equations
gives the coefficients
\[
 8n^2+8n(n-1)-8n(n+2),
 \qquad
 8(2n-1)(n-1)-8(n^2+n-1),
\]
respectively. These reduce to the same derivatives.

Here \(P_{2,0}=S_1+2DS_2+2sT_2\) by \eqref{split:reduction},
and \(T_2(0)=2u_0^{-3}\ne0\). Thus the square-root term is
nonzero on either label for \(t>0\) near zero, and the raw
coordinate has no finite derivative through this double.
This example is in the finite polynomial heat class. It
exhibits the limitation of the split labels while the regular
integer moments remain analytic.

% 18 November 2026. Prepared for Edward Baker with OpenAI Codex LLM assistance.
% Actual model and effort unavailable; inherited helper settings. E26 source integration.
\section{A coupled hierarchy indexed by the adjacent pair}
\label{sec:lattice}

The split-pole moments can be attached to every adjacent pair of the
same divisor. This adds an index locating the pair, while imposing
compatibility between rows that share a zero. The resulting description
retains correlations which independent moment rows would discard.
We develop the formulation proposed in \cite{BakerCM29}, using the
single-pair identities of \cite{BakerCM08,BakerCM09}. The conclusions
are exact representation and error identities. An estimate preventing
concentration into an individual gap remains a separate requirement.

\subsection{Edge coordinates and the complete gap interaction}

First suppose the complete divisor is real and simple on the time
interval under consideration, and order its zeros as $x_n(t)$,
$n\in\mathbb Z$. The integer labels are continued locally in time.
All infinite-lattice assertions below are componentwise; uniform
estimates concern a fixed finite set of edges on a compact interval
before the relevant loss of simplicity or separation. Define
\begin{equation}
 g_n=x_{n+1}-x_n>0,\qquad
 c_n=\frac{x_n+x_{n+1}}2,\qquad D_n=\frac{g_n^2}{4},
 \label{lat:coordinates}
\end{equation}
and, in the orientation of \eqref{split:def},
\begin{equation}
 M_n(p,q;t)=\sum_{k\ne n,n+1}
       (x_k-x_{n+1})^{-p}(x_k-x_n)^{-q},
 \quad p,q\ge0,\quad p+q\ge2.
 \label{lat:moments}
\end{equation}
The index $n$ belongs to the edge $(x_n,x_{n+1})$. It does not
specify a new exterior source. In particular,
$M_n(1,1)$ is the interaction $M$ of that edge,
$M_n(r,r)=S_r^{(n)}$, and
$\bigl(M_n(r-1,r)+M_n(r,r-1)\bigr)/2=T_r^{(n)}$ for $r\ge2$.

For these exact identities, a strictly separated count-two contour
suffices at each edge; the stronger collar ratio in
\eqref{eq:domain} is not imposed on a general adjacent-edge lattice.
On a compact subinterval, a finite cover by such contours gives
positive denominator bounds for each retained edge. The finite
algebra in Proposition~\ref{split:prop} and its complete-product
passage therefore apply with those bounds. Its recurrence
\eqref{split:mixed}--\eqref{split:axes} holds at every $n$, replacing
$P_{p,q}$ by $M_n(p,q)$ and $d_0$ by $g_n$. In particular, in
forward physical time,
\begin{equation}
 g_n'=\frac4{g_n}-2g_nM_n(1,1),\qquad
 D_n'=2-4D_nM_n(1,1).
 \label{lat:gap-law}
\end{equation}
There is no additional nearest-neighbor term to append to the
single-pair recurrence. Coupling in $n$ is already encoded by the
common zeros and the identities derived below.

The complete sums in \eqref{lat:moments} are absolute at their stated
orders. For targets in a fixed compact spatial set and $N$ sufficiently
large, the terms with $|\rho|>N$ are bounded by
$C_{p,q}|\rho|^{-p-q}$. The unit-count bound
\eqref{eq:count-safe} consequently gives the tail
\begin{equation}
 \sum_{|\rho|>N}
 |\rho-x_{n+1}|^{-p}|\rho-x_n|^{-q}
 =O_{p,q,I,\mathcal B}(N^{1-p-q}\log N),
 \label{lat:tail}
\end{equation}
uniformly for a fixed finite edge block $\mathcal B$ with the stated
clearances. Opposite and conjugate occurrences and all multiplicities
are included. No uniformity in a growing block, height or moment order
is inferred from this estimate.

\begin{proposition}[All-gap expression]\label{lat:all-gap-prop}
When the complete divisor is real and ordered, put
\begin{equation}
 L_{n,m}=\sum_{j=1}^{m}g_{n-j},\qquad
 R_{n,m}=\sum_{j=1}^{m}g_{n+j},\qquad m\ge1.
 \label{lat:distances}
\end{equation}
Then the complete interaction is
\begin{equation}
 M_n(1,1)=\sum_{m\ge1}
 \left\{\frac1{L_{n,m}(L_{n,m}+g_n)}
       +\frac1{R_{n,m}(R_{n,m}+g_n)}\right\}.
 \label{lat:all-gap}
\end{equation}
Thus \eqref{lat:gap-law} and \eqref{lat:all-gap} describe a nonlocal
evolution of the complete gap sequence. They do not specify the
absolute translation coordinate.
\end{proposition}

\begin{proof}
The left exterior zero $x_{n-m}$ has distances $-L_{n,m}$ and
$-(L_{n,m}+g_n)$ from the left and right targets. The right exterior
zero $x_{n+1+m}$ has distances $R_{n,m}+g_n$ and $R_{n,m}$.
These exhaust the exterior divisor. Substitution in
\eqref{lat:moments} proves the formula; absolute convergence from
\eqref{lat:tail} permits the partition. The gap law has already
eliminated the common center velocity, so it determines relative
motion only.
\end{proof}

Every gap enters many rows through cumulative distances in
\eqref{lat:distances}. This is a spatially nonlocal lattice, rather
than a closed nearest-neighbor equation. For the stationary cosine
control with every gap equal to $g$, each side of
\eqref{lat:all-gap} is
$g^{-2}\sum_{m\ge1}[m(m+1)]^{-1}=g^{-2}$. Hence
$M_n(1,1)=2/g^2$ and $g_n'=0$, checking the physical normalization.
This noncanonical control does not model the canonical population.

If only a finite consecutive real block is known, use
\eqref{lat:moments} with the remaining complete divisor, which may
be nonreal. The part of \eqref{lat:all-gap} belonging to known real
zeros is still valid, but the omitted contribution must be retained
as a complete exterior interaction. Nonreal groups cannot be replaced
by the positive real-gap terms of \eqref{lat:all-gap}. The row
identities and the compatibility statements below remain valid on
cleared charts of the real block in this setting.

\subsection{Overlap identities and a common meromorphic field}

\begin{proposition}[Adjacent-row compatibility]\label{lat:compat-prop}
For every integer $r\ge2$,
\begin{equation}
 M_n(r,0)+(-g_n)^{-r}
 =M_{n+1}(0,r)+g_{n+1}^{-r}.
 \label{lat:overlap}
\end{equation}
Within a row, for $p,q\ge1$ and $p+q\ge3$,
\begin{equation}
 g_nM_n(p,q)=M_n(p,q-1)-M_n(p-1,q).
 \label{lat:row-algebra}
\end{equation}
These are necessary constraints on rows arising from a common divisor.
They are not asserted to characterize all realizable moment arrays.
\end{proposition}

\begin{proof}
Both sides of \eqref{lat:overlap} equal
$\sum_{k\ne n+1}(x_k-x_{n+1})^{-r}$, with any nonreal exterior
occurrences also included. The left side restores $x_n$ to the
sum for edge $n$, and the right side restores $x_{n+2}$ to the sum
for edge $n+1$. All terms are absolute for $r\ge2$.
For \eqref{lat:row-algebra}, set
$X=\rho-x_{n+1}$, $Y=\rho-x_n=X+g_n$. The identity
\[
 X^{-p}Y^{-(q-1)}-X^{-(p-1)}Y^{-q}
 =(Y-X)X^{-p}Y^{-q}
\]
can be summed absolutely in the stated range. At $(p,q)=(1,1)$,
the same identity holds only as the common-cutoff difference of
the first-power sums: its combined summand is $g_n/(XY)$ and
decays quadratically. Separate absolute first-power sums are neither
required nor introduced.
\end{proof}

The common generating field is $B(t,z)=H_z(t,z)/H(t,z)$.
Removing the two target poles of edge $n$ gives
\begin{equation}
 R_n(t,z)=B(t,z)-\frac1{z-x_n(t)}-\frac1{z-x_{n+1}(t)}.
 \label{lat:resolvent}
\end{equation}
The removed singularities are removable at these simple zeros.
For three consecutive distinct labels, direct subtraction gives
\begin{equation}
 R_{n+1}(t,z)=R_n(t,z)+\frac1{z-x_n(t)}
                         -\frac1{z-x_{n+2}(t)}.
 \label{lat:resolvent-shift}
\end{equation}
This identity requires no reality assumption on the other zeros.
The tail conventions are those of the same complete field $B$ on
both sides; there is no independently chosen additive regularizer.

The Taylor coefficients of $R_n$ at the two targets recover the
pure-axis moments. For $r\ge2$,
\begin{equation}
 \frac{\partial_z^{r-1}R_n(t,x_{n+1})}{(r-1)!}
       =-M_n(r,0),\qquad
 \frac{\partial_z^{r-1}R_n(t,x_n)}{(r-1)!}
       =-M_n(0,r).
 \label{lat:axis-taylor}
\end{equation}
Indeed, differentiation of an exterior simple pole gives
$(-1)^{r-1}(z-\rho)^{-r}$; at the target this is
$-(\rho-z)^{-r}$. The derivatives remove the genus-one constant
regularizers, and the remaining series is absolute.

\begin{proposition}[Mixed generating function]\label{lat:generator-prop}
Write $a=x_n$, $b=x_{n+1}$ at a fixed time. In sufficiently small
cleared neighborhoods of the two targets,
\begin{equation}
 -\frac{R_n(b+u)-R_n(a+v)}{b-a+u-v}
 =\sum_{p,q\ge1}M_n(p,q)u^{p-1}v^{q-1}.
 \label{lat:generator}
\end{equation}
The quotient has its divided-difference interpretation when its
two arguments coincide within the analytic domain of $R_n$.
Moreover, off the zeros,
\begin{equation}
 B_t=-B_{zz}-2BB_z.
 \label{lat:burgers}
\end{equation}
\end{proposition}

\begin{proof}
For one exterior pole,
\[
 -\frac{(b+u-\rho)^{-1}-(a+v-\rho)^{-1}}
        {b-a+u-v}
 =\frac1{(\rho-b-u)(\rho-a-v)}.
\]
The constant regularizers cancel in the difference. Choose the
neighborhoods so that $|u|/|\rho-b|$ and $|v|/|\rho-a|$ are
uniformly bounded by a number less than one for every exterior
occurrence. Expanding the two geometric series gives
\eqref{lat:generator}. Their absolute sum is bounded by a constant
times $\sum_\rho|\rho-b|^{-1}|\rho-a|^{-1}$, which is finite by
the fixed-order tail estimate and local clearance. This justifies
the interchange of the two power series and the complete divisor.
If the arguments coincide, the left side is $-R_n'$ by analytic
continuation of the divided difference. Finally,
$H_t/H=-(B_z+B^2)$ follows from the heat equation. Differentiate
in $z$ to obtain \eqref{lat:burgers}.
\end{proof}

Equations \eqref{lat:resolvent-shift}--\eqref{lat:burgers} express
all rows through one field. A Taylor translation cannot cross an
unremoved neighboring pole: use overlapping cleared neighborhoods
or remove the entire chosen block first. The meromorphic equation
is another exact expression of the original heat evolution, not an
additional estimate on its canonical solution. An arbitrary array
of individually evolved or truncated moment rows need not satisfy
these overlap identities or arise from any divisor.

\subsection{Node forces, block cancellation and its error budget}

At a simple root define the regular part
\begin{equation}
 U_n(t)=\lim_{z\to x_n}
       \left(B(t,z)-\frac1{z-x_n}\right)
       =\frac{H_{zz}(t,x_n)}{2H_z(t,x_n)}.
 \label{lat:node-force}
\end{equation}
Taylor expansion of $H$ at $x_n$ gives the second equality. Implicit
differentiation of $H(t,x_n(t))=0$ and $H_t=-H_{zz}$ gives
$x_n'=2U_n$. This definition uses a finite part of the complete
source, not an unregularized infinite first-power sum.

\begin{proposition}[Block flux identity]\label{lat:flux-prop}
For a fixed finite block of edges $n=a,\ldots,b$ on a common
simple-root interval,
\begin{align}
 g_nM_n(1,1)&=\frac2{g_n}-(U_{n+1}-U_n),
                         \label{lat:edge-flux}\\
 \sum_{n=a}^{b}g_nM_n(1,1)
 &=2\sum_{n=a}^{b}\frac1{g_n}-(U_{b+1}-U_a).
                         \label{lat:block-flux}
\end{align}
The same formulas hold for a finite real block with nonreal exterior,
provided the complete interactions and source finite parts are used.
\end{proposition}

\begin{proof}
The node law gives $g_n'=2(U_{n+1}-U_n)$. Comparison with
\eqref{lat:gap-law} proves \eqref{lat:edge-flux}. Summing the
successive differences telescopes exactly, giving
\eqref{lat:block-flux}. No positivity of exterior summands enters
these two operations.
\end{proof}

The cancellation has a precise, limited consequence for error
propagation. Put $A_{a,b}=\sum_{n=a}^b g_nM_n(1,1)$. If the
gaps are exact and common approximations to the endpoint forces
satisfy $|\widehat U_j-U_j|\le\varepsilon_j$ for
$j=a,b+1$, reconstruct the aggregate by the right side of
\eqref{lat:block-flux}. Then
\begin{equation}
 |\widehat A_{a,b}-A_{a,b}|
 \le\varepsilon_a+\varepsilon_{b+1}.
 \label{lat:exact-gap-error}
\end{equation}
For a uniform node error $\varepsilon$, this is $2\varepsilon$,
independent of the number of edges. Bounding each edge-force
difference separately and summing absolute allowances would pay
up to $2(b-a+1)\varepsilon$. The smaller aggregate allowance uses
the same approximated node at each shared endpoint, so its interior
errors cancel before absolute values are taken.

If the gaps are also approximated, assume $\widehat g_n>0$ and
$g_n>0$, and define
$\widehat A_{a,b}=2\sum_{n=a}^b\widehat g_n^{-1}
-(\widehat U_{b+1}-\widehat U_a)$. The full bound is
\begin{align}
 |\widehat A_{a,b}-A_{a,b}|
 &\le\varepsilon_a+\varepsilon_{b+1}
       +2\sum_{n=a}^b
          \frac{|\widehat g_n-g_n|}{\widehat g_n g_n}
                                                        \notag\\
 &\le\varepsilon_a+\varepsilon_{b+1}
       +2\sum_{n=a}^b
          \frac{|\widehat g_n-g_n|}
               {\min(\widehat g_n,g_n)^2}.
 \label{lat:uncertain-gap-error}
\end{align}
This follows from the exact difference of two reciprocals. Certified
positive lower gap bounds may replace the denominators in the final
allowance. Small gaps can make this payment large. These bounds
concern an aggregate reconstructed by the identity; independently
approximating each product $g_nM_n(1,1)$ requires its additional
weight and moment errors.

The same-information direct comparator has exactly the same
cancellation. The block length satisfies
\begin{equation}
 \frac{\mathrm d}{\mathrm dt}(x_{b+1}-x_a)
       =2(U_{b+1}-U_a)
       =4\sum_{n=a}^b\frac1{g_n}-2A_{a,b}.
 \label{lat:block-length}
\end{equation}
Thus \eqref{lat:exact-gap-error} is a correlated-error advantage
over an independently summed allowance, not a demonstrated gain
over direct common-source evolution. A block can preserve its total
length while one interior gap decreases to zero. Neither
\eqref{lat:block-flux} nor its endpoint-error bound protects the
regional minimum. Such protection would need a further estimate
that prevents concentration of the defect into one edge.

\subsection{Compatible reconstruction and the collision boundary}

One finite polynomial heat field produces all its roots and moment
rows at once. Before the relevant collisions, its rows satisfy the
gap, overlap, resolvent and block identities above exactly. This is
the compatibility supplied by the shared finite-divisor construction
of \cite{BakerCM15}, developed further in Section~\ref{sec:finite}.
Its higher moments are those of the same finite divisor; they are
not independently set to zero. Independent truncation in $(p,q)$
at each edge is therefore a different construction, generally
without a common underlying divisor.

Compatibility alone does not establish approximation to the
canonical field. Sections~\ref{sec:common-time}--\ref{sec:counted-transfer}
prove a common-time finite-product theorem with one initialization and
explicit field, jet and fixed-moment errors. A finite-block application
must supply a compact set containing every relevant contour, a common
field error, matching labels and the minimum clearance across that block.
The pairwise transfer then applies with all these constants paid,
provided every retained pair satisfies the chart and strong-collar
hypotheses of \eqref{eq:domain} and Section~\ref{sec:counted-transfer}. No uniformity
as the block size, height or moment order grows is included. An
index cutoff also leaves complete nonlocal exterior forces and
boundary data; it cannot be implemented by dropping the omitted
zeros without an error estimate.

The adjacent-edge charts have a further limitation at mergers.
If $x_n=x_{n+1}$, the symmetric $S/T$ coordinates of that target
pair can remain regular, as in Section~\ref{sec:integer}. For the
neighboring edge $(x_{n-1},x_n)$, however, the exterior occurrence
$x_{n+1}$ meets a target; the analogous failure occurs for
$(x_{n+1},x_{n+2})$. Their exterior separation is lost. It is
therefore incorrect to infer that every adjacent split-pole chart
extends regularly through the same merger. A shared polynomial,
or suitable larger-cluster symmetric coefficients, can remain
regular while these individual charts must stop or be replaced.

Physical quadratic and cosine-plus-constant collision controls
satisfy the applicable simple-root identities before collision.
The lattice representation does not exclude those events. Its
role is to preserve the common-divisor constraints and expose
correlations such as \eqref{lat:block-flux}; it supplies neither
a universal positivity principle nor stability of a finite
closure. The next quantitative question is a common-source
estimate controlling concentration into individual gaps over a
specified physical interval, with complete exterior and boundary
payments and the direct evolution given the same information.

\section{Finite heads, boundary residuals and compatible approximation}
\label{sec:finite}

An infinite hierarchy is useful for approximation only when the error
introduced at its boundary can be transported to the observables of
interest. We now make that boundary explicit. This also separates two
notions of finite approximation: truncating the list of moments, and
evolving a finite divisor whose moments remain mutually compatible.
The distinction was motivated by the earlier hierarchy discussion
\cite{BakerCM15}; the calculations below are consequences of
Theorem~\ref{thm:hierarchy}.

\subsection{The four boundary rows}

Freeze the scale $\ell$ from \eqref{eq:domain} and put
\begin{equation}
 \sigma=(t-t_0)/\ell^2,\qquad \delta=D/\ell^2,\qquad
 s_n=\ell^{2n}S_n,\qquad v_n=\ell^{2n-1}T_n,\qquad m=s_1.
 \label{fin:normalization}
\end{equation}
The equations \eqref{eq:S-hierarchy}--\eqref{eq:T-hierarchy} have the
same coefficients after replacing $(D,S,T,t)$ by $(\delta,s,v,\sigma)$.
Fix $P\ge3$ and retain
$x_P=(\delta,s_1,\ldots,s_P,v_2,\ldots,v_P)$.
Define the polynomial vector field $f_P$ by setting all indices above
$P$ to zero in those equations. This defines an approximation, not an
identity for the actual higher moments.

\begin{proposition}[Boundary residual]\label{fin:residual}
The actual canonical head satisfies $x_P'=f_P(x_P)+r_P$, where primes
in this section denote $\sigma$ derivatives. Only four residual rows
can be nonzero:
\begin{align}
 r_{s,P-1}&=-4(P-1)(P+2)\delta s_{P+1},\notag\\
 r_{s,P}&=-4P(P+2)s_{P+1}-4P(P+3)\delta s_{P+2},\notag\\
 r_{v,P-1}&=-4(P-1)(P+2)\delta v_{P+1},\notag\\
 r_{v,P}&=-4(P^2+P-1)v_{P+1}-4P(P+3)\delta v_{P+2}.
 \label{fin:four-rows}
\end{align}
With $\gamma=q_0/a^2<1$ and $B_1=14\Cu\gamma$,
\begin{equation}
 \|r_P\|_\infty\le
 4P(P+3)\sqrt{q_0}B_1L\gamma^{P-1/2}
             \max\{d,1+d\gamma\}.
 \label{fin:residual-bound}
\end{equation}
In particular the residual tends to zero geometrically at fixed height
and collar parameters.
\end{proposition}

\begin{proof}
Every nonlinear factor in row $n$ has index at most $n$. The only
indices larger than $P$ therefore occur in the two linear higher-moment
terms of rows $P-1$ and $P$, giving \eqref{fin:four-rows}.
The shell bound used below in \eqref{eq:integer-bound} gives
$\mathcal A_1=\sum_j|\mu_j|\le B_1L$; it follows directly from
\eqref{eq:count-safe} and \eqref{eq:collar-kernel}. Thus
\[
 |s_n|\le B_1L\gamma^{n-1},\qquad
 |v_n|\le\sqrt{q_0}B_1L\gamma^{n-3/2}\quad(n\ge2).
\]
For the odd estimate use
$|w_j|/\sqrt{|Q_j|}\le\sqrt{q_0}$ and sum
$|\mu_j|^{n-1/2}$. Apply these inequalities to the four rows,
using $|\delta|\le d$, $P^2+P-1\le P(P+3)$ and $\gamma<1$.
The $P-1$ rows contribute the separate factor $d$ in
\eqref{fin:residual-bound}; they have no extra factor $\gamma$.
\end{proof}

\subsection{What an independent absolute error allowance loses}

Suppose $y'=f_P(y)$ and a convex box contains both $x_P$ and $y$
throughout $[0,\Sigma]$. If $K_P$ bounds the induced max-norm Jacobian
of $f_P$ on that box and $\epsilon_P$ is the right side of
\eqref{fin:residual-bound}, the usual differential inequality gives
\begin{equation}
 \|x_P(\tau)-y(\tau)\|_\infty
 \le \e^{K_P\tau}\|x_P(0)-y(0)\|_\infty
       +\epsilon_P\frac{\e^{K_P\tau}-1}{K_P}.
 \label{fin:gronwall}
\end{equation}
At $K_P=0$ the quotient means $\tau$. The containing box is a
premise of this bound. In the present unweighted coordinates,
\[
 \left|\frac{\partial (f_P)_{s,P-1}}{\partial s_P}\right|
       =4(P^2-1),
\]
so a scalar Jacobian majorant grows at least quadratically with $P$.
Combining that particular allowance with a geometrically small boundary
residual does not prove convergence on a fixed positive duration.
A componentwise version makes the source of this loss more precise.

\begin{proposition}[A scoped positive-comparison limitation]
\label{fin:positive-box}
Let $A_P$ be a constant entrywise nonnegative absolute-Jacobian majorant
on the containing box. Give its residual majorant an independent constant
component
\[
 \bar r_{s,P}=4P(P+2)B_1L\gamma^P
             +4P(P+3)dB_1L\gamma^{P+1},
\]
with all other components nonnegative. The comparison solution with
zero initial error,
$z(\tau)=\int_0^\tau\exp((\tau-u)A_P)\bar r\dd u$, satisfies
\begin{equation}
 z_m(\tau)\ge\frac{B_1L}{2}(P+2)(P+1)!(4\gamma\tau)^P.
 \label{fin:path-lower}
\end{equation}
For fixed $\tau>0$, $L$ and $\gamma>0$, this allowance diverges as
$P\to\infty$. The conclusion concerns this independent constant-tail,
absolute-Jacobian comparison, not the signed truncation error.
\end{proposition}

\begin{proof}
Row $s_n$ has derivative $-4n(n+2)$ with respect to $s_{n+1}$;
its nonlinear terms do not contain that component. Hence the
successive entries of $A_P$ along the path
$s_P\to s_{P-1}\to\cdots\to s_1$ are at least $4n(n+2)$.
Their product is
\[
 \prod_{n=1}^{P-1}4n(n+2)
       =4^{P-1}(P-1)!(P+1)!/2.
\]
Every matrix power and every component of $\bar r$ is nonnegative.
Keep this one path in the $(P-1)$st term of the exponential series
and the first term of $\bar r_{s,P}$. Integrating
$(\tau-u)^{P-1}/(P-1)!$ gives $\tau^P/P!$, proving
\eqref{fin:path-lower}. The ratio of successive factorial-exponential
terms tends to infinity, proving the stated divergence.
\end{proof}

The discarded signs and compatibility relations could matter to the
actual solution. Proposition~\ref{fin:positive-box} neither proves
instability of a moment truncation nor rules out a compatible closure.
It explains why a small local residual is not by itself a useful
fixed-duration error theorem. Shorter steps also need a controlled
restart and accumulation argument.

\subsection{Preserving the divisor and the intended error consumer}

A different finite object is obtained by truncating the complete
opposite-paired product at an initial time $t_0$, retaining conjugates,
ties and multiplicities, and evolving the resulting polynomial $P_N$
by the heat equation. Its higher moments are then determined by one
actual divisor; they are not set independently to zero. At the purely
polynomial level the evolution has the exact expression
\begin{equation}
 P_N(t_0+\tau,z)
 =\exp(-\tau\partial_z^2)P_N(t_0,z)
 =\mathbb E\,P_N(t_0,z+\ii\sqrt{2\tau}\,X),\qquad \tau\ge0,
 \label{fin:polynomial-heat}
\end{equation}
where $X$ is standard normal. To verify the second identity, expand the
polynomial in its finite Taylor series. Odd Gaussian moments vanish,
and $\mathbb EX^{2k}=(2k)!/(2^kk!)$ converts the even terms exactly
to $(-\tau)^k\partial_z^{2k}P_N/k!$.

The source in \eqref{eq:source} gives the same Gaussian formula for $H$:
expectation multiplies $\cos(zu)$ by $\exp(\tau u^2)$. Its
super-exponential decay justifies Fubini for bounded real $\tau\ge0$
and bounded complex $z$. What remains unproved here is a common-time
error estimate for $P_N-H$ and the transfer to the selected pair and
moments. The slice-jet passage used in Theorem~\ref{thm:hierarchy}
reinitializes at each reference time and does not supply that estimate.
A possible route is a common product growth envelope, the complete
omitted inverse-square tail $O(\log N/N)$, Gaussian domination, and
Cauchy and Rouch\'e estimates on the independently cleared tube.
This is a specified approximation problem, not a theorem of the present
paper. The cutoff $N$ is a divisor cutoff, distinct from the moment
order $P$; initial canonical divisor data and its tail must be supplied.

The error needed by transport is simple to state. Let $m_N$ be real,
as for a conjugation-compatible approximation. If it satisfies $|m-m_N|\le\eta$ on a scaled interval of length $\Sigma$, then
for every subinterval $[x,y]$,
\begin{equation}
 \left|\log\frac{W(x,y)}{W_N(x,y)}\right|
 \le4(y-x)\eta\le4\Sigma\eta,
 \qquad W(x,y)=\exp\left(-4\int_x^y m\dd\sigma\right).
 \label{fin:transport-error}
\end{equation}
This follows by subtracting the two integrals. The same uniform
subinterval control bounds the inhomogeneous cost multiplicatively;
the physical cost retains its factor $\ell^2$. Forcing error must also
retain the product in $12s_2+16\delta s_3$. None of these consequences
supplies the independent tube or gap anchors.

The analogy with BBGKY truncation is therefore organizational: it asks
how unresolved information enters a retained hierarchy. The older
Jacobi-chain memory variable was observation height, whereas the present
variable is Newman time; the omitted moments here generally have nonzero
initial data. Inter-level compatibility does not by itself imply stability
or physical realizability, and an exact projection-memory formalism needs
a specified state space and flow \cite{AkbariEtAl2012Hierarchy,ChorinHaldKupferman2000}.
No density-matrix positivity, independence law or arbitrary-sequence
well-posedness is assumed for these divisor moments.

\section{Complete complex powers on local time charts}
\label{sec:complex}

The rational moments are regular at a merger. To extend the denominator
power to a complex index, the logarithm must be chosen on a spatial
neighborhood of each coalescing cluster. The two infinite ends of the
divisor admit a common principal branch. We use the stated domain
\eqref{eq:domain}, the strip and estimates \eqref{eq:count-safe}, and the
frozen \(\ell\) throughout this section \cite{BakerBase,BakerGT41,BakerCM03}.

Use \(w_j=\rho_j-c\), \(Q_j=Q(t,w_j)\) and
\(\mu_j=\ell^2/Q_j\) from \eqref{eq:quadratic}, with every exterior
occurrence counted. The collar gives \(|Q_j|\ge(15/16)|w_j|^2\).
For a specified spatial logarithm of \(Q(t,w)\), put
\[
 f_p(t,w)=\exp\!\big(p(2\log\ell-\Log Q(t,w))\big),\qquad
 h(t,w)=\frac{H_z(t,c(t)+w)}{H(t,c(t)+w)} .
\]
Here \(\log\ell\) is real. In a time chart, define
\(Z(p,t)=\sum_j f_p(t,w_j)\). At positive integers \(n\) this agrees
with \(\ell^{2n} S_n(t)\).

\subsection{Spatial contours and the time law}

At each reference time choose left and right centered cuts
\(-R_-\in[-7,-6]\) and \(R_+\in[6,7]\) with no zero on their
vertical connectors \(|\Im w|\le4\). These cuts remain cleared on
a smaller real-time interval. Appendix~\ref{app:clearance} proves the
choice from the unit count and gives the bounds needed to remove remote
vertical closures. On the two closed half-strips beyond the cuts,
\(\Re Q\ge(\Re w)^2/2\), so we use the principal log. This defines
the power at every exterior occurrence beyond either cut.

Inside the central rectangle, enclose each distinct reference
exterior location in a small disk. The disks have disjoint closures
and avoid the target and \(Q=0\). Their boundaries \(\gamma_C\) are
zero-free. The argument principle preserves each count and the total
central count on a smaller interval, so the disks account for all
central exterior occurrences. At a disk center \(w_*\) and reference
time \(t_*\) choose \(L_*\) with \(\exp L_*=Q_*=Q(t_*,w_*)\) and
\(\Im L_*\in[-\pi,\pi]\). On the disk use
\begin{equation}
 \Log Q=L_*+\Log_{\rm pr}(Q/Q_*),\qquad |Q/Q_*-1|<1/2 .
 \label{cp:local-log}
\end{equation}
Shrinking the disks and the time interval ensures the inequality on
neighborhoods of their closures. These logs have
\(|\Im\Log Q|\le\Theta=3\pi/2\). Each disk uses one log through any
merger of its counted occurrences. The compact established interval admits a
finite cover by such charts, with estimates on compact subintervals
that still cover it.

For an integrand \(B\) with the log prescribed on each piece, let
\begin{equation}
\begin{aligned}
 \mathcal S(B)=\frac1{2\pi\ii}\Bigg(&
   \sum_C\oint_{\gamma_C} B(w)\dd w
   +\int_{R_+}^{\infty}\!\big[B(x-4\ii)-B(x+4\ii)\big]\dd x\\
 &+\int_{-\infty}^{-R_-}\!\big[B(x-4\ii)-B(x+4\ii)\big]\dd x\\
 &+\int_{R_++4\ii}^{R_+-4\ii} B(w)\dd w
  +\int_{-R_--4\ii}^{-R_-+4\ii} B(w)\dd w\Bigg).
\end{aligned}
\label{cp:spatial-boundary}
\end{equation}
The circles are counterclockwise. Both lower rays go rightward and
both upper rays leftward. The finite right connector goes down, and
the finite left connector goes up.

\begin{theorem}[Complete local complex-power law]\label{thm:complex-time}
On each chart the complete sum \(Z\) is holomorphic in \(p\) for
\(\Re p>1/2\) and is \(C^1\) in real physical time, with inward
derivatives at the interval endpoints. With derivatives at fixed \(w\),
\begin{equation}
\begin{aligned}
 Z&=\mathcal S(f_ph),\\
 Z_t&=\mathcal S((f_p)_t h+f_ph_t)
   =\mathcal S\big(\big[(f_p)_t-(f_p)_{ww}-c'(f_p)_w\big]h
                                       +(f_p)_wh^2\big),\\
 h_t&=-h_{ww}-2h h_w+c' h_w,\qquad (f_p)_t=pD'f_p/Q .
\end{aligned}
\label{eq:complex-time}
\end{equation}
The derivative is holomorphic on the same half-plane and continuous
through target and exterior mergers on the chart.
\end{theorem}

\begin{proof}
The complete value bound in Appendix~\ref{app:value-sums} gives normal
convergence in \(p\). For the spatial formula, first close the two
half-strips at distant cleared verticals, upward on the right and
downward on the left. The residue of \(h\) is the zero multiplicity,
and \(f_p\) is holomorphic in the region of each contour. Hence the
finite contours count all their exterior occurrences. The remote
edges vanish by \eqref{app:edge-bound}. The omitted divisor values
tend to zero and the horizontal rays converge absolutely.

The centered heat equation yields
\(h_t=\partial_w(-H_{zz}/H)+c'h_w\). Substituting
\(H_{zz}/H=h_w+h^2\) gives the last line of
\eqref{eq:complex-time}. Differentiation under \(\mathcal S\) is
justified independently of divisor truncation by the safe-line and
finite-contour estimates in Appendix~\ref{app:time-bounds}. For the
second derivative formula, the difference of the two integrands is
the spatial derivative of
\[
 -f_ph_w+(f_p)_wh-f_ph^2+c'f_ph .
\]
The integral of this derivative vanishes on each circle. On each
tail boundary the primitive agrees where the rays meet the finite
connector and tends to zero at the infinite ends. The square of
the full field \(h\) in \eqref{eq:complex-time} retains the
spatial interaction in the integer law \eqref{eq:low-orders}.
\end{proof}

Two time charts use the same principal logs for \(|\Re w|\ge7\).
Near each overlap time, refine the finitely many other assignments
by small disks around the distinct affected locations, with fixed
counts. In each refinement disk the two logs differ by a constant
\(2\pi\ii k_C\), \(k_C\in\mathbb{Z}\). The branch difference is
\begin{equation}
 Z_B-Z_A=\sum_C\big(\exp(-2\pi\ii k_Cp)-1\big)Z_{A,C},
 \label{cp:transition}
\end{equation}
where \(Z_{A,C}\) is the trace of the power in chart \(A\) over
that disk. Its fixed circle proves validity also after one
physical time differentiation. All these finite corrections
vanish at integers, including their time derivatives. This is a
local atlas for the noninteger powers.

\subsection{Conjugation and the negative-axis defect}

For a function of \(p\), write
\(Z^\sharp(p,t)=\overline{Z(\bar p,t)}\). The field and selected
pair are conjugation invariant at real time, so the central disks
can be chosen in conjugate pairs, with self-conjugate disks at
real centers. For real exterior \(w\) the denominator \(Q\) is
positive. A negative-real \(Q\) can arise only at \(w=\ii y\ne0\)
since \(\Im Q=2\Re w\,\Im w\). Its two conjugate locations are
distinct. Choose opposite arguments \(+\pi,-\pi\) there and take
conjugate \(L_*\) on all paired disks, with the real log at real
centers. In \eqref{cp:local-log} the ratio log has argument of
magnitude less than \(\pi/4\), so \(5\pi/4\) suffices for \(\Theta\)
on these charts. Both far principal logs are conjugation
compatible. Conjugation now permutes the entire occurrence
list and conjugates the logs. Absolute convergence gives
\begin{equation}
 Z^\sharp=Z
 \label{cp:schwarz}
\end{equation}
locally in real time, including mergers. This symmetry also holds
for \(Z_t\) by Theorem~\ref{thm:complex-time}.

At a fixed slice let \(Z_+\) have principal logs off the negative
axis and argument \(+\pi\) at all negative-real \(Q_j\). Its exact
defect is
\begin{equation}
 Z_+-Z_+^\sharp=-2\ii\,\ell^{2p}\sin(\pi p)
                        \sum_{j:Q_j<0}|Q_j|^{-p}.
 \label{cp:axis-defect}
\end{equation}
The sum is finite and includes both conjugate locations with
all multiplicities. This follows by subtracting the two
exponentials \(\exp(-\ii\pi p)\) and \(\exp(+\ii\pi p)\) at each
such occurrence. Resetting to this slice prescription can cross
a branch cut. The time law uses the continuations of the local logs.
At that fixed slice the algebra of \eqref{cp:transition} compares this
prescription with any chart: group the finitely many affected occurrences
by their logarithm shifts. This value comparison does not assert that
resetting to $Z_+$ defines an analytic time chart. The differentiated
transition formula applies to pairs of local analytic charts
\cite{BakerCM04,BakerCMReview}.


\section{Damped Cauchy evolution in the power index}
\label{sec:cauchy}

Fix \(\epsilon>0\). The normalized variables are
\(\delta=D/\ell^2\), \(m=\ell^2M=Z(1,t)\), and
\(\sigma=(t-t_0)/\ell^2\). The two residues contributing to the
forcing are \(12Z(2,t)\) and \(16\delta Z(3,t)\). A separate
contour evaluates \(m\). We treat the two kernels together
\cite{BakerCM05,BakerCMReview}.

Let
\[
 K(p,t)=\sum_{n\in\mathcal N}\frac{c_n(t)\exp(\epsilon(p-n)^2)}{p-n},
 \qquad b\le \Re p\le r ,
\]
where \(\mathcal N\) is a finite set of residue indices. For
\(F\) take \(K=K_\epsilon\), \(\mathcal N=\{2,3\}\),
\(c_2=12,\ c_3=16\delta\) and \((b,r)=(1,4)\). For \(m\) take
\(K=G_\epsilon=\exp(\epsilon(p-1)^2)/(p-1)\) and
\((b,r)=(3/4,3/2)\). If \(\delta=0\), the term at \(3\) is
absent. Set
\[
 \begin{array}{c|c|c}
  K & C_\epsilon & \lambda\\ \hline
  K_\epsilon & (12+16d)\exp(4\epsilon) & 1\\
  G_\epsilon & \exp(\epsilon/4) & 1/4
 \end{array}
\]
so \(|K(q+\ii v,t)|\le C_\epsilon\exp(-\epsilon v^2)/
\sqrt{\lambda^2+v^2}\) for either \(q=b,r\).
Along \(\Im p=\pm Y\) in the strip we have
\begin{equation}
 |K(p,t)|\le C_\epsilon\exp(-\epsilon Y^2)/Y,\qquad Y>0 .
 \label{cau:kernel-bounds}
\end{equation}
Indeed \(\Re (p-n)^2=(\Re p-n)^2-(\Im p)^2\)
and \(|\Re p-n|\ge\lambda\) on either vertical.

For \(V_q(Y)\) upward from \(q-\ii Y\) to \(q+\ii Y\), use
\(V_q\) for its infinite limit and
\[
 \mathcal V_{b,r}(B)=\frac1{2\pi\ii}\left(
                      \int_{V_r} B(p)\dd p-\int_{V_b}B(p)\dd p\right).
\]
The counterclockwise rectangle uses the right upward line minus
the left upward line, the lower edge directed rightward,
and the upper edge directed leftward. Both real edges and
\(\epsilon\) remain fixed.

\begin{theorem}[Complete two-line formulas]\label{thm:damped-cauchy}
Under the hypotheses of Theorem~\ref{thm:complex-time}, one has
\begin{equation}
 \begin{aligned}
 F(t)&=\mathcal V_{1,4}(K_\epsilon Z)
                  =12Z(2,t)+16\delta(t)Z(3,t)
                  =\ell^4(12S_2+16DS_3),\\
 m(t)&=\mathcal V_{3/4,3/2}(G_\epsilon Z)=Z(1,t),\\
 M'&=2M^2-\ell^{-4}F,\qquad
 \frac{\dd m}{\dd\sigma}=2m^2-F,\qquad
 \frac{\dd\delta}{\dd\sigma}=2-4m\delta .
 \end{aligned}
 \label{eq:forcing-contour}
\end{equation}
The line integrals converge absolutely and allow physical time
differentiation under the integral on each compact subchart. The
two-line differences take the same values on chart overlaps.
\end{theorem}

\begin{proof}
The value estimate \eqref{app:value-envelope} gives
\(|Z(p,t)|\le A_b\exp(\Theta|\Im p|)\) on the strip. Multiply by
\eqref{cau:kernel-bounds}. If
\(\Phi_\theta(Y)=\exp(-\epsilon Y^2+\theta Y)\) and
\(\kappa_\theta(Y)=2\epsilon Y-\theta\), the normalized sum of the
two horizontal integrals is bounded by
\((r-b)C_\epsilon A_b\Phi_\Theta(Y)/(\pi Y)\).
For \(Y>\Theta/(2\epsilon)\), the absolute error in the vertical
difference after cutting to \(|\Im p|\le Y\) is bounded by
\begin{equation}
 \frac{2C_\epsilon A_b}{\pi}
            \frac{\Phi_\Theta(Y)}{Y\kappa_\Theta(Y)}.
 \label{eq:vertical-tail}
\end{equation}
This follows from
\(\Phi_\theta(v)\le\Phi_\theta(Y)
\exp(-\kappa_\theta(Y)(v-Y))\) for \(v\ge Y\), bounding
\(1/\sqrt{\lambda^2+v^2}\) by \(1/v\).
Residues on the finite rectangle give the stated values since
the damping equals one at its poles. The horizontal bounds and
vertical tails give the two-line limit. The laws for \(M,m,\delta\)
follow from \eqref{eq:pairlaw} and \eqref{eq:low-orders}, with the
local \(2m^2\) outside the contour.

On a compact chart and fixed real-part strip,
\eqref{app:chart-time} gives
\(|Z_t|\le C_Z(1+|p|)^2\exp(\Theta|\Im p|)\).
The complete physical derivatives of the integrands are
\begin{equation}
 \begin{aligned}
 \partial_t(K_\epsilon Z)&=K_\epsilon Z_t+
          \frac{16(D'/\ell^2)\exp(\epsilon(p-3)^2)}{p-3}Z,\\
 \partial_t(G_\epsilon Z)&=G_\epsilon Z_t .
 \end{aligned}
 \label{cau:kernel-time}
\end{equation}
The kernel bounds and Gaussian moments justify the interchange
and the continuity of the derivative. The exact bounds on
differentiated horizontal edges and vertical tails are in
Appendix~\ref{app:cauchy-errors}.

On an overlap the finite difference \eqref{cp:transition} and
its time derivative vanish at the residue positions.
Consequently \(K(Z_B-Z_A)\) has removable poles, and its
physical derivative has removable poles as well. Their complete
rectangle integrals vanish. Each bound used for \(KZ\), or
its derivative, applies to the difference with the sum of the two
chart constants. Passing to the vertical limits proves equality
of the full right-minus-left values and their derivatives. In
particular the integer evolution glues through the compact time
cover.
\end{proof}

\subsection{Divisor cutoff and independent spatial boundary cutoff}

For integer \(N>4T+8\), let \(Z_N\) retain \(|\rho_j|\le N\)
with all ties and multiplicities. The estimates used to prove the
theorem apply also to \(Z_N\). The strip-wide bound in
Appendix~\ref{app:cauchy-errors} gives absolute convergence in
the line integrals as \(N\to\infty\). The sharper error is the
signed residue of \(Z-Z_N\). If \(F_N=\mathcal V_{1,4}(K_\epsilon Z_N)\)
then
\begin{equation}
 \begin{aligned}
 F-F_N&=12[Z(2)-Z_N(2)]+16\delta[Z(3)-Z_N(3)],\\
 |F-F_N|&\le12R_2(N)+16d R_3(N),\\
 R_n(N)&=2\Cu(\ell^2q_0)^n E_1(2n-1)3^{2n-1}
                               N^{1-2n}\Lambda(N).
 \end{aligned}
 \label{cau:residue-error}
\end{equation}
Here \(q_0=16/15\), \(\Lambda(x)=\log(3+2T+|x|)\), and \(E_1\)
is the explicit dyadic constant of \eqref{app:dyadic}.
This follows by applying the residue argument to the whole
difference before using \eqref{app:divisor-tail} at integer \(n\).
The rates are \(\ell^4N^{-3}\log N\) and \(d\ell^6N^{-5}\log N\)
at fixed \(T\). The separate \(m\) evaluation has error at most
\(R_1(N)\). For truncated vertical integrals add
\eqref{eq:vertical-tail} to these allowances.

For a spatial cutoff \(X\ge7\), truncate only the horizontal rays
of \eqref{cp:spatial-boundary} at \(|\Re w|=X\), retaining the
finite connectors and circles. Its derivative uses
\((f_p)_t h+f_p h_t\) from \eqref{eq:complex-time}. The
integrated-by-parts form on the cut rays carries the primitive
endpoint terms. The independent derivative errors
\eqref{app:spatial-tails} and \eqref{app:weighted-spatial-tail}
pay this cutoff. They are uniform at fixed height on compact
time subcharts. The complete value limit may take \(N\) first
on each finite rectangle and then \(Y\), or use the absolute
bounds with \(N,Y\) tending to infinity together at fixed
\(\epsilon,T\). For the time derivative one may take \(X\) first
on finite rectangles and then \(Y\), or take \(X\to\infty\) in
the full vertical difference. These distinct limits use the
derivative of the spatial boundary formula and the values of
the sharp divisor cutoff.

\subsection{Folded integrals and branch corrections}

For either kernel and real physical time, \(K^\sharp=K\).
On the compatible charts \eqref{cp:schwarz} the weighted
integrand \(B=KZ\) has \(B^\sharp=B\) and
\begin{equation}
 \mathcal V_{b,r}(B)=\frac1\pi\int_0^\infty
              \Re\big[B(r+\ii v)-B(b+\ii v)\big]\dd v .
 \label{cau:fold}
\end{equation}
On finite rectangles the two horizontal integrals still enter.
To fold a full rectangle, split it into the upper arc from \(r\)
to \(b\) and the lower arc from \(b\) to \(r\), including their
horizontal edges. Reflection sends the upper arc to the reversed
lower arc, while conjugation of \(1/(2\pi\ii)\) changes its sign.
Thus for a Schwarz integrand the rectangle integral is twice the
real part of the normalized upper-arc integral. In the defective
case add the normalized lower-arc integral of
\(B_+-B_+^\sharp\).
For the \(+\pi\) slice prescription, the weighted defect of
\(B_+=KZ_+\) is
\[
 \mathcal D_B=B_+-B_+^\sharp
   =-2\ii K(p,t)\ell^{2p}\sin(\pi p)\sum_{j:Q_j<0}|Q_j|^{-p}.
\]
In fact \(B_+(q-\ii v)=\overline{B_+(q+\ii v)}+\mathcal D_B(q-\ii v)\).
For an upward line this gives
\begin{equation}
 \frac1{2\pi\ii}\int_{V_q(Y)} B_+(p)\dd p
 =\frac1{\pi}\int_0^Y\Re B_+(q+\ii v)\dd v
  +\frac1{2\pi}\int_0^Y\mathcal D_B(q-\ii v)\dd v .
 \label{cau:weighted-defect}
\end{equation}
The two-sided contour takes the right-minus-left difference
of both terms. The removable singularities of \(\mathcal D_B\)
at the integer residues give zero integral of the defect on a
complete rectangle. Its horizontal limit is zero by the same
Gaussian bound (the sum of the bounds for \(B_+\) and \(B_+^\sharp\)
suffices). Thus the \emph{full} defect line integrals agree on
the two lines. The lower-half defect integrals in
\eqref{cau:weighted-defect} are kept in a folded formula.


\section{Signed forcing and gap transport}
\label{sec:bounds}

The hierarchy and the complex-power formula describe the evolution
exactly. Their quantitative use is governed by the real forcing
\begin{equation}
 F=12Z(2)+16\delta Z(3),\qquad
 m_\sigma=2m^2-F,\qquad \delta_\sigma=2-4m\delta,
 \label{eq:consumer-system}
\end{equation}
where $m=\ell^2M$, $\delta=D/\ell^2$ and
$\sigma=(t-t_0)/\ell^2$. A lower bound for $m$ requires an
\emph{upper} bound for $F$. Reality of $F$ alone does not provide it.
We first compare the direct and contour estimates at the same information,
then derive the conditional transport bound. The comparison and static
example extend \cite{BakerGR05} as proved in \cite{BakerCM06}.

\subsection{Absolute estimates and preservation of joint signs}

The elementary shell count in \eqref{eq:count-safe} gives, for $n\ge1$,
\begin{equation}
 \mathcal A_n:=\sum_j|\mu_j|^n\le B_nL,
 \qquad B_n=14\Cu\gamma^n,\qquad \gamma=q_0/a^2.
 \label{eq:integer-bound}
\end{equation}
For completeness, the band $|\Re w|<1$ contains at most $4\Cu L$
occurrences, and the shell $k\le|\Re w|<k+1$ contains at most
$2\Cu(2L+\log(k+1))$. The kernel there is bounded by
$q_0\ell^2/k^2$. Using
$\sum k^{-2n}\le2$, $\sum k^{-2n}\log(k+1)\le3$ and $a\le L$
bounds the far shells by $10\Cu\gamma^nL$ and proves
\eqref{eq:integer-bound}.
Since $|\mu_j|\le\gamma$, the direct estimate is
\begin{equation}
 |F|\le (12\gamma+16d\gamma^2)\mathcal A_1\le C_FL,
 \qquad C_F=12B_2+16dB_3.
 \label{eq:direct-forcing}
\end{equation}

Fix the damping $\epsilon>0$ of Section~\ref{sec:cauchy}. Write
$C_\epsilon=(12+16d)\e^{4\epsilon}$ and let $\Theta$ be its branch
argument bound. On the two verticals $q=1,4$,
\[
 |K_\epsilon(q+\ii v,\delta)Z(q+\ii v)|
 \le C_\epsilon\mathcal A_1
        \frac{\e^{-\epsilon v^2+\Theta|v|}}{\sqrt{1+v^2}}.
\]
The absolute two-line envelope is therefore bounded by
\begin{equation}
 \frac1{2\pi}\int_\R
       (|K_\epsilon Z(4+\ii v)|+|K_\epsilon Z(1+\ii v)|)\dd v
 \le C_{\rm env}\mathcal A_1,
 \quad C_{\rm env}=\frac{2C_\epsilon}{\pi}
       \sqrt{\frac\pi\epsilon}\e^{\Theta^2/(4\epsilon)}.
 \label{eq:contour-envelope}
\end{equation}
Completing the square proves this estimate. At fixed collar parameters
$a,d$ and fixed damping $\epsilon$, its height order is the same as
\eqref{eq:direct-forcing}, and its displayed coefficient is
larger: minimizing $\e^{4\epsilon}/\sqrt\epsilon$ gives
$C_{\rm env}/(12+16d)\ge2\sqrt{8\e/\pi}>1$, whereas
$12\gamma+16d\gamma^2<12+16d$.
This comparison concerns these explicit envelopes, not every possible
signed contour estimate.

\begin{proposition}[Comparison for a fixed grouping]\label{prop:group-bound}
Partition the exterior occurrences into prescribed finite groups $G$,
and retain the assigned logarithms. Put
\[
 Z_G(p)=\sum_{j\in G}\mu_j^p,\quad
 I_{q,G}(v)=K_\epsilon(q+\ii v,\delta)Z_G(q+\ii v),\quad
 F_G=\sum_{j\in G}(12\mu_j^2+16\delta\mu_j^3),
\]
where the power notation uses those logarithms. Then
\begin{equation}
 F_G=\frac1{2\pi}\int_\R(I_{4,G}-I_{1,G})\dd v,
 \qquad
 |F_G|\le\frac1{2\pi}\int_\R|I_{4,G}-I_{1,G}|\dd v
 \le\frac1{2\pi}\int_\R(|I_{4,G}|+|I_{1,G}|)\dd v.
 \label{eq:group-bound}
\end{equation}
All group sums and integrals in these estimates interchange absolutely.
Thus estimating each group by an absolute contour norm costs at least
$\sum_G|F_G|$, the direct signed-residue norm for the same grouping.
For groups closed under conjugation as occurrence multisets, with compatible
logs, the same statement holds for the folded
norm $\pi^{-1}\int_0^\infty|\Re(I_{4,G}-I_{1,G})|\dd v$.
\end{proposition}

\begin{proof}
Apply the positive-rectangle residue formula to the finite entire
sum $Z_G$. The horizontal edges tend to zero by Gaussian damping.
The residues at $2$ and $3$ are independent of the assigned logarithms.
The triangle inequalities give \eqref{eq:group-bound}. On either
vertical, the sum over groups of the absolute integrands is bounded
by the integrable majorant preceding \eqref{eq:contour-envelope}.
The corresponding horizontal group sum is bounded by a constant times
$\e^{-\epsilon Y^2+\Theta Y}/Y$. This proves all complete-divisor
passages. For a compatible group, conjugation changes $v$ to $-v$,
giving the folded expression and the same triangle inequality.
\end{proof}

Proposition~\ref{prop:group-bound} leaves joint signs between groups
available on either side of the identity. Discarding them in the direct
formula but retaining them in the contour formula would compare
different information. The sharp zero-cutoff residue error in
Section~\ref{sec:cauchy} likewise belongs to both evaluations.

\subsection{A complete static example with linear forcing}

The following construction tests what the slice hypotheses can force.
It changes the canonical divisor and is not an example of its heat
evolution. The explicit constants ensure that the class is fixed as
$T$ grows.

\begin{proposition}[Scope of a static forcing bound]\label{prop:static}
There is a family of real-even entire functions of order one with a
counted double target, collar $a=4,d=1$, width-one strip, growing zero
count, unit-count and static safe-line bounds with fixed enlarged
constants, whose complete complex-power sums are holomorphic on
$\Re p>1/2$, have compatible Schwarz symmetry and admit the fixed-damping
contours of Section~\ref{sec:cauchy}, but satisfy
\begin{equation}
 m=-\frac1{32}+O(L^{-1}),\qquad \frac FL\longrightarrow\frac3{256}>0.
 \label{eq:static-asymptotic}
\end{equation}
In particular, this static class admits no uniform upper bound $F=o(L)$,
even on the subclass with $m\le-1/64$.
\end{proposition}

\begin{proof}
Take $c=T=\e^L$, $L\ge16$, $D=0$, $b=8$ and $J=\lfloor L\rfloor$.
For this separate slice construction use the inherited strip, unit-count
and safe-line bounds of the canonical field at time zero. The background
bounds hold on the closed slab $[0,1/2]$ independently of whether the
particular pair tube $I$ contains zero \cite{BakerBase,BakerGR05}.
Remove from $H(0,z)$ every occurrence in the closed bands
$|\Re\rho\mp c|\le3$ by division by their complete monic polynomial.
Insert at $+c$ the double target and $J$ occurrences at each
displacement $\pm b\ell$, and $J+1$ at each displacement
$\pm\ii b\ell$. Reflect all insertions at $-c$. The division is
entire, and conjugation and negation invariance make the resulting
function real and even. Finite polynomial multiplication and division
preserve the order bound; the retained growing count forces order one.
The strip and the required collar hold, with the remaining background
at horizontal distance greater than $3$.

The inherited canonical growing count and its logarithmic discrepancy
are the static inputs described in \cite[Section 3]{BakerGR05}.
There are at most $12\Cu L$ removed occurrences per band and at most
$4J+4\le5L$ insertions per band. Hence the unit-count constant can
be enlarged to $\Cu+5$, and the cumulative-count discrepancy constant
by $24\Cu+10$. The logarithmic derivative changes by a signed sum
of at most $(24\Cu+10)L$ reciprocal distances. On the line $v+4\ii$,
if $|v|\ge c/2$ use vertical distance at least $3$ and
$\log(2+|v|)\ge L/2$; otherwise use horizontal distance at least
$c/3$. The safe-line constant can be enlarged to
$\Cu+5(24\Cu+10)$, uniformly in $c$.
The corresponding lower safe line follows by conjugation.

For the near core assign conjugate logarithms $\pm\pi$ to the negative
values of $Q_j$ and real logarithms to the positive ones. Its complete
contribution is
\begin{equation}
 Z_{\rm core}(p)=b^{-2p}\{2J+2(J+1)\cos(\pi p)\}.
 \label{eq:static-core}
\end{equation}
All remaining occurrences are retained in a remainder $R$. The retained
background and every opposite-band insertion, including the reflected
double, satisfy $|\Re w|>3$, $|\Im w|\le1$. Hence
$\Re(w^2)>8$, and we assign principal logs on the entire noncore.
These are conjugation compatible and have argument magnitude less than
$\pi/2$. Unit-shell counting, including the opposite insertions and
opposite target, gives
for $1\le u\le r<\infty$
\begin{equation}
 |R(u+\ii v)|\le (10\Cu+5)L^{1-2u}\e^{\pi|v|/2}.
 \label{eq:static-tail}
\end{equation}
Indeed the physical absolute power sum of exponent $2u$ over the
retained background is at most $10\Cu L$; the opposite insertions
cost at most $5Lc^{-2u}$. Multiplication by $\ell^{2u}$ proves the
bound. The same shell argument for any $u>1/2$ proves normal
convergence on compact subsets of that half-plane. Compatible logs
give Schwarz symmetry, and Gaussian damping pays the complete vertical
integrals and vanishing horizontal edges. Evaluating
\eqref{eq:static-core}--\eqref{eq:static-tail} at $1$ and $2$ gives
\[
 m=-2/b^2+R(1),\qquad
 F=12(4J+2)/b^4+12R(2),
\]
which proves \eqref{eq:static-asymptotic} and the stated exclusion.
\end{proof}

The same construction retains the inherited broad transverse upper
bounds with fixed enlarged constants \cite[Section 3]{BakerGR05}.
If their original coefficient is $B_\beta$ for $0<\beta<1/2$, the
added nonreal occurrences lie
at height $8/L$ and number at most $3L$ per band. For every dyadic
interval $[U,2U]$ meeting the positive inserted nonreal projection,
$c\in[U,2U]$, including endpoint cuts, one has
$L/2\le\log U\le2L$. The comparison
$U\log U\,\e^{-\beta v\log U}\ge4L\e^{-16\beta}$ whenever
$v<8/L$ allows coefficient $B_\beta+\e^{16\beta}$; deletion only
reduces the counts. This is still a slice construction. It supplies
neither the unchanged theta source nor a common independently established
canonical time tube. In particular Proposition~\ref{prop:static} does
not exclude a signed estimate using the full canonical time law.

\subsection{A conditional Riccati comparison}

\begin{theorem}[Forcing upper bound and transport cost]\label{thm:transport}
Suppose $F(\sigma)\le Q$ on the entire established scaled tube
$0\le\sigma\le\Sigma=(t_1-t_0)/\ell^2$, for a real constant $Q$
that may depend on $T$.
Set $A_0=B_1L$, $K_+=\max(Q,0)$ and $h=\sqrt{K_+/2}$. Define
\begin{equation}
 \begin{array}{ll}
 f(\sigma)=\cosh(2h\sigma)+(A_0/h)\sinh(2h\sigma),&
 g_{\rm tr}(\sigma)=\sinh(2h\sigma)/(2h),\quad h>0,\\[2pt]
 f(\sigma)=1+2A_0\sigma,&g_{\rm tr}(\sigma)=\sigma,\quad h=0.
 \end{array}
 \label{eq:comparison-functions}
\end{equation}
Then, with primes here denoting scaled-time derivatives,
\begin{equation}
 m(\sigma)\ge-\frac{f'(\sigma)}{2f(\sigma)},\qquad
 m(\sigma)\ge-h-\frac1{2\sigma}\quad(\sigma>0).
 \label{eq:riccati-barrier}
\end{equation}
For $0\le x\le y\le\Sigma$, put
$W(x,y)=\exp(-4\int_x^y m(u)\dd u)$. Then
\begin{equation}
 \begin{split}
 W(x,y)&\le[f(y)/f(x)]^2,\\
 \int_x^y W(u,y)\dd u&\le\frac{f(y)}{f(x)}g_{\rm tr}(y-x),\\
 W(x,y)\delta(x)&=\delta(y)-2\int_x^y W(u,y)\dd u.
 \end{split}
 \label{eq:transport-cost}
\end{equation}
These bounds retain the finite initial layer at $\sigma=0$.
\end{theorem}

\begin{proof}
The absolute count gives $m(0)\ge-A_0$, and
$m'\ge2m^2-K_+$. The function $v=-f'/2f$ solves
$v'=2v^2-K_+$ with $v(0)=-A_0$. Since
$(m-v)'-2(m+v)(m-v)\ge0$, an integrating factor gives $m\ge v$.
For $h>0$, $f'/2f\le h\coth(2h\sigma)\le h+1/(2\sigma)$;
for $h=0$ use the displayed linear $f$. Integrating the first bound
gives the multiplier bound in \eqref{eq:transport-cost}.

Both $f$ and $g_{\rm tr}$ solve $u''=2K_+u$. Their Wronskian,
with the second solution translated by $x$, is
\[
 f(u)g_{\rm tr}'(u-x)-f'(u)g_{\rm tr}(u-x)=f(x).
\]
Consequently integrating $(g_{\rm tr}(u-x)/f(u))'$ gives
\[
 f(y)^2\int_x^y\frac{\dd u}{f(u)^2}
        =\frac{f(y)}{f(x)}g_{\rm tr}(y-x),
\]
which proves the cost bound. Finally
$\partial_u[W(u,y)\delta(u)]=2W(u,y)$ by
\eqref{eq:consumer-system}. Integrating proves the exact last identity.
\end{proof}

An independent terminal lower anchor
$\delta(y)\ge2f(y)g_{\rm tr}(y-x)/f(x)$ suffices for
$\delta(x)\ge0$. Equality with this upper allowance guarantees
nonnegativity, including $x=0$. Equality with the exact cost
$2\int_x^y W(u,y)\dd u$ gives $\delta(x)=0$, whereas strict
separation of the terminal value from that actual cost gives strict
positivity. The forward version uses an independent incoming
$\delta(x)$ and yields
\[
 |\delta(y)|\le[f(y)/f(x)]^2|\delta(x)|
                 +2f(y)g_{\rm tr}(y-x)/f(x).
\]
In physical variables $W$ integrates $M$ over physical time, and the
inhomogeneous cost for $D$ carries the additional factor $\ell^2$.
Neither the forcing estimate nor the initial moment bound supplies
either gap anchor or the tube itself.

Using only \eqref{eq:direct-forcing}, fixed collar parameters give
$K_+=C_FL$, an $O(\sqrt L)$ interior negative allowance after a
fixed positive scaled buffer, and an
$O(L)\exp(O(\sqrt L))$ upper multiplier on a fixed scaled duration.
The contour representations proved here do not improve this allowance
by the absolute estimates in Proposition~\ref{prop:group-bound}.

% 18 November 2026. Prepared for Edward Baker with OpenAI Codex assistance.
% Actual model/effort unavailable; inherited helper settings, no attestation.
\section{Joint forcing dynamics and deflated logarithmic derivatives}
\label{sec:joint}

The direct forcing bound in Section~\ref{sec:bounds} estimates each moment
separately. The hierarchy permits a more structured question: what is the
exact evolution of their forcing combination, and what does its coupling
to the gap actually control? We first keep that combination intact, then
identify an equivalent representation in derivatives of the pair-deflated
field. The identities and the two scoped tests below develop
\cite{BakerCM13}; they retain the strong-collar tube
\eqref{eq:domain} throughout. They supply no additional canonical sign
estimate.

\subsection{An exact derivative formula for the forcing combination}

Use the fixed normalized variables
\begin{equation}
 s_n=\ell^{2n}S_n=Z(n),\qquad
 v_n=\ell^{2n-1}T_n,\qquad
 m=s_1,\quad \delta=D/\ell^2,\quad
 \sigma=(t-t_0)/\ell^2.
 \label{joint:normalization}
\end{equation}
Here $s_n$ is a moment, whereas the unindexed $s$ in $z_\pm=c\pm s$
remains a local square root of $D$. In this subsection and the next,
primes denote $\sigma$ derivatives. The forcing and gap equations are
\begin{equation}
 F=12s_2+16\delta s_3,\qquad m'=2m^2-F,\qquad
 \delta'=2-4m\delta.
 \label{joint:system}
\end{equation}

\begin{proposition}[Joint forcing remainder]\label{joint:forcing-prop}
On the complete tube, including mergers,
\begin{align}
 F'-8mF={}&\mathcal R,\label{joint:forcing-law}\\
 \mathcal R={}&48v_2^2+192\delta v_2v_3+144\delta s_2^2
       +192\delta^2s_2s_3\notag\\
       &\quad-352s_3-1440\delta s_4-1152\delta^2s_5.
       \label{joint:remainder}
\end{align}
Consequently, with the multiplier $W$ of \eqref{eq:transport-cost},
\begin{equation}
 F(y)=W(x,y)^{-2}F(x)
       +\int_x^y W(u,y)^{-2}\mathcal R(u)\dd u.
 \label{joint:forcing-integral}
\end{equation}
\end{proposition}

\begin{proof}
The second and third even rows of Theorem~\ref{thm:hierarchy}, after the
constant normalization \eqref{joint:normalization}, read
\begin{align*}
 s_2'&=8ms_2-32s_3+4\delta s_2^2-40\delta s_4+4v_2^2,\\
 s_3'&=12ms_3+6s_2^2+12v_2v_3+12\delta s_2s_3
                   -60s_4-72\delta s_5.
\end{align*}
Apply the product rule to $F$ and substitute $\delta'=2-4m\delta$.
The terms involving $m$ are
$96ms_2+(-64+192)m\delta s_3=8mF$.
The coefficients of $s_3$, $\delta s_2^2$ and $\delta s_4$ are,
respectively, $-384+32=-352$, $48+96=144$ and $-480-960=-1440$.
The other terms give \eqref{joint:remainder}. Only five fixed even orders
and two fixed odd orders occur; their analytic contour expressions and
complete convergence were proved in Theorem~\ref{thm:hierarchy}.
Thus no growing-order or differentiated sharp-cutoff passage is used.
Finally the integrating factor for \eqref{joint:forcing-law} is
$\exp(-8\int_x^\sigma m)=W(x,\sigma)^2$. Integrating and dividing by
$W(x,y)^2$ gives \eqref{joint:forcing-integral}.
\end{proof}

This exact combination makes the missing estimate visible. If
$u_j=w_j/\ell$ and $\mu_j=\ell^2/Q_j$, then
$|u_j||\mu_j|^{1/2}\le\sqrt{q_0}$ and $|\mu_j|\le\gamma$.
Therefore \eqref{eq:integer-bound} gives
\begin{equation}
 |v_n|\le\sqrt{q_0}\,\gamma^{n-3/2}\mathcal A_1
          \le\sqrt{q_0}\,\gamma^{n-3/2}B_1L,
 \qquad n\ge2.
 \label{joint:odd-bound}
\end{equation}
At fixed collar parameters these separate bounds produce an $O(L^2)$
allowance for $\mathcal R$, not a smaller forcing scale. This is a
statement about that allowance, not a lower bound on the actual remainder.
Nor is there a useful sign hidden in the first two terms: their completion
is $48(v_2+2\delta v_3)^2-192\delta^2v_3^2$. A signed use of
\eqref{joint:forcing-integral} must also control its entry value and
weights; assuming the desired transport bound for those weights would
not establish it.

\begin{remark}[Exact local sign test]\label{joint:quartic-test}
For a finite polynomial test only, take $\ell=1$ and initialize
$P(0,z)=z^2(z^2+b^2)$, $b>0$. Its exact heat evolution is
$P(0,z)-tP_{zz}(0,z)+t^2P_{zzzz}(0,z)/2$. A small time neighborhood
isolates the two occurrences at the origin from the exterior roots
$\pm\ii b$. At the initial double,
\[
 m=-2/b^2,\quad s_2=2/b^4,\quad s_3=-2/b^6,\quad v_2=0,
 \qquad F=24/b^4,\quad \mathcal R=704/b^6.
\]
Initializing instead $z^2(z^2-b^2)$ changes the exterior roots to
$\pm b$ and gives $\mathcal R=-704/b^6$. Thus the finite hierarchy,
conjugation and local collar alone do not force $\mathcal R\le0$.
These polynomials lack the unchanged theta source, its positive value at
the origin and its growing complete population. They are local tests of
the algebra, not counterexamples to a canonical estimate.
\end{remark}

\subsection{The regular gap-energy identity and its weight}

A second way to retain the coupling is to eliminate $m$ from
\eqref{joint:system}. Division by $\delta$ would introduce an artificial
singularity at a merger. The following identity avoids it.

\begin{proposition}[Signed weighted gap energy]\label{joint:energy-prop}
On every closed scaled subinterval $[x,y]$ of the tube,
\begin{align}
 2\delta\delta''-(\delta')^2+4&=8F\delta^2,
       \label{joint:energy-point}\\
 8\int_x^yF\delta^2\dd u
 &=2[\delta\delta']_x^y-3\int_x^y(\delta')^2\dd u+4(y-x)
       \label{joint:energy-integral}\\
 &\le2[\delta\delta']_x^y+4(y-x).\notag
\end{align}
These formulas hold through $\delta=0$.
\end{proposition}

\begin{proof}
Differentiating the last equation of \eqref{joint:system} and substituting
the other one gives
$\delta''=8m^2\delta-8m+4F\delta$.
Multiplication by $2\delta$ and subtraction of
$(\delta')^2=(2-4m\delta)^2$ proves \eqref{joint:energy-point}.
All terms are regular, including at a zero of $\delta$.
Integrate and use
$\int_x^y\delta\delta''=[\delta\delta']_x^y-
\int_x^y(\delta')^2$ to obtain \eqref{joint:energy-integral}.
Reality of $\delta'$ proves the inequality.
\end{proof}

The boundary value is $\delta\delta'=2\delta-4m\delta^2$ and hence
still contains the interaction. The weight vanishes at mergers.
Even on an interval where it is bounded below, an upper bound for a
weighted \emph{signed} integral cannot be divided by the minimum weight
to bound the unweighted integral without controlling its negative part.
The exact transport cost in \eqref{eq:transport-cost} involves a different
weight.

\begin{proposition}[Limitation in the scalar consumer class]
\label{joint:scalar-screen}
The weighted allowance \eqref{joint:energy-integral}, together with the
two scalar equations \eqref{joint:system}, does not bound $W$ or its
inhomogeneous cost uniformly, even when $|\delta|\le d$ for fixed $d>0$.
This conclusion concerns only scalar solutions of those two equations.
\end{proposition}

\begin{proof}
For $h\ge1/(2d)$ set $m=-h$, $F=2h^2$ and $\delta=-1/(2h)$.
Both equations hold identically. The weighted integral is
$\int_x^yF\delta^2\dd u=(y-x)/2$, and
\eqref{joint:energy-integral} is an equality independent of $h$.
In contrast,
\begin{equation}
 W(x,y)=\e^{4h(y-x)},\qquad
 2\int_x^yW(u,y)\dd u=\frac{\e^{4h(y-x)}-1}{2h},
 \label{joint:scalar-cost}
\end{equation}
which are unbounded as $h\to\infty$ at fixed $y-x>0$.
\end{proof}

No divisor or full-hierarchy realization of this scalar family is asserted.
It therefore does not exclude an estimate using their additional structure.
For comparison with the baseline $C_FL$ allowance, taking
$h=\sqrt{C_FL/2}$ retains that forcing scale for large $L$ and the same
exponential transport scale. The linearization
$q=\exp(-2\int m)$, $q''=2Fq$ and $W(x,y)=[q(y)/q(x)]^2$ is equivalent
to the scalar system and adds no signed information.

\subsection{Deflation and an absolutely convergent derivative expansion}

The same forcing also has a representation using only local derivatives
of an entire function with the target pair removed. In the moving
coordinate $w$, define
\begin{equation}
 G(t,w)=\frac{H(t,c(t)+w)}{w^2-D(t)},\qquad
 U_n(t)=\sum_j w_j(t)^{-n},\quad n\ge2.
 \label{joint:deflated-field}
\end{equation}
The quotient removes both target occurrences, also at $D=0$, and is
entire in $w$. It is zero-free on $|w|<r_0$. Thus a holomorphic logarithm
exists there; its derivatives of positive order do not depend on the
chosen additive logarithm constant.

\begin{proposition}[Deflated derivative transfer]\label{joint:deflation-prop}
For $n\ge2$,
\begin{equation}
 U_n=-\frac{\partial_w^n\log G(t,0)}{(n-1)!}.
 \label{joint:U-derivative}
\end{equation}
For each fixed integer $n\ge1$,
\begin{align}
 S_n&=\sum_{k\ge0}\binom{n+k-1}{k}D^kU_{2n+2k},
       \label{joint:S-expansion}\\
 F&=4\ell^4\sum_{k\ge0}(k+1)(2k+3)D^kU_{2k+4}.
       \label{joint:F-expansion}
\end{align}
The series are absolutely convergent, uniformly on the established real
time tube. Truncation of \eqref{joint:F-expansion} after $k=K$ has error
at most
\begin{equation}
 4\ell^4\left(\sum_j|w_j|^{-4}\right)
       \sum_{k>K}(k+1)(2k+3)q^k,
 \qquad q=|D|/r_0^2\le1/16.
 \label{joint:deflation-tail}
\end{equation}
In particular, at a merger,
\begin{equation}
 F=12\ell^4U_4=-2\ell^4\partial_w^4\log G(t,0).
 \label{joint:merger-derivative}
\end{equation}
\end{proposition}

\begin{proof}
The deflated entire function has order at most one and precisely the
exterior zeros $w_j$. Its genus-one logarithmic product has a possible
affine prefactor and linear regularizers. Their derivatives vanish after
two differentiations. Differentiating the remaining factors at $w=0$
gives $-(n-1)!w_j^{-n}$, proving \eqref{joint:U-derivative}.
The sum and its derivatives converge on smaller disks by the complete
inverse-square tail and the exterior clearance.

For each occurrence, the collar gives $|D/w_j^2|\le1/16$. Apply the
binomial series to
$Q_j^{-n}=w_j^{-2n}(1-D/w_j^2)^{-n}$.
The absolute double sum is bounded by
\[
 \sum_j|w_j|^{-2n}
       \sum_{k\ge0}\binom{n+k-1}{k}16^{-k}
   =(1-1/16)^{-n}\sum_j|w_j|^{-2n}.
\]
The last sum is bounded uniformly on the fixed tube by the collar and
unit count. This proves \eqref{joint:S-expansion} with the stated
complete and time-uniform passage. Since
$F=\ell^4(12S_2+16DS_3)$, the coefficient of $D^kU_{2k+4}$ is
$12(k+1)+8k(k+1)=4(k+1)(2k+3)$, including $k=0$.
This proves \eqref{joint:F-expansion}. Taking absolute values only in
its omitted terms gives \eqref{joint:deflation-tail}. At $D=0$ retain
$k=0$ and use $3!=6$ in \eqref{joint:U-derivative}.
\end{proof}

\subsection{The fixed-source superzeta correspondence}

For clarity about the relation to classical superzeta functions, denote
zeros of $\xi$ by $\varrho$ and write the first-kind integer sum as
\[
 \mathfrak Z(n\mid\tau)=\sum_{\varrho}
              (1/2+\tau-\varrho)^{-n},\qquad n\ge2.
\]
This is the zero-based notation of \cite{Voros2014SuperzetaReview}; an
irrelevant constant normalization of the completed function does not
change it. If the independent pair tube contains $t=0$ and $D\ne0$, the
normalization \eqref{eq:source} gives
\begin{equation}
 U_n(0)=(\ii/2)^n\mathfrak Z(n\mid-\ii c(0)/2)
                     -s(0)^{-n}-(-s(0))^{-n}.
 \label{joint:superzeta-transfer}
\end{equation}
Indeed, for a zero $z_j=c+w_j$ of $H(0,\cdot)$ and its corresponding
$\varrho=(1-\ii z_j)/2$, one has
$(1-\ii c)/2-\varrho=\ii w_j/2$. Summing over the complete divisor
and removing the two target terms proves the formula. For even $n=2r$
the removed term is $2D^{-r}$. These separate pieces become singular at
a merger; the deflated expression \eqref{joint:U-derivative} is the
regular combination. Separate absolute bounds for the pieces would lose
that cancellation.

The auxiliary translation $\tau$ in a fixed-source superzeta sum is not
heat time. To see the distinction directly, put
$\Xi_t(\zeta)=8H(t,\ii(2\zeta-1))$. The chain rule and
\eqref{eq:heat} give
\begin{equation}
 \Xi_0=\xi,\qquad \partial_t\Xi_t=\tfrac14\partial_\zeta^2\Xi_t.
 \label{joint:Xi-heat}
\end{equation}
Defining a superzeta for this heat-deformed field returns the same local
derivatives in \eqref{joint:U-derivative}. The fixed-$\xi$ continuation
formulas of \cite{Voros2014SuperzetaReview} do not themselves provide a
signed estimate for these moving-center, pair-deflated derivatives on a
heat interval. Moreover, the integer forcing
\eqref{joint:F-expansion} already lies in the domain of absolute
convergence; continuation to $\Re p\le1/2$ is unnecessary for this
identity.

Together, \eqref{joint:forcing-integral}, \eqref{joint:energy-integral}
and \eqref{joint:F-expansion} distinguish three exact reorganizations
of the same quantitative problem. A useful further estimate must control
a signed remainder with its weights and entry value, a nondegenerate
quantity that bounds the transport cost, or the joint deflated derivatives
on the established tube. None of those estimates follows merely from the
representations, and direct residues have access to the same information.

% 18 November 2026. Drafted for Edward Baker with OpenAI Codex assistance.
% Actual model/effort unavailable; helper settings inherited.
\section{Backward comparison and regional gap statistics}
\label{sec:statistics}

The preceding hierarchy describes the exterior of a specified pair. A
statistical question instead asks how the smallest gap in a region evolves,
allowing its minimizing pair to change. These viewpoints can be connected
without assigning randomness to the zeros: choose a region at random,
retain the complete deterministic evolution in each region, and average
its observables. The essential additional information is the environment
of a small gap. In the all-real regime the hierarchy gives that environment
a useful sign, leading to a comparison determined by three initial gaps.
We develop the local comparison before forming its statistical consumers,
following the conditional results of \cite{BakerCM16,BakerCM17}.

\subsection{What minimum selection does and does not provide}

The strong collar already used in this paper makes the target occurrences
strict mutual nearest neighbors. Indeed, under \eqref{eq:domain},
\begin{equation}
 |\rho_j-z_\pm|\ge r_0-|s|\ge\frac{3r_0}{4}
 >\frac{r_0}{2}\ge |z_+-z_-|.
 \label{stat:mutual-neighbors}
\end{equation}
This compares each endpoint with its exterior; it does not compare two
exterior zeros. Nor does it require an attained global minimum.

\begin{proposition}[Complete versus regional minima]
\label{stat:minimum-selection}
At a simple canonical slice in $0\le t\le1/2$, the infimum of distances
between distinct zero occurrences is zero and is not attained. A finite
set of at least two occurrences has a minimum; if its points are real and
ordered, a minimizing pair is consecutive in that finite ordering.
If a conjugation-invariant target $c\pm s$ is shortest among a specified
comparison set and $s\ne0$, every compared exterior occurrence satisfies
\begin{equation}
 |w_j-s|,|w_j+s|\ge2|s|,
 \qquad |Q_j|\ge4|D|,
 \qquad |w_j|\ge\sqrt{3|D|}.
 \label{stat:relative-collar}
\end{equation}
\end{proposition}

\begin{proof}
The complete canonical count on the inherited slab is
$N_t(X)\sim X\log X/(4\pi)$ for positive real parts, with all zeros in
$|\Im z|\le1$; the closed-slab count passage is given in
\cite{BakerGT33}. If all distinct zeros were separated by $g>0$, their
open disks of radius $g/2$ would be disjoint. Disks centered at zeros with
$0<\Re z\le X$ lie in a rectangle of area
$(X+g)(2+g)$. Comparison of areas would give $N_t(X)=O_g(X)$,
contrary to the count. Simplicity excludes an attained zero distance.
Existence of a finite minimum is immediate. For real ordered points,
a nonconsecutive distance is a sum of consecutive positive gaps and
cannot be smaller than all of them. Finally, shortest-pair selection
gives the first inequalities in \eqref{stat:relative-collar}; their
product gives the bound for $Q_j$. Adding their squares gives
$2|w_j|^2+2|s|^2\ge8|s|^2$, proving the last bound.
\end{proof}

The separation in \eqref{stat:relative-collar} vanishes with the target
gap. It does not supply the independent collar in \eqref{eq:domain},
which gives the stronger fixed estimate $|Q_j|\ge15r_0^2/16$.
For a window minimum, the comparison says nothing about unexamined
exterior occurrences. At a complex slice, a closest pair also need not be
conjugation invariant: the points $\ii h,\varepsilon+\ii h,-\ii h,
\varepsilon-\ii h$, with $0<\varepsilon<2h$, have horizontal closest
pairs. Even an admissible closest pair need not detect all nonreality;
$\{-\varepsilon,\varepsilon,-\ii,\ii\}$ with
$0<\varepsilon<1/2$ has a real closest pair and a nonreal pair.
These finite examples test selection alone and do not satisfy all the
canonical hypotheses.

For perspective, at an all-real ordered heat-flow slice satisfying the
pair law, with an attained positive global minimum gap $g=2s$, the $k$th exterior zero on either side
of that pair has center-distance at least $(2k+1)s$. Thus, with fewer
terms allowed on a finite side,
\begin{equation}
 0\le M\le2\sum_{k\ge1}\frac1{4k(k+1)D}=\frac1{2D},
 \qquad D_t=2-4DM\ge0.
 \label{stat:telescoping}
\end{equation}
This useful extremal calculation belongs to its stated class. The
canonical simple divisor has no such attained global minimum, and
forward nondecrease does not furnish a backward margin.

\subsection{Adjacent real pairs require only strict isolation}

We now change the pair domain explicitly. Fix
\begin{equation}
 0\le t_*-\Delta\le t_*\le\tfrac12,
 \qquad u=t_*-t.
 \label{stat:slab}
\end{equation}
Here $u$ is \emph{physical backward time}, with no factor $\ell^2$.
Suppose the complete divisor at $t_*$ is real and simple. Follow an
adjacent pair $x_j(u)<x_{j+1}(u)$ on a common interval $0\le u<\tau$
on which the complete divisor remains real and simple and the branches
continue. This is a stopped regime, not a conclusion of initial reality
at one time. Set
\begin{equation}
 \begin{gathered}
 c_j=\frac{x_j+x_{j+1}}2,\qquad
 D_j=\frac{(x_{j+1}-x_j)^2}{4},\\
 Q_{j,\rho}=(\rho-c_j)^2-D_j
             =(\rho-x_j)(\rho-x_{j+1}),\qquad
 M_j=\sum_{\rho\ne x_j,x_{j+1}}Q_{j,\rho}^{-1},\qquad
 S_{j,n}=\sum_{\rho\ne x_j,x_{j+1}}Q_{j,\rho}^{-n}.
 \end{gathered}
 \label{stat:real-moments}
\end{equation}
All sums include every exterior occurrence. The canonical product and
inverse-square tails are still those of the complete source.
We do not impose the factor-four collar or a prescribed height interval
on the adjacent pair in this subsection.

\begin{lemma}[Strict-isolation passage]
\label{stat:strict-isolation}
On every compact subinterval of the stopped regime,
\begin{equation}
 \frac{dD_j}{du}=-2+4D_jM_j,
 \qquad
 \frac{dM_j}{du}=-2M_j^2+12S_{j,2}+16D_jS_{j,3},
 \qquad M_j,S_{j,2},S_{j,3}\ge0.
 \label{stat:backward-system}
\end{equation}
These identities do not require the quantitative collar of
\eqref{eq:domain}.
\end{lemma}

\begin{proof}
At a reference time the pair is adjacent, so every exterior real zero
has $|\rho-c_j|>\sqrt{D_j}$. Local finiteness supplies a radius
\[
 \sqrt{D_j}<R<\min_{\rho\ne x_j,x_{j+1}}|\rho-c_j|.
\]
The circle of radius $R$ about this reference center contains exactly
the two target occurrences and is zero-free. Hold it fixed while
differentiating in time. The argument-principle moments
\eqref{eq:coordinates} define $c_j,D_j$ locally, and the deflated
logarithmic derivative gives
\[
 M_j=-\frac1{2\pi\ii}\oint_\Gamma
              \frac{H_z/H}{(z-c_j)^2-D_j}\dd z.
\]
To pass the finite polynomial identity to $H$, freeze this time and
truncate the complete canonical product in opposite and conjugate groups,
retaining the target and all its occurrences. Use each truncated product
as initial data for its exact polynomial heat evolution and differentiate
at the reference time. The products and all fixed spatial derivatives
converge locally uniformly. On $\Gamma$,
$H$ and the target denominator are bounded away from zero. Hence the
logarithmic derivatives and their time derivatives converge, using
\[
 \partial_t(H_z/H)=-H_{zzz}/H+H_zH_{zz}/H^2.
\]
The contour moments and their derivatives consequently converge as well.
The finitely many nearby exterior denominators are nonzero, while the
remaining summands are $O(|\rho|^{-2n})$ and have convergent complete
tails. Thus the finite identities for $D_t$ and $M_t$, obtained in
\eqref{eq:pairlaw} and \eqref{eq:low-orders}, pass to the limit.
Their polynomial derivation uses strict separation, not the factor-four
ratio. Reversing time gives \eqref{stat:backward-system}. Finally,
$Q_{j,\rho}>0$ because the exterior is real and lies outside the adjacent
pair; this proves the signs. Finitely many local charts cover any compact
pre-collision interval.
\end{proof}

Only the identities have been extended. None of the earlier quantitative
collar estimates is asserted on this larger domain. The polynomial
passage is made at a frozen time through spatial jets; it is not a
common-time convergence theorem for finite polynomial heat evolutions.
A uniform isolating circle need not survive to a multiple collision.
The comparisons below approach such an event by continuity of the real
gap, without differentiating its square-root labeling at the event.

\subsection{A backward comparison from initial exterior mass}

\begin{theorem}[Initial-environment comparison]
\label{stat:backward-comparison}
For one pair in the stopped real regime, suppose
$0<d_0\le D(0)$ and $0\le b_0\le M(0)$. Put
$a(u)=1+2b_0u$. Then
\begin{equation}
 M(u)\ge\frac{b_0}{1+2b_0u},\qquad
 \boxed{D(u)\ge a(u)\bigl[d_0a(u)-2u\bigr]}
 \quad(0\le u<\tau).
 \label{stat:comparison}
\end{equation}
If
\begin{equation}
 d_0(1+2b_0\Delta)>2\Delta,
 \label{stat:strict-margin}
\end{equation}
this pair cannot be a colliding pair at a first finite collision time
$u_c\le\Delta$ up to which the stipulated all-real regime holds.
\end{theorem}

\begin{proof}
Positivity in \eqref{stat:backward-system} gives $M_u\ge-2M^2$.
The function $b(u)=b_0/(1+2b_0u)$ satisfies $b_u=-2b^2$, whence
\[
 (M-b)_u+2(M+b)(M-b)\ge0.
\]
Multiplying by its positive integrating factor, on each compact
pre-collision interval, proves $M\ge b$. Since $D\ge0$ there,
\[
 (a^{-2}D)_u=a^{-2}(D_u-4bD)\ge-2a^{-2}.
\]
The exact integral $\int_0^u a(v)^{-2}\dd v=u/a(u)$, including $b_0=0$,
gives \eqref{stat:comparison}. The bracket
$d_0+(2b_0d_0-2)u$ is affine, positive at zero and, under
\eqref{stat:strict-margin}, positive at $\Delta$. It is therefore
strictly positive on the full interval. If the pair were to collide
at $u_c\le\Delta$, its continuous gap would give $D(u)\to0$ as
$u\uparrow u_c$, contradicting the strictly positive lower bound.
\end{proof}

When $b_0d_0<1$, the explicit sufficient duration is
\begin{equation}
 \Delta<\frac{d_0}{2(1-b_0d_0)}.
 \label{stat:duration}
\end{equation}
When $b_0d_0\ge1$, the scalar inequality places no finite cap on
$\Delta$; the stopped-regime and coverage hypotheses still apply.
In particular, the theorem does not continue a pair's positivity
argument past an earlier unprotected collision elsewhere.

The same calculation bounds the exact backward transport cost. With
$A(u)=\int_0^u M(v)\dd v$, integration of the first equation in
\eqref{stat:backward-system} yields
\begin{equation}
 \begin{split}
 D(u)&=\e^{4A(u)}
       \left[D(0)-2\int_0^u\e^{-4A(v)}\dd v\right],\\
 2\int_0^u\e^{-4A(v)}\dd v
     &\le2\int_0^u(1+2b_0v)^{-2}\dd v
       =\frac{2u}{1+2b_0u}.
 \end{split}
 \label{stat:cost}
\end{equation}
Thus the environmental lower envelope is acquired from the initial
interaction and the real hierarchy. It is not an assumed inequality
throughout the desired time interval.

\begin{corollary}[Three initial gaps]
\label{stat:three-gaps}
Let the target gap at $u=0$ be $g=x_{j+1}-x_j$, and let its two neighboring
gaps be $h_L=x_j-x_{j-1}$ and $h_R=x_{j+2}-x_{j+1}$. Then admissible
initial data for Theorem~\ref{stat:backward-comparison} are
\begin{equation}
 d_0=\frac{g^2}{4},\qquad
 b_0=\frac1{h_L(h_L+g)}+\frac1{h_R(h_R+g)}.
 \label{stat:neighbor-mass}
\end{equation}
Consequently the computable score
\begin{equation}
 R_j(\Delta)=\frac{g_j^2}{4}
               \bigl(1+2b_{0,j}\Delta\bigr)-2\Delta
 \label{stat:score}
\end{equation}
excludes this pair from a first collision through time $\Delta$ whenever
$R_j(\Delta)>0$, subject to the stopped-regime hypotheses.
\end{corollary}

\begin{proof}
For the exterior zero $x_{j-1}$,
$Q=(x_{j-1}-x_j)(x_{j-1}-x_{j+1})=h_L(h_L+g)$.
The right exterior zero contributes $1/[h_R(h_R+g)]$ in the same way.
Every omitted exterior term is positive, so their sum is a lower bound
for $M(0)$. Apply \eqref{stat:strict-margin}.
\end{proof}

One may omit an unavailable neighbor and retain the other positive term.
If only enclosures of the gaps are known, bounds for
\eqref{stat:neighbor-mass} must respect its monotonicities; arbitrary
substitution of enclosure endpoints is not justified. For
$h_L=h_R=g$, we obtain $b_0=1/g^2$, $b_0d_0=1/4$, and
\begin{equation}
 \Delta<\frac{g^2}{6},
 \qquad\text{whereas }b_0=0\text{ gives }\Delta<\frac{g^2}{8}.
 \label{stat:equal-gaps}
\end{equation}
The factor $4/3$ is a local comparison gain. If $g/h_L$ and $g/h_R$
tend to zero, then $b_0d_0\to0$ and the relative gain tends to one.
Thus the equal-gap illustration does not settle the exceptional case of
an isolated narrow gap. Moreover, the smallest initial gap need not
minimize the environment-adjusted score: every eligible pair must be
covered.

Strictness cannot be dropped. For the exact heat polynomial
$H(t,z)=z^2-2(t-t_c)$, one has $D=2(t-t_c)$ and $M=0$.
The comparison is exact and its margin vanishes at
$\Delta=t_*-t_c$, precisely when collision occurs. This finite field
lacks the canonical population but verifies the endpoint interpretation
without numerical evidence.

\subsection{The minimum envelope and its statistical evolution}

Choose $K$ finite buffered blocks. In block $k$, fix a nonempty finite
set $J_k$ of consecutive-pair labels at $u=0$ and follow those labels
throughout the common stopped regime. Boundary-crossing pairs needed
for a covering application must be included; labels are not silently
reselected when zeros move through a window boundary. Set
\begin{equation}
 q_k(u)=\min_{j\in J_k}D_j(u),\qquad
 \mathbb P(X=k)=\frac1K.
 \label{stat:minimum-law}
\end{equation}
Only the block index $X$ is random. The $D_j$ and the canonical divisor
are deterministic.

\begin{proposition}[Weak evolution of the regional minimum]
\label{stat:weak-evolution}
On a compact subinterval of the stopped regime, every $q_k$ is locally
Lipschitz. Choose an active minimizing label $j_k(u)$, with any fixed
rule at ties. For a $C^1$ test function $\varphi$ on the relevant range,
\begin{equation}
 \frac{d}{du}\mathbb E\varphi(q_X(u))
 =\mathbb E\!\left[
    \varphi'(q_X(u))
       \{-2+4q_X(u)M_{j_X(u)}(u)\}\right]
 \quad\text{for almost every }u.
 \label{stat:weak-law}
\end{equation}
Furthermore, throughout that interval,
\begin{equation}
 q_k(u)\ge q_k(0)-2u,
 \qquad
 \mathbb P(q_X(u)\le r)
     \le\mathbb P(q_X(0)\le r+2u),\quad r\ge0.
 \label{stat:tail-transfer}
\end{equation}
\end{proposition}

\begin{proof}
A finite minimum of $C^1$ functions is locally Lipschitz. At a time
where the minimum is differentiable, its derivative equals that of
an active branch; at almost every tie, all tied active derivatives
agree. One can see the latter assertion by applying the elementary
fact that a differentiable function has derivative zero at almost
every point of its zero set to each difference $D_i-D_j$. The finite
union of exceptional null sets remains null. The chain rule and
\eqref{stat:backward-system} therefore give \eqref{stat:weak-law};
the finite average requires no additional interchange assumption.
Positivity gives $D_{j,u}\ge-2$, hence
$D_j(u)\ge D_j(0)-2u$ for every eligible branch. Taking minima proves
the first inequality in \eqref{stat:tail-transfer}. The event
$q_X(u)\le r$ is consequently contained in
$q_X(0)\le r+2u$, which proves the second.
\end{proof}

The minimum may switch branches without invalidating this envelope law.
It does not become one persistent pair solution, and its old pair's
contour, exterior moments and gap anchors do not transfer at a switch.
Equation~\eqref{stat:weak-law} also displays the statistical closure
problem: the minimum marginal depends on the joint law of the minimum
and the complete exterior interaction of its active pair. This is the
relevant BBGKY analogy. No independence or finite marginal closure has
been assumed. Corollary~\ref{stat:three-gaps} identifies a more accessible
initial statistic that bounds a sufficient margin: three consecutive
gaps, with their adjacency conditions retained.

\subsection{Pair counts, exceptional blocks and global coverage}

A one-sided minimum-tail bound does not require the entire minimum law.
At a fixed time let $C_k(r)$ count all unordered distinct-label pairs at
distance at most $r$ in a finite real-zero window. Let $G_k$ be its minimum
gap, with $G_k=+\infty$ when the window has fewer than two occurrences. Then
\begin{equation}
 \{G_k\le r\}=\{C_k(r)\ge1\},\qquad
 \mathbb P(G_X\le r)\le\mathbb E C_X(r).
 \label{stat:pair-majorant}
\end{equation}
The equality follows because a closest pair is consecutive; the inequality
is the pointwise bound $\mathbf1_{\{C_k\ge1\}}\le C_k$.
An exact minimum distribution needs absence-of-points, or void,
information, whereas \eqref{stat:pair-majorant} can be bounded by a
nonnegative two-point majorant. The three-gap score requires more:
four ordered locations must also have no intervening zeros. An
unrestricted four-point marginal cannot silently replace consecutive-gap
data.

The complete marked-pair theory in \cite{BakerOC17,BakerOCReview} gives
time moduli at a fixed positive smoothing scale, retaining multiplicity
and exterior tails. Those results provide neither an initial canonical
profile nor a sharp near-zero gap tail. To use them for
\eqref{stat:pair-majorant}, a smooth nonnegative majorant must control its
resolution, smoothing, boundary and height errors. Their constants grow
as the resolution shrinks, so a sharp-indicator limit is a separate
estimate. Distinct labels at one multiple zero contribute a collision
pair; deleting them as a diagonal would remove the event to be detected.

For the initial score \eqref{stat:score}, define the conservative count
\begin{equation}
 B_\Delta=\#\left\{k:\min_{j\in J_k}R_j(\Delta)\le0\right\}.
 \label{stat:bad-count}
\end{equation}
\begin{proposition}[Finite zero-exception threshold]
\label{stat:integer-exclusion}
If a proved estimate gives
\begin{equation}
 \mathbb P\!\left(\min_{j\in J_X}R_j(\Delta)\le0\right)<\frac1K,
 \label{stat:integer-threshold}
\end{equation}
then every covered pair has a strictly positive score and cannot be a
colliding pair at a first finite collision through backward time
$\Delta$ under the common stopped-regime hypotheses.
\end{proposition}

\begin{proof}
The probability in \eqref{stat:integer-threshold} equals $B_\Delta/K$.
Since $B_\Delta$ is a nonnegative integer, the strict bound implies
$B_\Delta=0$. Corollary~\ref{stat:three-gaps} then applies to every
eligible pair.
\end{proof}

A small probability above this threshold permits exceptions. A continuous
random-center formulation instead needs a proved positive lower measure
of centers that see any bad event, including events near window edges.
Likewise, averaging over observation times can miss an isolated collision;
the event must concern the whole time interval, or carry a proved positive
occupation contribution.

To infer an upper bound $\Lambda\le t_*-\Delta$, these finite comparisons
must be supplemented by a global continuation argument. Its inputs are
a known simple all-real start, coverage of all potentially colliding
pairs including block boundaries, and a uniform high-height exclusion or
complete tail estimate preventing the first failure from escaping to
infinity. An independently proved real/simple high-height statement
uniform over the entire slab could reduce the remaining task to a
buffered compact region. A different cutoff for each fixed time does
not automatically supply this uniform premise. Neither this uniform height
bound nor a canonical estimate strong enough for
\eqref{stat:integer-threshold} is acquired here; density one is not
zero exceptions.

Rodgers and Tao \cite{RodgersTao2020Newman} provide a precedent for using
averaged gap dynamics, but with a different conclusion. Under the
contradictory assumption $\Lambda<0$, their nearly equal spacing
consequence conflicts with a positive proportion of small gaps, proving
$\Lambda\ge0$. Their negative-time motion theorem is not used to justify
our positive-time identity; Lemma~\ref{stat:strict-isolation} supplies
that passage on \eqref{stat:slab}. Here the proposed application is a
lower \emph{upper bound}, obtained by uniform backward preservation of
reality. The established gain is the explicit local comparison and its
statistical consumer, not a new numerical bound on $\Lambda$.

% 18 November 2026. Drafted for Edward Baker with OpenAI Codex assistance.
% Actual model and effort unavailable; helper settings inherited.
\section{A joint comparison in physical backward time}
\label{sec:jointtime}

The initial-environment comparison discards the positive forcing in the
equation for $M$. This can lose a cancellation that persists throughout
the evolution. Section~\ref{sec:canonicaltime} will show that, at high
canonical zeros, even exact initial interaction data give that comparison
only an $O(g^2)$ duration. We first develop a different interface: a lower
bound for a combination of the evolving gap and its complete exterior
moments. The comparison and its finite-head acquisition are from
\cite{BakerCM23,BakerCM25}; the role of the source, rather than additional
initial data alone, will be explicit.

\subsection{The defect equation and its barrier}

Use the physical backward variable $u=t_*-t$, and initially work in the
stopped all-real simple regime of Section~\ref{sec:statistics}. For one
adjacent pair define the physical forcing and two dimensionless quantities
\begin{equation}
 \mathcal F=12S_2+16DS_3,\qquad
 V=1-2DM,\qquad
 K=2D^2\mathcal F=24D^2S_2+32D^3S_3.
 \label{jt:variables}
\end{equation}
Thus $\mathcal F$ differs by a factor $\ell^{-4}$ from the normalized
forcing $F$ of Section~\ref{sec:joint}. All sums in
\eqref{jt:variables} concern the same pair and the complete divisor.
There is no moment cutoff in the following identities.

\begin{proposition}[Joint defect equation]
\label{jt:defect}
Where $D>0$ and the complete pair laws hold,
\begin{equation}
 D_u=-2V,\qquad V_u=\frac{1-V^2-K}{D}.
 \label{jt:system}
\end{equation}
\end{proposition}

\begin{proof}
The first identity is the pair equation. Using
$M_u=-2M^2+\mathcal F$ in the derivative of $V$ gives
\[
 V_u=-2(D_uM+DM_u)=4M-4DM^2-2D\mathcal F.
\]
Since $2DM=1-V$, the first two terms equal $(1-V^2)/D$,
and the last equals $-K/D$.
\end{proof}

\begin{theorem}[Joint forcing barrier]
\label{jt:barrier}
Suppose the real complete pair laws hold on a stopped interval
$0\le u<\tau$ with $D>0$, and suppose
\begin{equation}
 K(u)\ge1-\eta^2,\qquad V(0)\le v_0,
 \qquad \eta\ge0.
 \label{jt:barrier-input}
\end{equation}
For any $v\ge\max\{v_0,\eta\}$,
\begin{equation}
 V(u)\le v,\qquad D(u)\ge D(0)-2vu.
 \label{jt:barrier-output}
\end{equation}
In particular, if $D(0)\ge d_0>0$ and
\begin{equation}
 2v\Delta<d_0,
 \label{jt:strict-margin}
\end{equation}
the selected pair cannot collide by $u=\Delta$ while the stipulated
regime and inputs remain available up to a first event. When $v=0$,
$D$ is nondecreasing on that regime.
\end{theorem}

\begin{proof}
On a compact subinterval before the stop, $D$ has a positive minimum.
For $w=V-v$, \eqref{jt:system} gives
\[
 w_u+\frac{V+v}{D}w
 \le\frac{\eta^2-v^2}{D}\le0.
\]
The coefficient is continuous. Its integrating factor is positive,
regardless of the sign of $V+v$, and $w(0)\le0$. Therefore $w\le0$.
Integration of $D_u=-2V$ proves \eqref{jt:barrier-output}.
Exhausting the stopped interval by compact subintervals preserves the
bound. At a putative collision time $u_c\le\Delta$, continuity of the
gap gives $D(u)\to0$, contradicting
$D(u)\ge d_0-2v\Delta>0$. No division by $D$ is made at the event.
\end{proof}

For a gap with $D(0)\asymp L^{-2}$, sufficient rates for this comparison
to pay a fixed available duration are
\begin{equation}
 V(0)_+=o(L^{-2}),\qquad
 1-K(u)\le\eta^2=o(L^{-4})\quad\hbox{uniformly in }u.
 \label{jt:rates}
\end{equation}
The square root of the forcing deficiency is essential. A statement
$K=1+o(1)$ alone does not give \eqref{jt:rates}. These are sufficient
rates for the comparison, not necessary conditions for actual
noncollision. In the global all-real application, a different earlier
unprotected collision ends the sign regime; the theorem does not
continue the argument past it.

\begin{remark}[Local signed complete version]
\label{jt:local-version}
The proof of Theorem~\ref{jt:barrier} uses neither $M\ge0$ nor positivity
of the individual exterior summands. It therefore applies to a
conjugation-invariant count-two tube with a possibly nonreal exterior,
provided its complete product, contour and first-time-jet passages have
been paid on the stated time domain. To see the identity transfer,
use the fixed cleared contour and the complete product as in
Lemma~\ref{lem:slice-passage}. The finite polynomial identities for $D_u$
and $M_u$ pass through uniform spatial jets on that contour and the
absolute inverse-square tails. Reality follows from conjugation.
Before the first local loss of $D>0$, the same integrating-factor proof
then applies to a directly acquired signed complete bound for $K$.
Positive omission of an exterior tail is not available in this version.
Section~\ref{sec:compact} formulates its source acquisition on a stopped
trial-coordinate domain, without assuming positivity on the desired
whole interval. At a regular target merger $\mathcal F$ is finite and
$K=2D^2\mathcal F=0$; thus an unconditional positive floor for $K$ on
the whole tube would itself already exclude that merger.
\end{remark}

\subsection{Finite neighboring heads and an explicit perturbation bound}

Consider first an equally spaced real divisor, a target pair $0,g$, and
$J$ exterior zeros on each side. Put $D=g^2/4$. The $m$th exterior
occurrence on either side has denominator $Q=m(m+1)g^2$. Denote the
interaction subsum by $b_J$ and the forcing subsum by $K_J$.

\begin{lemma}[Exact lattice head]
\label{jt:lattice}
For every integer $J\ge1$,
\begin{equation}
 2Db_J=\frac{J}{J+1},\qquad K_J=1-\frac1{(J+1)^3}.
 \label{jt:lattice-values}
\end{equation}
\end{lemma}

\begin{proof}
The two occurrences at level $m$ contribute
$2/[g^2m(m+1)]$ to $b_J$, so the first equality follows from
$1/[m(m+1)]=1/m-1/(m+1)$. Their combined contribution to $K_J$ is
\[
 \frac3{m^2(m+1)^2}+\frac1{m^3(m+1)^3}
 =\frac1{m^3}-\frac1{(m+1)^3},
\]
because $3m(m+1)+1=(m+1)^3-m^3$. Summing proves the second equality.
\end{proof}

For an all-real adjacent pair, omitted terms of $K$ and $M$ are
nonnegative. A finite lower bound for $K$ therefore needs no subtraction
of an unknown upper tail. This differs from a complete asymptotic
equality for $M$, whose remaining tail must be bounded. The following
version keeps explicit constants for a perturbed head.

\begin{proposition}[Robust finite-head bounds]
\label{jt:robust-head}
At a real-simple slice, let the target gap be $g>0$, let $d>0$, and
let $0\le e<1/3$. Suppose $g\ge(1-e)d$. For the $m$th selected
exterior occurrence on each side, $1\le m\le J$, let $r$ be its
distance to the nearer endpoint, and suppose
\begin{equation}
 r\le(1+2e)md,\qquad r+g\le(1+2e)(m+1)d.
 \label{jt:head-distances}
\end{equation}
With $b_J$ the interaction of these selected occurrences,
\begin{equation}
 \begin{split}
 K&\ge1-18e-(J+1)^{-3},\\
 1-2DM&\le1-2Db_J\le6e+(J+1)^{-1}.
 \end{split}
 \label{jt:robust-bounds}
\end{equation}
If the distance hypotheses hold with one $e_*$ throughout a stopped
all-real interval, an admissible barrier parameter is
\begin{equation}
 v=\max\left\{6e_*+(J+1)^{-1},
              \sqrt{18e_*+(J+1)^{-3}}\right\}.
 \label{jt:head-barrier}
\end{equation}
Only the initial instance of the second inequality is needed here.
\end{proposition}

\begin{proof}
Put $\gamma=(1-e)/(1+2e)\in[0,1]$. Since $Q=r(r+g)$,
\[
 \frac{D}{Q}\ge\frac{\gamma^2}{4m(m+1)}.
\]
The quadratic and cubic terms in $D/Q$ contributing to $K$ are at
least $\gamma^4$ and $\gamma^6$ times the lattice terms. Both are at
least $\gamma^6$ times those terms. Positivity of the omitted complete
exterior and Lemma~\ref{jt:lattice} give
\[
 K\ge\gamma^6[1-(J+1)^{-3}],\qquad
 2DM\ge2Db_J\ge\gamma^2\frac{J}{J+1}.
\]
Now $1-\gamma\le3e$, so
$1-\gamma^6\le18e$ and $1-\gamma^2\le6e$. This proves
\eqref{jt:robust-bounds}. Insert its uniform forcing bound and its
initial defect bound into Theorem~\ref{jt:barrier} to obtain
\eqref{jt:head-barrier}.
\end{proof}

With exact lattice head inequalities, $e_*=0$ and
$v=(J+1)^{-1}$. The strict duration becomes
\begin{equation}
 \Delta<\frac{g(0)^2(J+1)}8.
 \label{jt:head-duration}
\end{equation}
The growing allowance uses control of the head throughout time. Adding
the same neighbors only to the initial mass does not supply it. Nor is
this a finite divisor evolution: the selected labels and $D$ are
evaluated in the unchanged complete heat flow.

\subsection{A ratio-forcing floor retains a different loss}

To distinguish the joint quantity from a weaker datum, suppose $M>0$
and write $\kappa=\mathcal F/(2M^2)$. Equation~\eqref{jt:system} becomes
\begin{equation}
 V_u=2M\{(1-\kappa)+(1+\kappa)V\},
 \qquad K=(1-V)^2\kappa.
 \label{jt:ratio-system}
\end{equation}
Even near $\kappa=1$, an upper comparison for a positive defect has
coefficient approximately $4M$. The following exact integration records
the resulting allowance, rather than asserting instability of the
underlying heat flow.

\begin{proposition}[Ratio-floor comparison]
\label{jt:ratio-comparison}
In the stopped all-real regime suppose
$\mathcal F\ge2(1-\varepsilon)M^2$, where $0<\varepsilon<1$,
and put $b=M(0)>0$ and $A(u)=1+2\varepsilon bu$. Then
\begin{equation}
 \begin{split}
 M(u)&\ge b/A(u),\\
 D(u)&\ge A(u)^{2/\varepsilon}
 \left[D(0)-
 \frac{1-A(u)^{-(2-\varepsilon)/\varepsilon}}
      {b(2-\varepsilon)}\right].
 \end{split}
 \label{jt:ratio-bound}
\end{equation}
If $r=1-(2-\varepsilon)bD(0)>0$, positivity of this lower bound
holds precisely while
\begin{equation}
 u<\frac{r^{-\varepsilon/(2-\varepsilon)}-1}
          {2\varepsilon b}.
 \label{jt:ratio-horizon}
\end{equation}
For $\varepsilon=0$ the corresponding lower bound is
\begin{equation}
 D(u)\ge\e^{4bu}
       \left[D(0)-\frac{1-\e^{-4bu}}{2b}\right].
 \label{jt:ratio-zero}
\end{equation}
\end{proposition}

\begin{proof}
The hypothesis gives $M_u\ge-2\varepsilon M^2$. The same
integrating-factor comparison as in
Theorem~\ref{stat:backward-comparison} yields $M\ge b/A$.
Before collision $D\ge0$, so multiplying
$D_u\ge-2+4Db/A$ by $A^{-2/\varepsilon}$ and integrating gives
\eqref{jt:ratio-bound}; explicitly,
\[
 2\int_0^u A(v)^{-2/\varepsilon}\dd v
 =\frac{1-A(u)^{-(2-\varepsilon)/\varepsilon}}
       {b(2-\varepsilon)}.
\]
Solving the strict bracket inequality proves
\eqref{jt:ratio-horizon}. When $\varepsilon=0$, $M\ge b$ and
the integrating factor is $\e^{-4bu}$, giving
\eqref{jt:ratio-zero} directly.
\end{proof}

For a positive initial defect $\delta=1-2bD(0)$, the last bound has
horizon $\log(1/\delta)/(4b)$. At $bD(0)=1/2$,
\eqref{jt:ratio-horizon} is asymptotic to
$\log(2/\varepsilon)/(4b)$ as $\varepsilon\downarrow0$.
Thus with $b\asymp L^2$, a conservative positive defect allowance of
size $L^{-1}$ gives a duration of order $(\log L)/L^2$; an allowance
of size $\e^{-cL}$ gives order $L^{-1}$. This particular absolute-defect
comparison would require accuracy of order $\e^{-c_\Delta L^2}$
to pay a fixed duration. A favorable proved sign or additional
correlation is outside this limitation. No positive sign of the actual
canonical defect, or lower bound on its actual error, is asserted.

An exact heat control separates initial near-spacing from a whole-time
forcing floor. For $k>0$ and $0<\alpha_0<1$, consider
\begin{equation}
 H(u,z)=\e^{-k^2u}\cos(kz)+\alpha_0,
 \qquad H_u=H_{zz}.
 \label{jt:cosine-control}
\end{equation}
This is an even finite positive atomic source, not theta. Initially all
zeros are real and simple. Put $\alpha(u)=\alpha_0\e^{k^2u}$ and
$\theta(u)=\arccos\alpha(u)$. The short pair around $\pi/k$ has
\[
 D(u)=\theta(u)^2/k^2,\qquad
 V(u)=\frac{\theta(u)\alpha(u)}{\sqrt{1-\alpha(u)^2}},
 \qquad u_c=k^{-2}\log(1/\alpha_0).
\]
These formulas follow by solving $\cos(kz)=-\alpha(u)$ and
differentiating; in particular $D_u=-2V$. As $\alpha_0\downarrow0$,
the initial gaps differ from $\pi/k$ by $O(\alpha_0/k)$ and
$V(0)=(\pi/2)\alpha_0+O(\alpha_0^2)$. Choosing
$\alpha_0=\e^{-ck}$ makes those errors exponentially small, yet
$u_c=c/k$. For each fixed $k$, the other clusters stay separated from
the selected pair at $u_c$, and their periodic inverse-power tails
converge uniformly. Hence $S_2,S_3$ stay bounded while $D\to0$, giving
$K\to0$. The near-unit whole-time forcing premise fails.

The stationary cosine heat field has $V=0$, $K=1$ and is preserved
exactly by the joint barrier. The quadratic collision field has
$M=K=0$, $V=1$ and requires $\eta\ge1$, reproducing its gap-only
lifetime. These controls satisfy the local equations but lack the
unchanged canonical source and population. They test the hypotheses
and strict endpoint interpretation, not the source-specific acquisition
proved next.

% 18 November 2026. Drafted for Edward Baker with OpenAI Codex assistance.
% Actual model and effort unavailable; helper settings inherited.
\section{Canonical comparisons on a fixed positive-time slab}
\label{sec:canonicaltime}

We now acquire the preceding inputs from the unchanged canonical source.
There are three distinct conclusions. Complete initial interaction data
quantify the finite allowance of Theorem~\ref{stat:backward-comparison};
a growing head of evolving neighbors supplies the stronger joint barrier;
and spatial source jets give a sharper estimate of the actual gap drift.
All three use a fixed slab bounded away from zero. The published source
theorem already proves that its sufficiently high zeros are real and
simple throughout that slab. Consequently the conclusions refine the
dynamical comparisons on this known exterior, without extending its
coverage. We follow \cite{BakerCM21,BakerCM24,BakerCM25}, distinguishing
the published inputs recorded in \cite{BakerCMSource04} from their
moment and comparison consequences.

\subsection{The source clock and common labels}

Fix
\begin{equation}
 I=[t_0,t_*],\qquad 0<t_0\le t_*\le\tfrac12,
 \qquad \Delta=t_*-t_0.
 \label{ct:slab}
\end{equation}
The clock used below is
\begin{equation}
 G(x,t)=\frac{x}{4\pi}\log\frac{x}{4\pi}-\frac{x}{4\pi}
       +\frac{11}{8}+\frac{t}{16}\log\frac{x}{4\pi},
 \qquad G(q_n(t),t)=n.
 \label{ct:clock}
\end{equation}
Theorem~1.5(i)--(iv) of \cite{Polymath2019DBN}, with its absolute-constant
convention and the multiplicity count in its proof, supplies the following
inputs. For $0<t\le1/2$, sufficiently high zeros are exhausted by one
real simple occurrence $x_n(t)$ near each clock point, with
\begin{equation}
 x_n(t)=q_n(t)+O(q_n(t)^{-c_*t}).
 \label{ct:source-position}
\end{equation}
Here $c_*>0$ and the error constants are absolute; choose $c_*$ small
enough for both this localization and the source estimates used below.
The height and index
onsets are both exponential in $1/t$, and the theorem also supplies
the complete unit-window occurrence bound used in
\eqref{eq:count-safe}. We use the weaker displayed localization
\eqref{ct:source-position}; the source proof contains a stronger
displacement estimate. No numerical onset constant is extracted here.

\begin{lemma}[One cutoff and fixed moving labels]
\label{ct:labels}
After a single sufficiently large choice of onset depending on $t_0$,
the localization is uniform on $I$, the high labels are consecutive,
and each $x_n(t)$ is one $C^1$ branch throughout $I$. With spatial
buffers retained above that onset,
\begin{equation}
 q_n'(t)=-\frac{\pi}{4}
       \left(1+\frac{\pi t}{4q_n(t)\log(q_n(t)/(4\pi))}\right)^{-1},
 \qquad |q_n'(t)|\le\pi/4.
 \label{ct:clock-velocity}
\end{equation}
The same exterior statement holds at negative heights by evenness.
\end{lemma}

\begin{proof}
On $I$, the height threshold $\exp(C/t)$ is at most
$\exp(C/t_0)$ and $x^{-c_*t}\le x^{-c_*t_0}$. Increase the onset
to meet the index threshold as well: by \eqref{ct:clock}, an index at
height $x$ has size comparable to $x\log x$. This is an index-to-height
conversion, not an identification of the two thresholds. The same
increase pays any fixed buffer used below.

The source proof counts one occurrence in each sufficiently small disk
of radius $c/\log q_n$ and exhausts all zeros in the high region.
Clock separation is asymptotic to $4\pi/\log q_n$, so the disks are
disjoint after a further uniform enlargement. A count of one with
multiplicity proves simplicity; conjugation proves reality. Differentiating
the explicit equation $G(q_n(t),t)=n$, using
\[
 G_x(x,t)=\frac1{4\pi}\log\frac{x}{4\pi}+\frac{t}{16x},
 \qquad G_t(x,t)=\frac1{16}\log\frac{x}{4\pi},
\]
gives \eqref{ct:clock-velocity}. The moving disks are continuous and
their simple roots patch by the implicit-function theorem. Uniqueness
in each disk prevents a permutation of labels. Choosing the initial
label with the motion bound and the later window buffers included keeps
all disks in the source domain for the entire slab.
\end{proof}

This result does not assume reality of the low divisor. The later use
of positive exterior summands will explicitly impose the globally real
stopped regime. At $t_0=0$, the onset bound diverges and the error in
\eqref{ct:source-position} ceases to decay; the argument gives no
time-zero localization theorem, even though upper counts may remain
available there.

\subsection{Local geometry and the complete interaction}

Choose a buffered high label $n$, and freeze
\begin{equation}
 X=q_n(t_*),\qquad L=\log\frac{X}{4\pi}.
 \label{ct:scale}
\end{equation}
By Lemma~\ref{ct:labels}, $x_n(t)=X+O(1)$ on $I$. At a slice put
$a=x_n(t)$, $b=x_{n+1}(t)$, $g=b-a$, $D=g^2/4$, and
$d=1/G_x(a,t)$. In the next lemma take either $R=1$ or $R=L^2$, and set
\begin{equation}
 E_R=LX^{-c_*t_0}+\frac{R^2}{X}.
 \label{ct:window-error}
\end{equation}
All constants are independent of $X$ and the slice in the fixed slab.
They may depend on the stated source constants and slab. Every
asymptotic assertion is for $X$ sufficiently large; the threshold is
not evaluated numerically.

\begin{lemma}[Clock head and complete tail]
\label{ct:head-tail}
Uniformly for every actual high neighbor satisfying
$|x_{n+k}(t)-a|\le2R$,
\begin{equation}
 \frac{x_{n+k}(t)-a}{d}=k+e_k,
 \qquad |e_k|\le C E_R,
 \qquad e_0=0.
 \label{ct:normalized-position}
\end{equation}
In particular, the target and its two neighboring gaps are
$d(1+O(E_1))$, and $d\sim4\pi/L$.
At a slice where the complete divisor is real and simple, its complete
exterior interaction satisfies
\begin{equation}
 d^2M=2+O\left(E_R+\frac1{RL}\right),\qquad
 2DM=1+O\left(E_R+\frac1{RL}\right).
 \label{ct:complete-mass}
\end{equation}
Every exterior occurrence, including negative and bounded-height zeros,
is included in these estimates.
\end{lemma}

\begin{proof}
Throughout the radius-$2R$ window, $R=o(X)$ and
$G_x=O(L)$, $G_{xx}=O(1/X)$. Applying $G$ to
\eqref{ct:source-position} therefore gives
\[
 G(x_{n+k}(t),t)=n+k+O(LX^{-c_*t_0}).
\]
Subtract the corresponding equation at $a$. Taylor expansion about
$a$, with its full $O(R^2/X)$ curvature remainder, gives
\eqref{ct:normalized-position}. This uses values of the source error,
not derivatives of it. Consecutivity follows from source exhaustion;
subtracting neighboring normalized positions proves the gap assertions.

For the complete interaction, split the divisor at $|\rho-a|=R$,
assigning equality to the far set. The near indices form an interval
whose endpoints have magnitudes comparable to $RL$, with $0,1$ removed.
Indeed $G_x\asymp L$ on the window, and
\eqref{ct:normalized-position} has error $o(1)$. For such an exterior
index,
\begin{equation}
 \frac{d^2}{(\rho-a)(\rho-b)}
 =\frac1{(k+e_k)(k-1+e_k-e_1)}
 =\frac{1+O(E_R)}{k(k-1)}.
 \label{ct:head-summand}
\end{equation}
Both reference integer factors are nonzero and have magnitude at least
one, so the relative estimate is uniform over the whole head. The
positive reference sum and its two omitted tails are
\[
 \sum_{k\in\mathbb Z\setminus\{0,1\}}\frac1{k(k-1)}=2,
 \qquad
 \sum_{|k|>cRL}\frac1{|k(k-1)|}=O((RL)^{-1}).
\]
Consequently $d^2M_{\mathrm{near}}=2+O(E_R+(RL)^{-1})$.

For the far set let $r=|\rho-a|\ge R$. Eventually $g<1/4$, and
$|(\rho-a)(\rho-b)|\ge(3/4)r^2$. In each pair of unit shells
$R+j\le r<R+j+1$, $j\ge0$, the complete count bounds the number
of occurrences by $C\log(3+X+R+j)$. Hence
\begin{equation}
 |M_{\mathrm{far}}|
 \le C\sum_{j\ge0}
       \frac{\log(3+X+R+j)}{(R+j)^2}
 \le C'\frac{L+\log(R+1)+1}{R}
 =O(L/R).
 \label{ct:tail-payment}
\end{equation}
The middle inequality follows by splitting the logarithm into a
$\log(3+X)$ term and a $\log(1+R+j)$ term, and comparing the two
sums with the integrals of $r^{-2}$ and $(\log(1+r))r^{-2}$.
This accounts for every unbounded shell. If a shell crosses zero,
split its unit intervals and use reflection; the change is only in
the count constant. Boundary occurrences are retained by the stated
near/far convention.

Since $d^2=O(L^{-2})$, \eqref{ct:tail-payment} contributes
$O((RL)^{-1})$ to the normalized mass. This proves the first part
of \eqref{ct:complete-mass}; $g=d(1+O(E_R))$ and $2D=g^2/2$
give the second.
\end{proof}

The passage from local spacing to complete $M$ is essential. Local
clock data alone do not identify the complete interaction; the error
in \eqref{ct:tail-payment} pays the rest of the divisor. Conversely,
the positive lower-head argument for $K$ below will not claim an
asymptotic equality by silently discarding its tail.

\subsection{The exact initial-data allowance at high height}

At an all-real simple slice in $I$, let $h_L,h_R$ be the two neighboring
gaps and put
\[
 b_3=\frac1{h_L(h_L+g)}+\frac1{h_R(h_R+g)},\qquad
 \eta_X=LX^{-c_*t_0}+X^{-1},\qquad
 \varepsilon_X=L^{-1}+LX^{-c_*t_0}.
\]
The slice condition here is explicit: the following asymptotics are
uniform over all such slices, without presuming reality of the entire
slab.

\begin{theorem}[Initial-mass horizon]
\label{ct:initial-horizon}
For high pairs at the stated slices,
\begin{equation}
 g=\frac{4\pi}{L}(1+O(\eta_X)),\qquad
 DM=\frac12+O(\varepsilon_X),\qquad
 Db_3=\frac14+O(\eta_X).
 \label{ct:initial-asymptotics}
\end{equation}
Among all admissible initial lower data
$0<d_0\le D$ and $0\le b_0\le M$, the largest score in
\eqref{stat:strict-margin} is the exact-data score. Its strict duration
and the two-neighbor duration satisfy
\begin{equation}
 \begin{split}
 \tau_{\mathrm{full}}
 &=\frac{D}{2(1-DM)}
   =\frac{g^2}{4}(1+O(\varepsilon_X))
   =\frac{4\pi^2}{L^2}(1+O(\varepsilon_X)),\\
 \tau_3
 &=\frac{D}{2(1-Db_3)}
   =\frac{g^2}{6}(1+O(\eta_X))
   =\frac{8\pi^2}{3L^2}(1+O(\eta_X)).
 \end{split}
 \label{ct:horizons}
\end{equation}
In particular, for every fixed $\Delta>0$, even the exact-data score
is negative at all sufficiently high such pairs.
\end{theorem}

\begin{proof}
Apply Lemma~\ref{ct:head-tail} with $R=1$. The replacement
$d=(4\pi/L)(1+O(X^{-1}))$ is absorbed in $\eta_X$.
The complete tail error is $O(L^{-1})$, giving the stated error for
$DM$. The neighboring gaps are $d(1+O(\eta_X))$, so each term
of $b_3$ is $(2d^2)^{-1}(1+O(\eta_X))$; this gives
$Db_3=1/4+O(\eta_X)$.

For fixed $\Delta\ge0$, the score
$R(d_0,b_0;\Delta)=d_0(1+2b_0\Delta)-2\Delta$ is nondecreasing
in both nonnegative lower inputs. Exact $D,M$ therefore maximize it
in this information class. Eventually $DM<1$ and $Db_3<1$;
solving the strict score inequalities and using
\eqref{ct:initial-asymptotics} proves \eqref{ct:horizons}. More
explicitly, at fixed positive $\Delta$,
\begin{equation}
 R(D,M;\Delta)=D-\Delta+O(\Delta\varepsilon_X)\longrightarrow-\Delta,
 \qquad
 R(D,b_3;\Delta)\longrightarrow-3\Delta/2.
 \label{ct:score-limits}
\end{equation}
The errors are uniform over the stated slices.
\end{proof}

The full initial mass improves the asymptotic allowance by a factor
$3/2$ over two neighbors, and by a factor $2$ over the gap-only
$g^2/8$. Both improved durations still vanish as $L^{-2}$. Thus a
nonempty finite block made entirely of sufficiently high eligible pairs
is bad for the fixed-duration sufficient score with probability one
under the finite block sampling of Section~\ref{sec:statistics}.
This concerns the score event, not an actual collision event.

The limitation is confined to the one-step Riccati lower envelope using
initial $D,M$. It excludes neither new time-dependent information nor
reinitialization, divided steps or direct source control. For an exact
check, the field
$H(t,z)=\e^{\pi^2t/g^2}\cos(\pi z/g)$ has stationary zeros,
$D=g^2/4$, $M=2/g^2$ and $D_u=0$, while that same comparison stops
at $g^2/4$. The discarded forcing balances $2M^2$ exactly. The field
has linear zero count and is not the canonical theta source; its purpose
is to distinguish an allowance from a dynamical lifetime.

\subsection{A growing head pays the joint comparison}

Return to the fixed common labels on $I$, let $u=t_*-t$, and suppose
the complete divisor at $t_*$ is real and simple. For this subsection
apply positive omission only on its globally real simple stopped part.
Set
\begin{equation}
 R=L^2,\qquad
 J=\left\lfloor\frac{RL}{64\pi}\right\rfloor,\qquad
 E_X=LX^{-c_*t_0}+\frac{L^4}{X}.
 \label{ct:growing-head}
\end{equation}

\begin{theorem}[Canonical acquisition of the joint barrier]
\label{ct:joint-acquisition}
There are constants $C_K,C_M$ such that, uniformly on that stopped
regime,
\begin{equation}
 K\ge1-C_K(E_X+L^{-9}),\qquad
 |V|\le C_M(E_X+L^{-3}).
 \label{ct:acquired-inputs}
\end{equation}
Consequently, with
\begin{equation}
 \eta=\sqrt{C_K(E_X+L^{-9})},\qquad
 v_X=\max\{C_M(E_X+L^{-3}),\eta\},
 \label{ct:barrier-parameter}
\end{equation}
the joint comparison gives
\begin{equation}
 D(u)\ge D(0)-2v_Xu,\qquad
 v_X=O(L^{-3}+\sqrt{E_X})=O(L^{-3}).
 \label{ct:joint-output}
\end{equation}
Every fixed duration within the stipulated slab satisfies the strict
margin eventually, since
\begin{equation}
 \frac{2v_X\Delta}{D(0)}=O(\Delta/L)\longrightarrow0.
 \label{ct:relative-margin}
\end{equation}
\end{theorem}

\begin{proof}
The fixed head contains labels $n-J,\ldots,n+J+1$. All lie within
distance $R$ of $a=x_n(t)$ for every $t\in I$, after increasing $X$.
Indeed $G_x\ge L/(8\pi)$ on the radius-$2R$ window, so a distance
$R$ allows at least $RL/(8\pi)$ clock indices. This is much larger
than $J+1$, and the source displacements are $o(1)$. The slack pays
both the endpoint asymmetry and the moving spatial buffers. Also
$J=O(L^3)=o(n)$, so the leftmost label still meets the source index
onset. The same finite labels, rather than a moving cutoff, are
therefore available throughout the slab.

Lemma~\ref{ct:head-tail} supplies
$(x_{n+k}-a)/d=k+e_k$ with $|e_k|\le CE_X$ for all these labels.
It follows that $g\ge(1-e)d$ and the two endpoint-distance
inequalities \eqref{jt:head-distances} hold with $e=C'E_X$:
subtracting two normalized positions costs at most $2CE_X$, and
each reference integer distance is at least one. Eventually $e<1/3$.
Proposition~\ref{jt:robust-head}, together with $J\asymp L^3$,
gives the first inequality in \eqref{ct:acquired-inputs}. Its omitted
forcing terms are positive on the stopped regime. No upper tail
subtraction or independent finite-divisor evolution has been used.

For $M$, use the complete statement \eqref{ct:complete-mass} at
$R=L^2$. It gives $2DM=1+O(E_X+L^{-3})$, including the unbounded
tail payment \eqref{ct:tail-payment}, hence the second inequality
in \eqref{ct:acquired-inputs}. Only its initial instance is needed
for Theorem~\ref{jt:barrier}. That theorem gives
\eqref{ct:joint-output}. At fixed $t_0>0$, $E_X$ is smaller than
every fixed negative power of $L$, while
$D(0)\sim4\pi^2/L^2$. These facts prove the last equality in
\eqref{ct:joint-output} and \eqref{ct:relative-margin}.
\end{proof}

When $v_X>0$, the algebraic duration budget $D(0)/(2v_X)$ is at
least a constant times $L$. This is a budget of the comparison; its
inputs have only been established on the fixed slab \eqref{ct:slab}.
It is not a claim of backward evolution for a time of order $L$.
The change from the shrinking allowance in
Theorem~\ref{ct:initial-horizon} is the acquired whole-time forcing
floor, not a more accurate initial mass.

The same-information comparison is already visible:
\eqref{ct:acquired-inputs} also bounds $|V|$ at every time, so
$|D_u|\le2C_M(E_X+L^{-3})$ follows directly, without taking the
square root of a forcing error. More fundamentally, the clock theorem
already proves exterior reality and simplicity independently of the
low divisor. The barrier has a weaker data interface---a full-time
lower $K$ bound plus initial $V$---but no cheaper acquisition or
improved canonical onset is demonstrated here.

\subsection{An analytic source residual and the velocity consequence}

There is a sharper way to retain the common motion of neighboring
zeros. Polymath's Remark~9.4 explicitly describes the spatial-jet
method for a velocity asymptotic and leaves its details to the reader
\cite{Polymath2019DBN}. The proof below supplies the local residual,
Cauchy and denominator estimates needed for this application; the
velocity method is not claimed as new.

Let $q=q_n(t)$, $L_q=\log(q/(4\pi))$, and define
\begin{equation}
 \lambda_q=L_q/4,\qquad \kappa_0=\pi/8,\qquad
 \epsilon_q=q^{-c_*t_0}+L_q/q.
 \label{ct:jet-scale}
\end{equation}
Write $\mathfrak M_t$ for the nonvanishing carrier denoted $M_t$ in
\cite{Polymath2019DBN}; it is not the exterior interaction $M$.
Set $A_q=\mathfrak M_t((1+\ii q)/2)\ne0$.

\begin{lemma}[Holomorphic local residual]
\label{ct:analytic-residual}
For a sufficiently small fixed $r>0$, uniformly for $t\in I$ and
all sufficiently high clock labels,
\begin{equation}
 H(t,q+w)=A_q\e^{-\kappa_0w}
          \{2\ii\sin(\lambda_qw)+E_{n,t}(w)\},
 \quad |w|\le r/L_q,
 \quad \sup|E_{n,t}|\le C\epsilon_q,
 \label{ct:residual}
\end{equation}
where $E_{n,t}$ is an actual holomorphic function on the disk.
\end{lemma}

\begin{proof}
Keep both analytic carrier terms in \cite[\S9, (102)]{Polymath2019DBN}.
For $w=(x-q)+\ii y$, their exponents have leading parts
\[
 (-\ii w)(\lambda_q-\ii\kappa_0)
       =-\ii\lambda_qw-\kappa_0w,
 \qquad
 (\ii w)(\lambda_q+\ii\kappa_0)
       =\ii\lambda_qw-\kappa_0w.
\]
The corresponding multipliers are $\overline{A_q}$ and $A_q$.
The relative errors in those two terms are $O(q^{-c_*t_0})$ on the
disk, with the fixed small source exponent chosen above.
Lemma~9.2 and (103) of that source give
$\overline{A_q}/A_q=-1+O(L_q/q)$ at the clock point.
The two exponentials have uniformly bounded modulus on
$|w|\le r/L_q$. Divide the source approximation by
$A_q\e^{-\kappa_0w}$: the leading terms combine to
$\e^{\ii\lambda_qw}-\e^{-\ii\lambda_qw}$ and the remaining
terms are bounded by $C\epsilon_q$.

Define $E_{n,t}$ by the exact left side of \eqref{ct:residual}
minus $2\ii\sin(\lambda_qw)$. The quotient uses a nonvanishing
entire factor, so this definition proves holomorphy, and the source
estimate proves its uniform bound. No derivative of a symbol $O(\cdot)$
is taken. The leading frequency is deliberately $L_q/4$; the smaller
time-dependent phase refinement is included in the stated error.
The weaker display (104) in the source, with an allowance involving
$|\Im w|\log q$, would not give a vanishing residual on this disk
and is not used for the Cauchy step.
\end{proof}

\begin{proposition}[Source jets and the complete gap defect]
\label{ct:velocity}
For the common high labels on $I$, put
$\epsilon_X=X^{-c_*t_0}+L/X$. Without an assumption that the low
divisor is real,
\begin{equation}
 \begin{split}
 |x_n'(t)+\pi/4|&\le C L\epsilon_X,\\
 |g_t|&\le C L\epsilon_X,\qquad
 |D_t|\le C\epsilon_X,\qquad
 |1-2DM|\le C\epsilon_X.
 \end{split}
 \label{ct:velocity-bounds}
\end{equation}
Here $M$ is the complete exterior sum for the adjacent high pair.
In particular, for $s,t\in I$,
\begin{equation}
 |g(t)-g(s)|\le C L\epsilon_X|t-s|,
 \qquad |D(t)-D(s)|\le C\epsilon_X|t-s|.
 \label{ct:integrated-jets}
\end{equation}
\end{proposition}

\begin{proof}
Source localization puts the actual real root
$w_n=x_n(t)-q_n(t)$ in $|w_n|\le r/(4L_q)$ for sufficiently large
$q$, uniformly in $t$. Equation~\eqref{ct:residual} at that root
implies $\sin(\lambda_qw_n)=O(\epsilon_q)$. On this fixed small
real interval the sine is bounded below in modulus by a constant
times its argument, so
\begin{equation}
 |w_n|\le C\epsilon_q/L_q.
 \label{ct:sharpened-position}
\end{equation}
Cauchy's formula on the larger disk gives, in its inner half,
\[
 |E'_{n,t}|\le CL_q\epsilon_q,\qquad
 |E''_{n,t}|\le CL_q^2\epsilon_q.
\]
Put $f(w)=2\ii\sin(\lambda_qw)+E_{n,t}(w)$. At the root,
$\lambda_qw_n=O(\epsilon_q)$, and therefore
\begin{equation}
 f'(w_n)=2\ii\lambda_q(1+O(\epsilon_q)),\qquad
 f''(w_n)=O(L_q^2\epsilon_q),\qquad
 |f'(w_n)|\ge c_1L_q.
 \label{ct:root-jets}
\end{equation}
The second bound uses the small sine at the root and the paid
second derivative of the residual. Spatially differentiating the
exact representation \eqref{ct:residual}, with $t,q,A_q$ fixed,
and then setting $f(w_n)=0$ gives
\[
 \frac{H_{zz}(t,x_n)}{H_z(t,x_n)}
 =\frac{f''(w_n)}{f'(w_n)}-2\kappa_0
 =-\frac\pi4+O(L_q\epsilon_q).
\]
Implicit differentiation of $H(t,x_n(t))=0$ and $H_t=-H_{zz}$
gives the exact physical velocity $x_n'=H_{zz}/H_z$.
This proves its asserted bound. The leading velocity $-\pi/4$
cancels for neighboring labels, so $g_t=O(L\epsilon_X)$;
since $g=O(L^{-1})$, $D_t=gg_t/2=O(\epsilon_X)$.

For completeness, this derivative estimate transfers to the complete
$M$, not only to a local head. The high pair is strictly isolated;
apply the finite-product pair identity at a reference time and pass
through the cleared pair contour as in Lemma~\ref{lem:slice-passage}.
Equivalently, subtract the two finite root-velocity sums before their
complete opposite/conjugate-group limit. The exterior difference is
\[
 g_t=\frac4g-2g\sum_{\rho\ne a,b}
                       \frac1{(\rho-a)(\rho-b)}.
\]
Every surviving tail term is inverse-square and hence absolutely
summable by the complete strip and unit count. The nearby denominators
are nonzero by strict isolation, including when lower zeros are
nonreal. Thus $D_t=2-4DM=2V$ holds with the complete real sum,
proving the last bound in \eqref{ct:velocity-bounds}. No separate
nonabsolute first-power sum is estimated. Integration on the common
label branches proves \eqref{ct:integrated-jets}.
\end{proof}

This proof uses complex spatial neighborhoods and the exact heat
equation, not a complex-time source estimate or a time derivative of
the localization error. Since $\epsilon_X L^2\to0$, its squared
half-gap allowance pays every fixed duration in $I$ at sufficiently
high labels. The corresponding allowance for the actual gap has the
additional factor $L$ displayed in \eqref{ct:integrated-jets}.

\subsection{The matched direct clock comparison}

The same source information also encloses the gaps at each time.
Let $h_n(t)=q_{n+1}(t)-q_n(t)$. From
\eqref{ct:sharpened-position},
\begin{equation}
 |g_n(t)-h_n(t)|\le C\epsilon_X/L.
 \label{ct:direct-enclosure}
\end{equation}
This is a common-time enclosure, not merely an initial-value estimate.
Write $v(q,t)$ for the explicit clock velocity in
\eqref{ct:clock-velocity}, and $L_q=\log(q/(4\pi))$. Differentiating
this explicit function gives
\[
 v_q(q,t)=-\frac{\pi^2t(L_q+1)}{16q^2L_q^2}
          \left(1+\frac{\pi t}{4qL_q}\right)^{-2}.
\]
Since $h_n=O(L^{-1})$, the mean-value theorem implies
$|h_n'(t)|\le C/(X^2L^2)$. Combining two endpoint enclosures
\eqref{ct:direct-enclosure} with this bound yields
\begin{equation}
 \begin{split}
 |g_n(t)-g_n(s)|
 &\le C\epsilon_X/L+C|t-s|/(X^2L^2),\\
 |D_n(t)-D_n(s)|
 &\le C\epsilon_X/L^2+C|t-s|/(X^2L^3).
 \end{split}
 \label{ct:matched-direct}
\end{equation}
For the second line, use $D=g^2/4$ and $g(t)+g(s)=O(L^{-1})$.
Both endpoint source errors are included in the constants. For a fixed
positive time separation this direct bound pays roughly $L^2$ less
source error than the integrated squared half-gap bound. For very short
times, \eqref{ct:integrated-jets} can be better because its allowance
vanishes with $|t-s|$. The minimum of the two bounds is legitimate;
neither dominates in every time regime.

These comparisons locate the gain and its boundary. The original
initial-$M$ envelope loses a cancellation; the acquired $K$ floor
retains enough of it to pay a fixed slab, and the analytic source
residual estimates it still more sharply. All source acquisitions
remain in the exterior already covered by the published theorem.
They give no numerical cutoff, no estimate uniform down to $t=0$,
no new compact gap law, and no improved bound on $\Lambda$. Extending
the useful domain requires new source information of the type considered
in Section~\ref{sec:compact}, rather than repeated evaluation of this
already protected high exterior.

% 18 November 2026; prepared for Edward Baker with OpenAI Codex assistance.
% Actual model/effort unavailable. Source: CM Note27 and E25 assessment.
\section{Compact source acquisition and a local stopping theorem}
\label{sec:compact}

The positive-time source of the preceding section supplies both a joint
estimate and pre-existing exterior reality. A compact application has a
different logical requirement: it must acquire the needed information
without assuming that the target gap remains positive on the interval
being tested. A separating contour gives such a source interface because
its count may remain two through a target merger. We first derive that
interface for trial pair coordinates, then state the stopping theorem and
its full error requirements \cite{BakerCM27}.

\subsection{Complete moments at trial coordinates}

Fix a compact real time interval and a positively oriented circle
$\Gamma=\{|z-a|=r\}$, with real $a$. Suppose $H$ is nonzero on a
neighborhood of $\Gamma$ throughout the interval and counts exactly two
occurrences inside. These hypotheses allow a double target or a conjugate
nonreal pair. No assertion about exterior reality is made. The canonical
order-one product and its locally uniform logarithmic derivative are
used on this contour. Let $B=H_z/H$ and
\[
 A_j=\frac1{2\pi\ii}\oint_\Gamma(z-a)^jB(t,z)\dd z,
 \qquad c=a+A_1/2,\qquad D=A_2/2-(A_1/2)^2.
\]
These are the actual coordinates. For separate trial real coordinates
$(\widehat c,\widehat D)$ put
$Q_{\!\widehat c,\widehat D}(z)=(z-\widehat c)^2-\widehat D$.
Both roots of this trial quadratic are required to lie inside $\Gamma$;
they need not coincide with the actual targets.

\begin{proposition}[Trial-coordinate exterior traces]
\label{cmp:trial}
For every integer $p\ge1$,
\begin{equation}
 S_p^\Gamma(t;\widehat c,\widehat D)
 =-\frac1{2\pi\ii}\oint_\Gamma
       \frac{B(t,z)}{Q_{\!\widehat c,\widehat D}(z)^p}\dd z
 =\sum_{\rho\text{ outside }\Gamma}
       \bigl((\rho-\widehat c)^2-\widehat D\bigr)^{-p}.
 \label{cmp:trace}
\end{equation}
The sum includes the complete exterior with multiplicity. It converges
absolutely and locally uniformly on compact trial sets with strict
contour separation. At the actual $(c,D)$ it equals $S_p$.
\end{proposition}
\begin{proof}
Set $f=Q_{\!\widehat c,\widehat D}^{-p}$. All its poles are inside
$\Gamma$, and $f=O(z^{-2p})$ at infinity. For an actual interior root
$\rho$, the rational function $f(z)/(z-\rho)$ has all poles inside and
decays as $O(z^{-2p-1})$, so its contour integral is zero. Coincident or
multiple poles cause no difficulty. For an exterior root the residue
outside is $f(\rho)$, and the same decay gives the interior integral
$-f(\rho)$. The constant term and genus-one constant regularizers of the
logarithmic derivative integrate to zero since $\oint f=0$.
The complete product converges normally on the cleared contour. Its
exterior contributions to \eqref{cmp:trace} are dominated by a multiple
of $\sum_\rho(1+|\rho|)^{-2}$, finite for the order-one canonical divisor;
compact trial sets permit the same dominating tail. Thus termwise
contour passage gives the result. This argument extends the rational
trace of Proposition~\ref{prop:trace} without identifying trial and actual
pair coordinates.
\end{proof}

Define the complete signed trial forcing
\begin{equation}
 \begin{split}
 \mathcal F_\Gamma&=12S_2^\Gamma+16\widehat D S_3^\Gamma,\\
 K_\Gamma&=2\widehat D^2\mathcal F_\Gamma
 =-\frac1{2\pi\ii}\oint_\Gamma
 \left(\frac{24\widehat D^2}{Q^2}
       +\frac{32\widehat D^3}{Q^3}\right)B(t,z)\dd z .
 \end{split}
 \label{cmp:forcing}
\end{equation}
Here and below $Q$ in a trial integral denotes the trial quadratic.
Conjugation makes these quantities real for real trial coordinates.
Unlike a positive-head estimate, \eqref{cmp:forcing} does not discard
any possibly nonreal exterior terms.

At the actual coordinates the local pair law and the first hierarchy
row give, in backward time $u=t_*-t$,
\begin{equation}
 D_u=-2+4DM,\qquad M_u=-2M^2+\mathcal F,
 \qquad \mathcal F=12S_2+16DS_3.
 \label{cmp:local-system}
\end{equation}
The cancellation proof of Theorem~\ref{thm:hierarchy} uses finite rational
identities before complete contour passage; neither step uses positivity
of exterior summands. It therefore applies to the present strictly
separating local tube, with its own compact contour bounds in place of
the strong-collar constants. These hypotheses must hold on the interval
at issue. An earlier positive-time certificate does not supply them on
an enlarged negative-time interval.

\subsection{Stopping without an assumed positive trajectory}

Suppose unconditional contour information gives
$|c(t)-a|\le C$ and $D(t)\le D_{\max}$ throughout the desired interval.
Its lower enclosure for $D$ is allowed to cross zero. Choose
$0<d_{\mathrm{stop}}<d_0$ and assume an independent initial enclosure
$D(0)\ge d_0$, $V(0)\le v_0$, where $V=1-2DM$ uses backward time.

\begin{theorem}[A conditional source bound prevents a lower exit]
\label{cmp:stop}
Assume the count-two contour and its clearance hold for
$0\le u\le\Delta$. Suppose all trial quadratics with
$|\widehat c-a|\le C$ and
$d_{\mathrm{stop}}\le\widehat D\le D_{\max}$ lie inside the contour and
\begin{equation}
 K_\Gamma(t_* -u;\widehat c,\widehat D)\ge\kappa
 \label{cmp:kfloor}
\end{equation}
uniformly on that trial set. Put
\[
 v=\max\{0,v_0,\sqrt{\max(0,1-\kappa)}\}.
\]
If
\begin{equation}
 d_0-2v\Delta-d_{\mathrm{stop}}>0,
 \label{cmp:exit-margin}
\end{equation}
then the actual pair satisfies $D(u)>d_{\mathrm{stop}}$ throughout the
interval. Exterior zeros need not be real.
\end{theorem}
\begin{proof}
Until a first lower exit $D=d_{\mathrm{stop}}$, the actual coordinates
belong to the trial set. Thus $K\ge\kappa$ there, and
\eqref{cmp:local-system} gives
$D_u=-2V$ and $V_u=(1-V^2-K)/D$. For $w=V-v$,
\[
 w_u+\frac{V+v}{D}w\le\frac{1-v^2-\kappa}{D}\le0.
\]
On every compact pre-exit interval the integrating factor is positive,
so $V\le v$ and $D\ge d_0-2vu$. At a hypothetical first exit,
\eqref{cmp:exit-margin} gives $D>d_{\mathrm{stop}}$, a contradiction.
The other coordinate and contour exits are already excluded by the
unconditional assumptions. No division by $D$ at a merger is used.
\end{proof}

The distinction between the trial statement and an unconditional positive
forcing floor is substantive. On a regular counted tube $\mathcal F$
is finite at a target merger, whereas $K=2D^2\mathcal F$ vanishes there.
Hence an unconditional $K\ge k>0$ already excludes a merger by continuity
from a positive start. Such a floor can be a valid source certificate,
but is not a weaker substitute for the sought positivity. Theorem
\ref{cmp:stop} instead proves the necessary interval of positivity from
an estimate on a conditional trial region.

\subsection{A complete analytic error contract}

Let $G$ be any fixed nonzero scalar multiple of $H$. Suppose source
approximants $h,h_1$ have uniform errors $e_0,e_1$ for $G,G_z$ on the
whole time-contour cell, and $|h|\ge b_0>e_0$, $|h_1|\le b_1$. Then
\begin{equation}
 \left|B-\frac{h_1}{h}\right|
 \le e_B:=\frac{e_1}{b_0-e_0}
          +\frac{b_1e_0}{b_0(b_0-e_0)}.
 \label{cmp:quotient-error}
\end{equation}
Indeed, subtract the quotients and bound the numerator error over $G$
and the reciprocal-denominator error separately. If
$|\widehat c-a|\le C$, $|\widehat D|\le D_{\max}$,
$0\le C<r$, $D_{\max}\ge0$ and
$q=(r-C)^2-D_{\max}>0$, integration around the circle gives the source
parts of the moment and forcing errors:
\begin{align}
 E_M&\le re_B/q,&
 E_{A_j}&\le r^{j+1}e_B,\notag\\
 E_K&\le re_B\left(24D_{\max}^2/q^2+32D_{\max}^3/q^3\right).
 \label{cmp:kernel-errors}
\end{align}
All parameter dependence belongs to the trial enclosure; discarding a
proved time-state correlation by replacing it with independent boxes is
an additional approximation, not part of these inequalities.

An annular logarithm gives an alternative interface when
$\mathcal L=\Log(G/(z-a)^2)$ is single valued on a cleared annulus.
Since $B=\mathcal L_z+2/(z-a)$ and the rational second term has zero
integral against $Q^{-p}$, integration by parts gives
\begin{align}
 S_p^\Gamma&=-\frac{p}{2\pi\ii}\oint_\Gamma
                  \mathcal L\frac{Q'}{Q^{p+1}}\dd z,\notag\\
 K_\Gamma&=-\frac1{2\pi\ii}\oint_\Gamma\mathcal L Q'
              \left(\frac{48\widehat D^2}{Q^3}
                    +\frac{96\widehat D^3}{Q^4}\right)\dd z.
 \label{cmp:log-kernel}
\end{align}
A constant may be subtracted from $\mathcal L$ without changing either
integral. A uniform log error $e_{\mathcal L}$ gives
\begin{align}
 E_M&\le 2r(r+C)e_{\mathcal L}/q^2,\notag\\
 E_K&\le 2r(r+C)e_{\mathcal L}
        \left(48D_{\max}^2/q^3+96D_{\max}^3/q^4\right).
 \label{cmp:log-errors}
\end{align}
For matching log branches and $|G-h|\le e_0<b_0\le|h|$, the log increment
has modulus at most $-\log(1-e_0/b_0)$ by its convergent power series.
Branch matching and annular clearance remain explicit obligations.

Finite contour quadrature has its own error. If the parameterized
integrand is holomorphic for $|\Im\theta|\le s$ and bounded there by
$A_s$, its normalized $N$-node trapezoidal error is at most
\begin{equation}
 \frac{2A_s}{\e^{Ns}-1}.
 \label{cmp:trapezoid}
\end{equation}
To prove this, shift the Fourier coefficient contour to the appropriate
strip edge, obtaining $|a_k|\le A_s\e^{-|k|s}$. The trapezoidal mean
aliases only the coefficients $a_{jN}$; summing their geometric bound
gives \eqref{cmp:trapezoid}. The source and every trial denominator must
be cleared throughout that complex strip, not only at real nodes.

The budgets $e_0,e_1$ must include complete source evaluation, source
integration tails, lattice or quadrature errors, coefficient reconstruction,
arithmetic, spatial approximation and whole-cell time variation. Time
remainders may use $\partial_t^jG=(-1)^j\partial_z^{2j}G$ with actual
Cauchy bounds. Time samples alone do not supply uniform estimates, and
an asymptotic error cannot be differentiated without an analytic residual.
The complete contour representation requires no separately omitted zero
tail after the full source error has been paid.

A negative-time source formula is available analytically. For $s\ge0$,
\begin{equation}
 H(-s,z)=\pi^{-1/2}\int_\R\e^{-y^2}
                         H(0,z+2\sqrt{s}\,y)\dd y.
 \label{cmp:negative-gaussian}
\end{equation}
Insert the complete Fourier representation of $H(0,\cdot)$; its positive
source and fixed $\Im z$ dominate the double integral. The Gaussian
Fourier integral supplies the multiplier $\e^{-s u^2}$, proving the
formula by Fubini. The real spatial shift in \eqref{cmp:negative-gaussian}
differs from the imaginary shift at positive heat time. Consequently an
existing positive-time xi quadrature estimate cannot be extrapolated to
negative time without new source and tail bounds.

\subsection{What a constant barrier can and cannot retain}

The stopping theorem preserves coupling in its source kernel, but its
single constant $v$ still discards the temporal profile. The following
screen quantifies that limitation independently of any numerical example.

\begin{proposition}[A small included forcing value limits the constant barrier]
\label{cmp:small-slice}
Suppose the initial gap square satisfies $D(0)\le d_+$ and an included
slice covered by the proposed forcing floor has $K\le k_*<1$. Every
certificate using the constant parameter of Theorem~\ref{cmp:stop} then
has $v\ge\sqrt{1-k_*}$ and can satisfy its strict margin only if
\begin{equation}
 \Delta<\frac{d_+}{2\sqrt{1-k_*}}.
 \label{cmp:duration-cap}
\end{equation}
In particular, at a slice with $|D|\le\delta$ and
$|\mathcal F|\le A$, one may take $k_*=2\delta^2 A<1$.
\end{proposition}
\begin{proof}
A valid floor must have $\kappa\le K\le k_*$. The definition of $v$
therefore gives the asserted lower bound. Dropping the positive stop
threshold and enlarging the initial lower anchor to the actual initial
value, then to $d_+$, can only enlarge the allowable duration. This gives
\eqref{cmp:duration-cap}. The last statement is $|K|=2D^2|\mathcal F|$.
\end{proof}

This is a limitation of the displayed constant-$v$ consumer, not of all
comparisons using a constant $K$ bound, and not of actual dynamics.
Retaining the initial transient, using a time-varying barrier, renewing
anchors or integrating the regular $(D,M)$ system changes the consumer.
The supporting example in Appendix~\ref{app:compact-example} applies the
proposition to inherited certified data; it is not needed for the theorem.

An independent rectangular forcing floor can lose even the equilibrium
advantage. At a start with $V_0=0$, suppose the chosen reduction uses
$\mathcal F\ge f$ with $0\le f\le1/(2d_0^2)$ and
$d_{\mathrm{stop}}=\alpha d_0$, $0<\alpha<1$. It replaces the trial
forcing by $\kappa=2d_{\mathrm{stop}}^2f\le\alpha^2$, so
$v\ge\sqrt{1-\alpha^2}$. The exit margin requires
\begin{equation}
 \Delta<\frac{d_0}{2}\sqrt{\frac{1-\alpha}{1+\alpha}}<\frac{d_0}{2}.
 \label{cmp:rectangular-screen}
\end{equation}
Thus that independent reduction can be worse than the gap-only deadline.
It does not exclude correlated trial regions or a source estimate for
the combined kernel itself.

Every proposed source acquisition must also supply the direct comparator
with the same information: the centered discriminant, complete $M$
transport and regular $(D,M)$ evolution. If those shared coordinates
already prove positivity, computing $K$ adds a magnitude or cost
comparison, not new zero coverage. Likewise the lattice block identity
of Section~\ref{sec:lattice} does not by itself prevent concentration
into a single small gap. These are distinct limitations of scalar and
aggregate consumers, not reasons to discard the shared-source structure.

% 18 November 2026. Prepared for Edward Baker with OpenAI Codex LLM assistance.
% Actual model and effort unavailable; helper settings inherited.
% E31 author A. Sources: CM Notes32/34, with E29/E30 scope retained.
\section{Forward finite-block statistics from a counted initial slice}
\label{sec:forward-statistics}

Known information at time zero can be propagated forward before a
backward comparison is attempted. The two operations require different
estimates. Forward root repulsion can improve a lower bound for a block
minimum, whereas backward protection needs a lower interaction or forcing
bound. We first establish the forward statements, retaining the complete
exterior in each one. The results of \cite{BakerCM32,BakerCM34}, assessed
in \cite{BakerCMReview29,BakerCMReview30}, are conditional finite estimates;
they do not supply a new time-wide canonical source certificate.

\subsection{What a complete initial census can imply}

\begin{lemma}[Reality from a complete singleton census]
\label{fwd:singleton}
Let $H(0,\overline z)=\overline{H(0,z)}$, and let $\Omega$ be a bounded
conjugation-invariant region with zero-free boundary. Suppose a complete
census of its zeros consists of finitely many cells whose horizontal
projections are pairwise disjoint. If every occupied cell has count one,
including multiplicity, every zero in $\Omega$ is real and simple.
\end{lemma}
\begin{proof}
If $\rho$ is a zero in one cell, its conjugate is also in $\Omega$.
Completeness places $\overline\rho$ in an occupied cell. Since the real
coordinates agree, disjoint horizontal projections place both occurrences
in the same cell. Its count one forces $\rho=\overline\rho$, and also
forces multiplicity one. No replacement of an imaginary interval by its
midpoint is used. The cells themselves need not be conjugation invariant.
\end{proof}

Thus a complete finite census may contain more useful information than
one scalar subsequently evaluated from it. Ordered disjoint real cells
give interval enclosures for all internal adjacent gaps and, away from the
two end edges, for neighboring triples of gaps. The numerical census of
\cite{BakerOCCensus22} illustrates this deduction in
Appendix~\ref{app:compact-example}, at its inherited winding/count scope.
The deduction supplies neither global reality at time zero nor a
positive-time source or root-velocity bound.

\subsection{A local forward interval without global reality}

Fix a bounded conjugation-invariant domain $\Omega$ with a simple,
piecewise smooth boundary $\Gamma$. Assume the complete heat field is
nonzero on a neighborhood of $\Gamma$ for $0\le t\le s$ and has exactly
$N\ge2$ zeros inside, counting multiplicity. All source regularity and
clearance assumptions concern this whole interval. Initially suppose
these $N$ zeros are real and simple. Exterior zeros may be nonreal.

On a simple-root subinterval write the ordered interior labels as
$x_1<\cdots<x_N$ and define the whole-block deflation
\begin{equation}
 R(t,z)=B(t,z)-\sum_{j=1}^{N}\frac1{z-x_j(t)},\qquad
 U_i=\sum_{j\ne i}\frac1{x_i-x_j}+R(t,x_i),\qquad
 x_i'=2U_i.
 \label{fwd:deflation}
\end{equation}
The definition uses the same complete $B=H_z/H$ as
Section~\ref{sec:lattice}. In particular $R$ retains all exterior poles
and the regularizing terms of the complete field. It is analytic inside
$\Gamma$ after the interior poles have been removed; no independent
infinite first-power sum is introduced.

\begin{proposition}[Forward preservation on a cleared finite domain]
\label{fwd:preservation}
Under the preceding hypotheses the $N$ interior zeros remain real and
simple throughout $0\le t\le s$.
\end{proposition}
\begin{proof}
A simple real zero continues as a real zero by conjugation and local
uniqueness. Suppose there were a first loss of simplicity, and isolate
one of the colliding clusters of $k\ge2$ real labels. Its limiting point
is inside $\Omega$: the time-wide boundary clearance prevents escape or
entry. Choose a small fixed circle separating this cluster from every
other zero near the putative event. Removing all its poles gives an
analytic regular part $R_{\rm cl}$, uniformly bounded near the cluster
up to the event. For example, the monic polynomial of the enclosed roots
has analytic symmetric coefficients obtained from the counted contour;
the remaining factor is analytic and nonzero there.

Set $\bar x=k^{-1}\sum_{i=1}^k x_i$ and
$\mathcal V=\sum_{i=1}^k(x_i-\bar x)^2$. Since the centered coordinates
sum to zero, the node law gives
\begin{align}
 \mathcal V'
 &=4\sum_i(x_i-\bar x)
       \left(\sum_{j\ne i}\frac1{x_i-x_j}+R_{\rm cl}(t,x_i)\right)
       \notag\\
 &=2k(k-1)+4\sum_i(x_i-\bar x)R_{\rm cl}(t,x_i).
 \label{fwd:cluster-variance}
\end{align}
For the second equality, pair the two ordered contributions of each
unordered pair: their quotient is
$[(x_i-\bar x)-(x_j-\bar x)]/(x_i-x_j)=1$.
At the proposed collapse the last term tends to zero, so
$\mathcal V'>0$ sufficiently close to the event. A nonnegative
$\mathcal V$ tending to zero from the left cannot have that property.
This excludes every cluster size, including simultaneous disjoint
collisions. The local continuation and the compact boundary clearance
then cover the entire interval.
\end{proof}

For an isolated pair the same sign appears as $D'=2$ at $D=0$.
Proposition~\ref{fwd:preservation} includes higher clusters and does not
assume reality of the complete divisor. It is nevertheless conditional
on the independent time-wide contour certificate. A time-zero census
does not acquire that certificate, spatial buffers, or any numerical
lower bound for the later gaps.

\subsection{Dissipation of common node-force dispersion}

Let $P=I-\mathbf1\mathbf1^{T}/N$ be the orthogonal projection onto
mean-zero vectors. Define a positive weighted graph Laplacian by
\begin{equation}
 (Lv)_i=\sum_{j\ne i}\frac{v_i-v_j}{(x_i-x_j)^2},\qquad
 h_i=(R_t+2U_iR_z)(t,x_i),\qquad
 a=PU,\quad r_U=\lVert a\rVert_2.
 \label{fwd:laplacian}
\end{equation}
The derivatives of $R$ are taken at a fixed spatial argument before
evaluation at $x_i$. All external and block-boundary effects remain in
$h$.

\begin{proposition}[Force dispersion in the two time directions]
\label{fwd:force-identity}
On a common simple-root interval,
\begin{align}
 U_t&=-2LU+h,\notag\\
 \frac{\dd}{\dd t}\frac{r_U^2}{2}
 &=-2\sum_{i<j}\frac{(U_i-U_j)^2}{(x_i-x_j)^2}
       +(PU)\mathbin{\cdot}(Ph).
 \label{fwd:energy}
\end{align}
In backward time $u=t_*-t$, put
$f_i=(R_u-2U_iR_z)(u,x_i)=-h_i(t_*-u)$. Then
\begin{align}
 U_u&=2LU+f,\notag\\
 \frac{\dd}{\dd u}\frac{r_U^2}{2}
 &=2\sum_{i<j}\frac{(U_i-U_j)^2}{(x_i-x_j)^2}
       +(PU)\mathbin{\cdot}(Pf).
 \label{fwd:backward-energy}
\end{align}
\end{proposition}
\begin{proof}
Differentiate each finite reciprocal in \eqref{fwd:deflation}, using
$x_i'=2U_i$:
\[
 \frac{\dd}{\dd t}\frac1{x_i-x_j}
 =-\frac{2(U_i-U_j)}{(x_i-x_j)^2}.
\]
The derivative of $R(t,x_i(t))$ is $h_i$, proving the vector equation.
Since $L\mathbf1=0$, $a\cdot LU=a\cdot La$; pairing $i,j$ gives
$a\cdot La=\sum_{i<j}(U_i-U_j)^2/(x_i-x_j)^2$.
Taking the scalar product with $a$ proves \eqref{fwd:energy}.
Reversing physical time proves \eqref{fwd:backward-energy}.
\end{proof}

The internal term is dissipative forward and antidissipative backward.
Neither sign is a sign theorem for the complete derivative: the signed
exterior term can reinforce or cancel it.

\begin{corollary}[Conditional forward force contraction]
\label{fwd:force-contraction}
Suppose an independent spatial enclosure gives
$x_N(t)-x_1(t)\le L_{\max}$ and $J(t)\ge\lVert Ph(t)\rVert_2$ on the
whole interval. With $\gamma=2N/L_{\max}^2$,
\begin{equation}
 r_U(t)\le \e^{-\gamma t}r_U(0)
       +\int_0^t\e^{-\gamma(t-v)}J(v)\dd v.
 \label{fwd:contraction}
\end{equation}
\end{corollary}
\begin{proof}
Every internal weight is at least $L_{\max}^{-2}$, and for a mean-zero
vector $\sum_{i<j}(a_i-a_j)^2=N\lVert a\rVert_2^2$. Thus
\eqref{fwd:energy} gives
$(r_U)_t\le-\gamma r_U+J$ wherever $r_U>0$.
At $r_U=0$ the upper right derivative of the norm is at most $\lVert Ph\rVert_2$,
so the same comparison holds in the upper-Dini sense. An integrating
factor proves \eqref{fwd:contraction}.
\end{proof}

For a constant $J$, the upper envelope is
$\e^{-\gamma t}r_U(0)+(J/\gamma)(1-\e^{-\gamma t})$.
It improves the initial norm only if its source floor is sufficiently
small. Without an acquired $J$, internal dissipation alone provides no
canonical contraction rate. The corresponding backward unsigned bound is
\begin{equation}
 (r_U)_u\le2\lambda_{\max}(L)r_U+\lVert Pf\rVert_2.
 \label{fwd:backward-norm}
\end{equation}
Replacing the denominators by one lower gap stop can make this scalar
allowance large; the fixed compact screen in
Appendix~\ref{app:compact-example} is scoped to that replacement and an
unsigned exterior norm.

Force dispersion is a velocity statistic. Indeed,
\begin{equation}
 g_i(t)-g_i(0)=2\int_0^t(U_{i+1}-U_i)\dd v,
 \qquad |U_{i+1}-U_i|\le\sqrt2\,r_U.
 \label{fwd:gap-integral}
\end{equation}
A small force dispersion need not imply a uniform gap profile: the
exterior may balance a nonuniform configuration. Direct common-source
node evolution retains the signed differences in
\eqref{fwd:gap-integral}; given the same node-force path information, the
norm replacement cannot improve that directional bound. No acquisition
cost advantage is asserted.

\subsection{A forward lower bound for the actual block minimum}

Assume the real segment connecting the two endpoint nodes lies inside
the independently cleared domain. Put
\begin{equation}
 g_i=x_{i+1}-x_i,\qquad D_i=g_i^2/4,\qquad
 q(t)=\min_{1\le i<N}D_i(t).
 \label{fwd:minimum}
\end{equation}
The minimum includes \emph{all} internal edges, including both end edges.
Split the complete interaction into the other real nodes in this block
and the complete exterior. Comparing \eqref{fwd:deflation} at two adjacent
nodes with the gap law gives
\begin{equation}
 M_i=M_{{\rm int},i}+M_{{\rm ext},i},\qquad
 M_{{\rm ext},i}=-\frac{R(t,x_{i+1})-R(t,x_i)}{g_i}
       =-\int_0^1R_z(t,x_i+\theta g_i)\dd\theta.
 \label{fwd:exterior}
\end{equation}
This is a complete signed exterior contribution. In particular,
$|R_z|\le C$ on every connecting segment implies
$M_{{\rm ext},i}\le C$, without exterior reality.

\begin{theorem}[Forward transport of a finite-block minimum]
\label{fwd:minimum-transport}
Suppose $M_{{\rm ext},i}\le C$ for every internal edge throughout the
forward interval, where $C\ge0$. Then, with
\begin{equation}
 a_N=\min_{1\le i<N}\left(\frac1i+\frac1{N-i}\right)
     =\frac{N}{\lfloor N/2\rfloor\lceil N/2\rceil},
 \label{fwd:finite-constant}
\end{equation}
the minimum obeys almost everywhere
\begin{equation}
 q_t\ge a_N-4Cq.
 \label{fwd:minimum-law}
\end{equation}
Consequently
\begin{equation}
 q(t)\ge\Phi_t(q(0)),\qquad
 \Phi_t(d)=
 \begin{cases}
 \e^{-4Ct}d+\dfrac{a_N}{4C}(1-\e^{-4Ct}),&C>0,\\[5pt]
 d+a_Nt,&C=0.
 \end{cases}
 \label{fwd:minimum-bound}
\end{equation}
\end{theorem}
\begin{proof}
At an active minimizing edge $i$, every block gap is at least $g_i$.
The $m$th retained node to the left or right of this pair is therefore
at distance at least $mg_i$ from its nearer endpoint. The positive
real-node terms in the interaction satisfy
\begin{align}
 M_{{\rm int},i}
 &\le\frac1{g_i^2}\left(
       \sum_{m=1}^{i-1}\frac1{m(m+1)}
       +\sum_{m=1}^{N-1-i}\frac1{m(m+1)}\right)\notag\\
 &=\frac1{g_i^2}\left(2-\frac1i-\frac1{N-i}\right).
 \label{fwd:internal-bound}
\end{align}
The empty sums at the end edges have value zero, so this also covers
$N=2$. Since $4D_i=g_i^2$, the physical pair law gives at an active edge
\[
 (D_i)_t=2-g_i^2M_i
 \ge\frac1i+\frac1{N-i}-4CD_i.
\]
Each $D_i$ is continuously differentiable, and the minimum of finitely
many such functions is locally Lipschitz. At almost every time its
derivative is the derivative of an active edge. The preceding bound
holds for every active edge, including ties, so it gives
\eqref{fwd:minimum-law} without a fixed minimizing label. Multiplication
by $\e^{4Ct}$ and integration proves \eqref{fwd:minimum-bound}.
The expression in \eqref{fwd:finite-constant} is minimized at the central
integer or integers, which gives the stated constant.
\end{proof}

For an acquired lower anchor $q(0)\ge d_0>0$, the scalar lower bound
$\Phi_t(d_0)$ improves $d_0$ at every $t>0$ precisely when
$C<a_N/(4d_0)$, with the stated continuous interpretation at $C=0$.
This is a test of a proposed source bound, not its acquisition. The
positive right side of \eqref{fwd:minimum-bound} can alternatively be used
in a first-lower-exit argument on an independently acquired outer domain.

The choice of minimum is essential. If only selected interior edges
were minimized, an excluded boundary gap could be smaller than the active
edge, invalidating the distance bound used in
\eqref{fwd:internal-bound}. A smaller observational window does not
automatically give its buffer to the outermost nodes of a larger census.
One must use the complete block minimum or pay a different boundary
comparison.

\begin{corollary}[Empirical lower-tail transport]
\label{fwd:empirical}
For finitely many blocks with the same $N$ and a common paid $C$, sample
the block index $X$ uniformly. For $C>0$,
\begin{equation}
 \mathbb P\{q_X(t)\le\xi\}
 \le\mathbb P\left\{q_X(0)\le
       \e^{4Ct}\xi-\frac{a_N}{4C}(\e^{4Ct}-1)\right\}.
 \label{fwd:tail}
\end{equation}
For $C=0$ the threshold on the right is $\xi-a_Nt$.
\end{corollary}
\begin{proof}
The increasing map $\Phi_t$ in \eqref{fwd:minimum-bound} gives the
pointwise event inclusion
$\{q_X(t)\le\xi\}\subseteq\{q_X(0)\le\Phi_t^{-1}(\xi)\}$.
Counting the finite block indices proves the claim. No independence or
random-matrix law is used.
\end{proof}

Different block sizes or exterior bounds use their respective maps
$\Phi_t$. A global zero-exception conclusion additionally needs coverage
of every eligible pair, including pairs crossing block boundaries, and
the finite integer threshold of Section~\ref{sec:statistics}. A small
empirical bad-block rate is not zero by itself; no varying-height or
all-height law is inferred from this finite statement.

\subsection{Complete-source acquisition of the exterior derivative}

The exterior estimate in Theorem~\ref{fwd:minimum-transport} can be
formulated on fixed trial locations, without assuming that the roots
remain positive-distance apart. For $w$ inside $\Gamma$,
\begin{equation}
 R(t,w)=\frac1{2\pi\ii}\oint_\Gamma\frac{B(t,z)}{z-w}\dd z,
 \qquad
 R_z(t,w)=\frac1{2\pi\ii}\oint_\Gamma
                         \frac{B(t,z)}{(z-w)^2}\dd z.
 \label{fwd:source-traces}
\end{equation}
To prove these formulas, split $B$ into $R$ and its finite set of
interior poles, with multiplicities. Cauchy's formula applies to $R$.
For each interior pole $\rho$, the rational functions
$[(z-\rho)(z-w)]^{-1}$ and
$[(z-\rho)(z-w)^2]^{-1}$ have all poles inside and decay at least as
$z^{-2}$ at infinity, so their contour integrals vanish. This remains
true at $w=\rho$ by the same rational argument or by continuation.
The count and full source determine the traces even if the interior
configuration is temporarily multiple or nonreal.

Suppose a circle of radius $r_\Gamma$ is used and every trial point $w$
is at distance at least $\delta>0$ from it. A uniform complete-source
error $e_B$ on the circle contributes at most
\begin{equation}
 E_R\le r_\Gamma e_B/\delta,\qquad
 E_{R_z}\le r_\Gamma e_B/\delta^2.
 \label{fwd:trace-error}
\end{equation}
These follow directly from contour length and the two kernels.
The quotient estimate \eqref{cmp:quotient-error} supplies $e_B$ only
after field, first-jet and positive contour-denominator errors are paid.
One must add quadrature, arithmetic, trial-location and genuine time-cell
errors; the angular-strip estimate \eqref{cmp:trapezoid} applies only
when the full integrand is cleared and bounded on that strip. A bound
for $|R_z|$ must hold on the whole trial segment containing all possible
connecting edges, not just at sampled nodes. The complete source error
already pays the exterior; no omitted zero tail is then silently set
to zero. No such favorable time-wide $C$ is acquired by these identities.

The force statistic has additional requirements. The fixed-coordinate
time derivative of the first trace in \eqref{fwd:source-traces}, when
its time-uniform source differentiation is justified, gives $R_t$ from
$B_t$ on the same contour. Alternatively one can evaluate the complete
source jets at simple roots. With a fixed scalar normalization $G$ of
$H$, differentiation along a root gives
\begin{equation}
 U_i=\frac{G_{zz}}{2G_z}(t,x_i),\qquad
 (U_i)_t=-\frac{G_{zzzz}}{2G_z}
       +\frac{G_{zz}G_{zzz}}{G_z^2}
       -\frac{G_{zz}^3}{2G_z^3}\quad\hbox{at }x_i.
 \label{fwd:source-jets}
\end{equation}
Indeed $\dd G_z(t,x_i)/\dd t=-G_{zzz}+G_{zz}^2/G_z$ and
$\dd G_{zz}(t,x_i)/\dd t=-G_{zzzz}+G_{zzz}G_{zz}/G_z$;
the quotient rule proves the second formula.
This route requires nonzero $G_z$ denominators and root-location errors.
For approximate node jets $h_1,h_2$ with errors $e_1,e_2$ and
$|h_1|\ge b_1>e_1$,
\begin{equation}
 \left|U-\frac{h_2}{2h_1}\right|
 \le\frac{e_2}{2(b_1-e_1)}
       +\frac{|h_2|e_1}{2b_1(b_1-e_1)}.
 \label{fwd:node-error}
\end{equation}
Higher quotient products in \eqref{fwd:source-jets} require their own
joint enclosure. If two exact and approximate positive node distances
are both at least $d_*$ and differ by at most $e$, their inverse-square
weights differ by at most $2e/d_*^3$. That error also enters a computed
Laplacian. Neither the common field nor mean subtraction removes these
reciprocal and boundary payments.

\subsection{What the forward estimate can pay on returning backward}

The internal estimate \eqref{fwd:internal-bound} is an \emph{upper}
interaction bound. Backward gap protection needs a lower interaction or
forcing bound, so it cannot reuse that inequality with its sign reversed.
For example, a two-sided exterior bound $|M_{{\rm ext},i}|\le C$ gives
only the coarse lower bound $M_i\ge-C$. On a paid backward interval it
therefore gives $D_u\ge-2-4CD$. If this is composed with
\eqref{fwd:minimum-bound} over equal forward and backward durations $s$,
and $z=\e^{-4Cs}$, the resulting scalar allowance is
\begin{equation}
 \mathcal Q_{\rm return}:=z^2q(0)
       +\frac{(1-z)(a_Nz-2)}{4C},\qquad C>0.
 \label{fwd:coarse-return}
\end{equation}
Here $\mathcal Q_{\rm return}$ is the lower allowance supplied by this
composition, not the actual return, which retraces the same trajectory.
To obtain \eqref{fwd:coarse-return}, the backward integrating factor
gives $D(s)\ge zD(0)-(1-z)/(2C)$; substitute the forward lower
anchor $zq(0)+a_N(1-z)/(4C)$. Since $a_N\le2$, the displayed right
side is at most $q(0)$. At $C=0$ the allowance is
$q(0)+(a_N-2)s$, with the same conclusion. Thus this particular pair
of scalar allowances does not turn a forward excursion into a better
returned lower bound.

The sharper comparison for the unchanged exact-initial-interaction
horizon is proved in Theorem~\ref{rf:horizon}. It likewise separates
an improved positive-time datum from an improved negative endpoint.
A useful forward acquisition must preserve additional joint information
for the backward consumer, such as the complete time-dependent forcing,
and give the direct source or regular $D/M$ comparison the same data.
The local forward theorem legitimately applies to noncanonical real
blocks as well; such controls may collide at an earlier backward time.
Its generic validity and the exact theta source identity are not, by
themselves, a canonical backward forcing estimate.

% 18 November 2026. Drafted for Edward Baker with OpenAI Codex assistance.
% Actual model and effort unavailable; authoring-helper settings inherited.
\section{Retained forcing and the exact backward return allowance}
\label{sec:retained-forcing}

A whole-time forcing estimate carries more information than an initial
interaction or a constant upper bound for the defect. We now keep that
information in two scalar comparisons and then identify exactly what is
lost when a forward evolution is compressed back to its initial-$M$
backward horizon. The derivations are from \cite{BakerCM32,BakerCM34},
with the internal assessments \cite{BakerCMReview29,BakerCMReview30}.
The numerical anchor tests are deferred to Appendix~\ref{app:compact-example}.
All statements in this section are analytical and independent of those
supporting values.

Work first in physical backward time $u=t_*-t$ on a counted pair tube
where the complete local equations hold. Write
\begin{equation}
 D_u=-2+4DM,\qquad M_u=-2M^2+\mathcal F,\qquad
 \mathcal F=12S_2+16DS_3,\qquad K=2D^2\mathcal F.
 \label{rf:pair-system}
\end{equation}
The pair is real and distinct while $D>0$. Its exterior may be nonreal:
the complete signed quantities in \eqref{rf:pair-system} are then
acquired through Section~\ref{sec:compact}, rather than by omitting an
uncounted tail. With $V=1-2DM$, the identities
\begin{equation}
 D_u=-2V,\qquad V_u=\frac{1-V^2-K}{D},\qquad
 \left(\frac{1-V^2}{D}\right)_u=4V\mathcal F
                                  =-2\mathcal F D_u
 \label{rf:transient-identity}
\end{equation}
follow from \eqref{rf:pair-system}. In the last identity,
$(1-V^2)/D=4M(1-DM)$; differentiating this expression proves the
formula without assigning separate bounds to its factors. We will not
divide its signed integral by an uncertain minimum weight. The
comparisons below use the differential equations directly.

\subsection{The sharp scalar comparison for a constant \texorpdfstring{$K$}{K} floor}

\begin{theorem}[Constant forcing floor with the initial transient retained]
\label{rf:constant-K}
Suppose \eqref{rf:pair-system} holds on a stopped interval with $D>0$,
and suppose $K\ge\kappa$ for a constant $\kappa\le1$. Let $d>0$ and
$v\in\R$ satisfy the two ordered initial conditions
\begin{equation}
 g(0)\ge2\sqrt d,\qquad
 g_u(0)\ge-2v/\sqrt d,\qquad g=2\sqrt D.
 \label{rf:ordered-data}
\end{equation}
Define
\begin{equation}
 d_\kappa(u)=d-2vu+\frac{v^2+\kappa-1}{d}u^2.
 \label{rf:constant-K-polynomial}
\end{equation}
Then $D(u)\ge d_\kappa(u)$ before the first zero of $d_\kappa$ and
before the stipulated pair tube ends. Equality holds for the scalar
comparison with constant $K=\kappa$ and equal initial data. Thus this
comparison retains the full initial transient of that scalar problem.
\end{theorem}

\begin{proof}
Differentiate $g=2\sqrt D$ using \eqref{rf:transient-identity}. The
terms quadratic in $V$ cancel and give
\begin{equation}
 g_{uu}=\frac{16(K-1)}{g^3}
       \ge f_\kappa(g),\qquad
 f_\kappa(g)=\frac{16(\kappa-1)}{g^3}.
 \label{rf:gap-acceleration}
\end{equation}
For $\kappa\le1$, $f_\kappa$ is nondecreasing on $(0,\infty)$.
Direct differentiation of $y=2\sqrt{d_\kappa}$ gives
$y_{uu}=f_\kappa(y)$, $y(0)=2\sqrt d$ and
$y_u(0)=-2v/\sqrt d$, on the initial interval where $d_\kappa>0$.

For completeness, the comparison can be proved on any compact
subinterval on which $g,y>0$. Put $a=g-y$ and $b=g_u-y_u$.
Then $a_u=b$ and
\[
 b_u\ge f_\kappa(g)-f_\kappa(y)=c(u)a,
\]
where the divided difference $c(u)$ is continuous, nonnegative and
bounded, using $f_\kappa'$ at equal arguments. The negative parts obey
$(a_-)_u\le b_-$ and $(b_-)_u\le c(u)a_-$ almost everywhere.
An integrating-factor inequality for $a_-+b_-$, initially zero, makes
both negative parts zero. Hence $g\ge y$, and squaring proves the
assertion. Exhaust the positive comparison interval by such compact
subintervals. No comparison is continued through the first zero of
$d_\kappa$ merely because its quadratic expression has values later.
\end{proof}

The initial conditions in \eqref{rf:ordered-data} are stated for both
position and velocity deliberately. If $D(0)\ge d$ and $V(0)\le v$
with $v\ge0$, they follow at once. A useful alternative uses
$D(0)\ge d$ and $M(0)\ge b\ge0$. Indeed,
\[
 g_u(0)=-\frac2{\sqrt{D(0)}}+4M(0)\sqrt{D(0)}
       \ge-\frac2{\sqrt d}+4b\sqrt d,
\]
because the last expression is increasing in $D$ and $M$ on that
rectangle. Thus one may take $v=1-2db$, without requiring this $v$ to
be nonnegative. These ordering checks also identify the favorable
endpoints when screening an uncertain anchor.

The equality claim is sharpness within the stated scalar differential
comparison, not a claim that every constant-$K$ scalar trajectory is a
canonical divisor. A strictly positive minimum of $d_\kappa-d_{\rm stop}$
on a prescribed finite interval excludes a lower stop, once all other
tube exits have been independently excluded. Conversely, a comparison
that reaches zero before the requested endpoint cannot certify that
endpoint by this theorem. It does not establish an actual collision.

\subsection{A constant physical forcing floor}

A constant floor for $\mathcal F$ is a different input from a constant
floor for $K=2D^2\mathcal F$. In particular, a small $K$ near a small
gap is compatible with a positive physical forcing. The regular
$D/M$ equations retain this distinction without division by $D$.

\begin{theorem}[Regular Riccati and gap comparison]
\label{rf:forcing-comparison}
Suppose \eqref{rf:pair-system} holds on a stopped positive-$D$ interval,
with complete signed $\mathcal F\ge f\ge0$, $D(0)\ge d>0$ and
$M(0)\ge b\ge0$. For $f>0$ put $\lambda=\sqrt{2f}$ and
\begin{equation}
 S(u)=\frac{\sinh(\lambda u)}{\lambda},\qquad
 W(u)=\cosh(\lambda u)+2bS(u),\qquad
 m(u)=\frac{W'(u)}{2W(u)}.
 \label{rf:model-functions}
\end{equation}
For $f=0$ use the continuous limits $S(u)=u$,
$W(u)=1+2bu$ and $m(u)=b/(1+2bu)$. In either case define
\begin{equation}
 d_f(u)=W(u)\bigl[dW(u)-2S(u)\bigr].
 \label{rf:physical-forcing-model}
\end{equation}
Then $M(u)\ge m(u)\ge0$ and $D(u)\ge d_f(u)$ throughout the stopped
interval. The direct regular $D/M$ comparison given the same
$(d,b,f)$ produces exactly these same functions.
\end{theorem}

\begin{proof}
For $u\ge0$, $W\ge1$ and $W'\ge0$. For $f>0$, $W''=2fW$;
the same formula holds at $f=0$ with $W''=0$. Therefore
\[
 m_u=\frac{W''W-(W')^2}{2W^2}=f-2m^2,
 \qquad m(0)=b.
\]
The complete forcing inequality gives
\[
 (M-m)_u+2(M+m)(M-m)\ge0.
\]
On each compact stopped subinterval an integrating factor is positive,
irrespective of the sign of the continuous coefficient. Since
$M(0)-m(0)\ge0$, this proves $M\ge m\ge0$; no exterior-summand
positivity was used.

The Wronskian identity $S'W-SW'=1$ follows directly from the hyperbolic
formulas, and is also valid for $f=0$. Consequently
\[
 (S/W)'=W^{-2},\qquad
 d_f=W^2\left[d-2\int_0^u W(s)^{-2}\dd s\right],\qquad
 (d_f)_u=-2+4m d_f.
\]
Subtracting from the actual equation yields
\[
 (D-d_f)_u-4m(D-d_f)=4D(M-m)\ge0
\]
while $D>0$. A second integrating factor proves $D\ge d_f$.
Both integrations use the regular pair system itself, which proves the
same-information statement as well. No monotonicity of the actual
$D$, and no sign of its actual $V$, is assumed.
\end{proof}

The formula covers the three Riccati regimes without a singular case
distinction. For $f>0$ put $a=\sqrt{f/2}$. If $b>a$, then $m$ decreases
to $a$; if $b=a$, it is constant; if $0\le b<a$, it increases to $a$.
These assertions follow from $m_u=2(a^2-m^2)$ and uniqueness, or from
\eqref{rf:model-functions}. For $f=0$, $m$ decreases to zero when
$b>0$ and is identically zero when $b=0$. In particular,
\begin{equation}
 d_0(u)=(1+2bu)\bigl[d(1+2bu)-2u\bigr],
 \label{rf:zero-forcing-limit}
\end{equation}
recovering the earlier initial-interaction comparison exactly.

\begin{corollary}[A whole-interval lower-stop test]
\label{rf:forcing-stop}
In addition to the hypotheses of Theorem~\ref{rf:forcing-comparison},
suppose an independently cleared count-two contour and unconditional
coordinate enclosures exclude every other exit on $0\le u\le\Delta$.
Acquire $\mathcal F_\Gamma\ge f$ on the corresponding stopped trial
region $D\ge d_{\rm stop}>0$. If
\begin{equation}
 \min_{0\le u\le\Delta}d_f(u)>d_{\rm stop},
 \label{rf:stop-margin}
\end{equation}
then the actual pair stays above that lower stop throughout the interval.
Positivity of $d_f(\Delta)$ alone always implies positivity of $d_f$
on $[0,\Delta]$. The stronger endpoint test
$d_f(\Delta)>d_{\rm stop}$ implies \eqref{rf:stop-margin} whenever
$b\ge\sqrt{f/2}$ and $4bd<2$.
\end{corollary}

\begin{proof}
At a hypothetical first lower exit, the actual coordinates have stayed
in the acquired trial region. Theorem~\ref{rf:forcing-comparison}
therefore gives $D\ge d_f>d_{\rm stop}$ there, a contradiction.
The normalized model satisfies
\[
 \frac{d_f(u)}{W(u)^2}=d-2S(u)/W(u),\qquad
 \left(\frac{d_f}{W^2}\right)'=-2/W^2<0.
\]
Thus a positive endpoint makes the normalized model, and hence $d_f$,
positive at every earlier time. For the stated stronger endpoint test,
$0\le m\le b$. The positive model cannot cross its initial upper
barrier $d$, since at such a crossing its derivative is at most
$-2+4bd<0$. While positive it consequently satisfies
$(d_f)_u\le-2+4bd<0$, so its minimum is its endpoint. In the other
Riccati regimes this last monotonicity is not asserted; use the actual
scalar minimum in \eqref{rf:stop-margin} instead.
\end{proof}

The count and coordinate assumptions in this corollary are the same
independent gates as in Theorem~\ref{cmp:stop}. In particular, coordinate
acquisition may permit negative $D$; the desired positive trajectory is
not an input to it. For a finite family of pair charts the margin must
hold for every covered label and every chart/boundary exit. Passing only
one initially minimizing label does not protect an evolving block
minimum. A regular merging-pair chart also does not regularize its
neighboring edge charts.

\subsection{Complete source acquisition and its error budget}

The trial trace in Proposition~\ref{cmp:trial} gives the combined
physical forcing directly:
\begin{align}
 \mathcal F_\Gamma(t;\widehat c,\widehat D)
 &=-\frac1{2\pi\ii}\oint_\Gamma B(t,z)
       \left(\frac{12}{Q^2}+\frac{16\widehat D}{Q^3}\right)\dd z,
       \notag\\
 &=-\frac1{2\pi\ii}\oint_\Gamma\mathcal L(t,z)Q'(z)
       \left(\frac{24}{Q^3}+\frac{48\widehat D}{Q^4}\right)\dd z,
 \label{rf:source-kernel}
\end{align}
where the second identity requires the matching cleared annular
logarithm of \eqref{cmp:log-kernel}. All interior rational contributions
cancel as in Section~\ref{sec:compact}; every exterior occurrence and
its integer multiplicity remain. Conjugation makes the value real, but
does not make the individual exterior summands nonnegative.

Use the geometry of \eqref{cmp:kernel-errors}: $0\le C<r$,
$|\widehat D|\le D_{\max}$ and $q=(r-C)^2-D_{\max}>0$.
Uniform source errors on the circle give
\begin{align}
 E_{\mathcal F,B}
 &\le re_B\left(\frac{12}{q^2}+\frac{16D_{\max}}{q^3}\right),
       \notag\\
 E_{\mathcal F,\mathcal L}
 &\le2r(r+C)e_{\mathcal L}
       \left(\frac{24}{q^3}+\frac{48D_{\max}}{q^4}\right).
 \label{rf:source-errors}
\end{align}
The quotient estimate \eqref{cmp:quotient-error} pays $e_B$ from
field/first-jet errors and a nonzero contour denominator. These inequalities
follow by taking the modulus of the combined kernel and integrating the
contour length. They need no reciprocal target gap.

If trial coordinates are approximated rather than enclosed as parameters,
suppose their errors are at most $e_c,e_D$ and all connecting parameter
segments retain the same $q$. With $B_{\max}$ bounding the exact field
$|B|$ on the contour, the additional error is at most
\begin{equation}
 \begin{split}
 E_{\rm par}\le rB_{\max}\bigg[&
 \left(\frac{48(r+C)}{q^3}
             +\frac{96D_{\max}(r+C)}{q^4}\right)e_c\\
 &+\left(\frac{40}{q^3}+\frac{48D_{\max}}{q^4}\right)e_D\bigg].
 \end{split}
 \label{rf:parameter-error}
\end{equation}
To verify the constants, differentiate
$J=12Q^{-2}+16\widehat DQ^{-3}$, using
$Q_{\widehat c}=-2(z-\widehat c)$ and $Q_{\widehat D}=-1$:
\[
 J_{\widehat c}=48(z-\widehat c)Q^{-3}
                    +96\widehat D(z-\widehat c)Q^{-4},\qquad
 J_{\widehat D}=40Q^{-3}+48\widehat DQ^{-4}.
\]
Integrate the parameter derivatives along the joining segments and then
around the contour. Evaluating all parameters by a rigorous enclosure
can instead include this payment directly, without adding it twice.

The remaining requirements are precisely those of
\eqref{cmp:trapezoid} and its following source ledger: quadrature on a
cleared angular strip, complete theta/xi and integration tails, arithmetic,
reconstruction, spatial uncertainty and a genuine whole-time remainder.
Heat identities for time derivatives do not remove the remainder.
A valid lower acquisition has the form
\[
 \inf\widehat{\mathcal F}_\Gamma
 -E_{\rm source}-E_{\rm par}-E_{\rm quad}
 -E_{\rm arithmetic}-E_{\rm time}-E_{\rm reconstruction}\ge f,
\]
with a nonoverlapping accounting of errors. It must cover the actual
stopped coordinates and all needed time cells. A bound for
$|\mathcal F|$ is not this signed lower estimate. A contour representation
does not leave an unpaid zero-tail cutoff once the complete source error
has been paid; a finite source implementation still has to pay that error.

Replacing the contour by root jets requires a lower bound for $|H_z|$
and root-location/reciprocal errors. Converting a $D$ error to a gap error
on $D\ge d_{\rm stop}$ costs at most $1/\sqrt{d_{\rm stop}}$ times
that error. Such denominators are not implicit in
\eqref{rf:source-errors}. Negative-time source bounds likewise need
their own domain verification as in \eqref{cmp:negative-gaussian}.
All these inputs must be given equally to direct regular $D/M$ evolution
or direct contour-discriminant acquisition. Theorem~\ref{rf:forcing-comparison}
claims no additional margin or cost advantage over those methods.

\subsection{Forward reanchoring and the initial-interaction horizon}

Return to forward physical time $t$. This subsection concerns one
specific consumer: at a starting slice, discard the positive forcing
after retaining the exact initial $D,M$, and use
\eqref{rf:zero-forcing-limit} backward. On $D>0$, $M\ge0$ and $DM<1$
its strict positive-gap horizon is
\begin{equation}
 h(t)=\frac{D(t)}{2(1-D(t)M(t))}.
 \label{rf:horizon-definition}
\end{equation}
This is an allowance, not the actual time remaining until a collision.
The consumer remains conditional on its backward source/sign regime and
coverage; an algebraic horizon alone does not acquire that regime.

\begin{theorem}[Exact evolution of the initial-$M$ horizon]
\label{rf:horizon}
On a forward local interval suppose the complete pair equations hold,
$D>0$, $M\ge0$ and $\mathcal F\ge0$. Assume $1-D(0)M(0)>0$.
Then $1-DM>0$ throughout that interval, and
\begin{equation}
 h_t=1-2\mathcal Fh^2,\qquad
 h(s)-s=h(0)-2\int_0^s\mathcal F(t)h(t)^2\dd t\le h(0).
 \label{rf:horizon-identity}
\end{equation}
Consequently, forwarding to $s$ and applying the unchanged
complete-initial-$M$ consumer there cannot reach an earlier endpoint
than the same exact-input consumer at time zero.
\end{theorem}

\begin{proof}
The forward equations are $D_t=2-4DM$ and
$M_t=2M^2-\mathcal F$. Put $y=1-DM$. Direct differentiation gives
\[
 y_t=-D_tM-DM_t=-2My+D\mathcal F.
\]
On every compact forward subinterval, multiplying by
$\exp(2\int_0^tM)$ proves $y(t)>0$ from $y(0)>0$, because
$D\mathcal F\ge0$. Thus the denominator of $h=D/(2y)$ does not
cross zero. Differentiating the quotient now gives
\[
 h_t=\frac{D_ty-Dy_t}{2y^2}
     =\frac{2y^2-D^2\mathcal F}{2y^2}
     =1-2\mathcal Fh^2.
\]
Integration proves \eqref{rf:horizon-identity}. The strict backward
allowance from time $s$ reaches only times greater than $s-h(s)$,
and the displayed identity gives $s-h(s)\ge-h(0)$.
\end{proof}

The hypothesis $M\ge0$ here is explicit throughout the forward
interval. It is not inferred from $\mathcal F\ge0$: at $M=0$ the
equation $M_t=-\mathcal F$ can point toward negative values. A complete
all-real adjacent-pair regime supplies both signs, while a local tube
with possibly nonreal exterior needs them separately. The raw quotient
identity holds more generally whenever $D>0$ and $DM\ne1$; the theorem
keeps the domain of the inherited nonnegative-initial-$M$ consumer.

For any persistent finite family of eligible labels, apply
\eqref{rf:horizon-identity} to each label. Every $h_j(t)-t$ is
nonincreasing, so is their minimum on each fixed block, and the sets of
block indices whose allowance exceeds a fixed deadline threshold can
only shrink. Repeated exact reanchoring at increasing forward times
within this same domain does not change that conclusion. A changed
minimizing label creates no improvement, since the pointwise inequalities
hold for all labels simultaneously. Coverage and outer-boundary exits
are still separate hypotheses.

The same ceiling applies to a nonnegative two-neighbor lower estimate
$b_3(s)\le M(s)$ only when that inequality is actually paid. For fixed
$D$, the horizon $D/[2(1-Db)]$ is increasing in $b$ on $Db<1$, so the
corresponding horizon is at most $h(s)$. In a complete all-real regime,
positive omission supplies this lower estimate. If the remaining exterior
may be nonreal, nonnegative complete $M$ alone does not make a positive
real-neighbor subsum a lower bound. Its omitted exterior must be paid.
This comparison is with the \emph{complete exact} initial-$M$ allowance
at time zero; it does not rule out improving a weaker or uncertain
initial estimate.

The quadratic heat collision has $\mathcal F=0$ and saturates the
round-trip identity. For stationary cosine the divisor does not move,
$h$ is constant and $h-t$ decreases: loss of an allowance is therefore
not loss of the actual gap. The exclusion is restricted to the unchanged
initial-$M$ consumer under the stated sign and common-label assumptions.
Signed forcing, the retained-$\mathcal F$ comparison above, time-dependent
joint evolution, and additional source information or acquisition-cost
gains are outside it. The generic identities apply to appropriate
changed-source controls; canonical use still requires the complete
unchanged source and every count, time and error premise.

% 18 November 2026. Prepared for Edward Baker with OpenAI Codex assistance.
% Actual recorded model and effort unavailable. B10 formulation, E31 integration.
\section{Recursive and continuous two-sided refinement}
\label{sec:two-sided}

The preceding sections suggest a more informative use of forward and
backward propagation. A forward source calculation may constrain the
environment of a small gap; a backward comparison may then reduce the
admissible earlier configurations; the reduced configuration set may
sharpen the next forward calculation. The object being refined is an
\emph{enclosure of one source-consistent trajectory}, not the physical
trajectory itself. We develop the formulation proposed in
\cite{BakerCM35}. Its elementary soundness and residual identities are
proved below. Strict improvement for the canonical source remains an
additional quantitative question.

There are two distinct parameters: physical Newman time $t$, and the
iteration number $k$ of an enclosure calculation. Introducing $k$ does
not reverse the physical heat equation, generate another divisor or
supply a finite moment closure. The complete source, the label/count
record and the errors must be shared by the two directions. In
particular, Theorem~\ref{rf:horizon} continues to screen repetition
of the unchanged initial-interaction consumer within its stated class.

\subsection{Common-source state sets and sound refinement}

Fix finite times $t_0<\cdots<t_m$ inside independently cleared counted
contours. Write $\omega$ for a common complete-source record, including
normalization, shared source-error parameters and count/label information.
Let $z_i$ consist of regular observables at $t_i$. For a separated pair
these may include $c,D,M$ and source traces; their admissible domain must
allow both signs of $D$ if positivity is the desired conclusion. For an
overlapping block, choose a valid common chart or larger-cluster symmetric
coordinates when a neighboring edge chart fails. An unexcluded chart or
boundary exit remains a possible failure; it cannot be discarded in the
definition of the admissible set.

Let $E_i^0$ be an outer set of augmented states $(\omega,z_i)$ containing
the true $(\omega_*,z_i^*)$. A validated transition relation
$\mathcal R_i$ contains the true transition and relates only states with
the same source record:
\begin{equation}
 ((\omega,z_i),(\omega,z_{i+1}))\in\mathcal R_i.
 \label{ts:transition}
\end{equation}
This relation may be an outer approximation. It is not asserted to
characterize exactly which finite moment data arise from a canonical
divisor. Define its two projections by
\begin{align}
 \operatorname{Post}_i(A)
 &=\{v:\text{there is }u\in A\text{ with }(u,v)\in\mathcal R_i\},
 \notag\\
 \operatorname{Pre}_i(B)
 &=\{u:\text{there is }v\in B\text{ with }(u,v)\in\mathcal R_i\}.
 \label{ts:projections}
\end{align}
Any new observation set $O_i^k$ must independently contain the true
augmented state. Take it to be the whole state space when no new source
observation is available. One simultaneous sweep is
\begin{equation}
 E_i^{k+1}=E_i^k\cap O_i^k
       \cap\operatorname{Post}_{i-1}(E_{i-1}^k)
       \cap\operatorname{Pre}_{i}(E_{i+1}^k),
 \label{ts:sweep}
\end{equation}
where nonexistent endpoint factors are omitted.

\begin{proposition}[Sound two-sided refinement]\label{ts:sound}
For every finite $k$, the sets in \eqref{ts:sweep} are nested and contain
the true augmented states. Define the full skeleton set
\begin{equation}
 \mathcal P^k=\{(\omega,z_0,\ldots,z_m):
       (\omega,z_i)\in E_i^k\text{ for every }i,
       \ ((\omega,z_i),(\omega,z_{i+1}))\in\mathcal R_i
                         \text{ for every }i<m\}.
 \label{ts:skeleton}
\end{equation}
Then $\mathcal P^{k+1}\subseteq\mathcal P^k$, and every $\mathcal P^k$
contains the true skeleton. For a fixed real statistic $Q$, wherever
the bounds are finite,
\begin{equation}
 \inf_{\mathcal P^{k+1}}Q\ge\inf_{\mathcal P^k}Q,
 \qquad
 \sup_{\mathcal P^{k+1}}Q\le\sup_{\mathcal P^k}Q.
 \label{ts:monotone-bounds}
\end{equation}
More precisely, for the fixed relations in \eqref{ts:transition},
\begin{equation}
 \mathcal P^{k+1}=\{p\in\mathcal P^k:
           (\omega,z_i)\in O_i^k\text{ for every }i\}.
 \label{ts:skeleton-invariance}
\end{equation}
Thus if all observation sets are universal, the full skeleton set and
its exact statistic bounds are unchanged by a sweep.
\end{proposition}

\begin{proof}
Assume the true states belong to all $E_i^k$. Their true neighboring
transitions witness membership in both projections in
\eqref{ts:sweep}, and soundness of $O_i^k$ supplies the remaining new
factor. The intersection therefore retains the true state and is a
subset of $E_i^k$. Induction starts at $k=0$. Substitution into
\eqref{ts:skeleton} proves its corresponding containment and
nonemptiness. Infima increase and suprema decrease when a nonempty
feasible set is restricted, proving \eqref{ts:monotone-bounds}.
For \eqref{ts:skeleton-invariance}, one inclusion follows from the
definition. Conversely, each skeleton in $\mathcal P^k$ supplies its
own neighboring witnesses for every projection in \eqref{ts:sweep}.
If it also satisfies the new observations, all its nodes survive.
\end{proof}

The proposition proves nonworsening enclosures, not strict contraction,
finite termination, a convergence rate or convergence to the exact
feasible set. An implementable projection may use an outward hull of
the exact image; intersecting that hull with the previous enclosure
remains sound if the true image is retained. Projection sweeps alone can
tighten marginal outer estimates by removing states without compatible
neighbors, but cannot improve the exact joint bounds over the already
retained skeleton. New observations or soundly refined transition
relations are needed to shrink that full feasible set. Outward hulls can
lose most of the marginal improvement. Similarly, discarding $\omega$ and choosing
unrelated forcing values on successive cells loses the common-source
constraint. Correlations across nodes and times belong in
\eqref{ts:skeleton}, even when the final reported statistic is scalar.

A statistic of a whole interval needs full admissible cell paths or a
paid between-node interpolation bound. For example, a lower bound for
$\min_{j,t}D_j(t)$ is not obtained from positive values at finitely many
times. A failure-count statistic may include chart exits, block-boundary
events and multiplicities; an upper bound below one then excludes those
events only after the population covers every eligible failure. Neither
random labeling nor a small average supplies this coverage.

\subsection{What reversal alone cannot add}

\begin{proposition}[Exact-flow consistency]\label{ts:inverse}
For a prescribed comparison flow $\Phi_{b,a}$ on a domain where the
forward and reverse solutions are unique,
\begin{equation}
 \Phi_{a,b}\bigl(\Phi_{b,a}(A)\bigr)=A
 \label{ts:inverse-identity}
\end{equation}
for every initial set whose trajectories stay in that domain.
\end{proposition}

\begin{proof}
Uniqueness identifies the reverse of each forward trajectory with its
original trajectory. Thus the two maps compose to the identity on the
stated initial set.
\end{proof}

This is a consistency test for an exact prescribed flow, not an
impossibility theorem for source acquisition. A terminal set obtained
only as the exact forward image of the initial set gives no extra
constraint upon exact reversal. In contrast, interval propagation may
have discarded dependencies which a second sweep recovers, or a later
source calculation may supply an additional checked observation.
Both are legitimate reasons for shrinking an outer set. They must be
compared with direct joint propagation using the same information.
No favorable stability of the inverse is implied by
\eqref{ts:inverse-identity}.

\subsection{An exact continuous forward--adjoint pairing}

For a regular counted pair, retain physical forward time and write
\begin{equation}
 y=\binom D M,\qquad
 y'=G_{\mathcal F}(t,y)
   =\binom{2-4DM}{2M^2-\mathcal F(t)}.
 \label{ts:pair}
\end{equation}
The function $\mathcal F(t)$ denotes the complete forcing along the
\emph{same actual source-consistent path}. Equation~\eqref{ts:pair}
does not say that two moments determine that forcing. Its dependence on
the source, center and exterior remains in $\omega$ and must be acquired
by the contour or hierarchy arguments already developed.

Choose a $C^1$ reference path $\widehat y$ and continuous reference
forcing $\widehat{\mathcal F}$. Put
\begin{equation}
 e=y-\widehat y,\quad
 \delta\mathcal F=\mathcal F-\widehat{\mathcal F},\quad
 r=\widehat y'-G_{\widehat{\mathcal F}}(t,\widehat y),\quad
 \widehat A=
 \begin{pmatrix}-4\widehat M&-4\widehat D\\0&4\widehat M\end{pmatrix}.
 \label{ts:reference}
\end{equation}
The complete error equation is
\begin{equation}
 e'=\widehat A e+
       \binom{-4e_De_M}{2e_M^2-\delta\mathcal F}-r.
 \label{ts:error}
\end{equation}
Indeed, expanding $(\widehat D+e_D)(\widehat M+e_M)$ and
$(\widehat M+e_M)^2$ gives every term displayed in
\eqref{ts:error}. The quadratic remainder has not been dropped.

\begin{proposition}[Two-sided adjoint residual identity]\label{ts:adjoint}
On a compact regular interval $[a,b]$, fix $\tau\in[a,b]$ and a vector
$\ell\in\R^2$. Let
\begin{equation}
 p'=-\widehat A^{\mathsf T}p,\qquad p(\tau)=\ell,
 \qquad
 \Psi=-4p_De_De_M+2p_Me_M^2-p_M\delta\mathcal F-p\cdot r.
 \label{ts:adjoint-def}
\end{equation}
Then
\begin{align}
 \ell\cdot e(\tau)
 &=p(a)\cdot e(a)+\int_a^\tau\Psi(t)\dd t\notag\\
 &=p(b)\cdot e(b)-\int_\tau^b\Psi(t)\dd t.
 \label{ts:two-sided}
\end{align}
The identity has no denominator $D$ and is valid through a target merger
whenever the complete counted-pair tube and regular quantities remain
defined.
\end{proposition}

\begin{proof}
The continuous coefficient matrix gives a unique linear adjoint solution
on the compact interval in both time directions. Apply the product rule
to $p\cdot e$, substitute \eqref{ts:error} and
\eqref{ts:adjoint-def}, and cancel the two linear terms. The result is
$(p\cdot e)'=\Psi$. Integration separately from $a$ to $\tau$ and from
$\tau$ to $b$, with $p(\tau)=\ell$, proves
\eqref{ts:two-sided}. All expressions remain polynomial in the pair
errors at $D=0$; contour separation, not division by $D$, is the domain
condition here.
\end{proof}

Taking $\ell=(1,0)$ reads out the discriminant error. If paid enclosures
of the two right sides are intervals $I_L(\tau)$ and $I_R(\tau)$, then
\begin{equation}
 D(\tau)\in\widehat D(\tau)+
                  \bigl(I_L(\tau)\cap I_R(\tau)\bigr).
 \label{ts:intersection}
\end{equation}
This supplies a continuous counterpart to the discrete sweep: a sound
smaller state enclosure can be fed into the source residual and the
quadratic terms, and the resulting two readouts can be intersected again.
The intervals enclose two expressions for the same error. They are not
independent measurements. Disjoint intervals mean that the asserted
enclosures or their source assumptions are inconsistent; they do not
prove positivity of an excluded canonical path.

\subsection{Error payment, statistics and strict improvement}

For transparent sufficient error payment, suppose
$|e_D|\le E_D$, $|e_M|\le E_M$,
$|\delta\mathcal F|\le E_{\mathcal F}$ and $|r_i|\le R_i$.
Then
\begin{equation}
 |\Psi|\le
 4|p_D|E_DE_M+2|p_M|E_M^2+|p_M|E_{\mathcal F}
                    +|p_D|R_D+|p_M|R_M.
 \label{ts:absolute}
\end{equation}
This follows term by term from \eqref{ts:adjoint-def}. It is not a
claim that separate absolute allowances are efficient. The signed
correlation of the forcing error, reference defect and quadratic terms
may be essential. The complete source contribution to
$E_{\mathcal F}$ includes uncertain trial coordinates, with the kernel
derivatives of Section~\ref{sec:retained-forcing}, together with source
tails, contour/quadrature, arithmetic and whole-time remainders. Endpoint
and adjoint-solution errors must also be included when those objects are
computed rather than exact. Large adjoint coefficients can erase any
apparent improvement.

Fixed finite weighted sums of \eqref{ts:two-sided} give the corresponding
identity for linear empirical error statistics, with all labels and the
same source retained. They do not give a sharp bad-gap probability by
themselves. A bad-event majorant needs its own derivative/resolution
payment and complete population coverage; an integer zero-exception
argument needs a bound below the appropriate integer threshold. An
average over time can miss an isolated event. Whole-interval use of
\eqref{ts:intersection} therefore requires a uniform bound for every
$\tau$, paid interpolation on time cells, or a first-exit argument whose
other exits are independently excluded.

The two-sided construction distinguishes a possible route from an
acquired gain. It preserves the physical forcing discarded by the
initial-$M$ horizon and can preserve correlations lost by an unsigned
force norm. It has not yet supplied a strict canonical intersection
width, a better whole-time gap margin, a convergence rate or a lower
acquisition cost. In particular it does not acquire the forcing floor
used in Appendix~\ref{app:compact-example} merely by repeating a sweep.

Before a quantitative execution, one must fix a common finite pair or
block, physical interval, independently cleared source/count tube,
regular charts, endpoint observations, reference path and complete
signed residual budget. A proposed improvement must be strict and
measured against direct regular $D/M$ or source-discriminant evolution
with precisely the same information, including the elapsed forward
excursion. Finite work and refinement limits and a separate check belong
to that contract. The present section supplies the analytical interface,
not its missing source estimate.

Quadratic and cosine-based heat controls satisfy the local identities on
their valid charts. A sound enclosure whose input admits their collision
must retain that event. The canonical application uses the separate
source-equality, order/count, physical-readout and integer-multiplicity
premises summarized in Section~\ref{sec:discussion}; none is an automatic
favorable residual bound. Thus the recursive and continuous descriptions
are mathematically specified without being promoted to a proved recursive
improvement of Newman bounds.

\section{Interpretation and further analytical questions}
\label{sec:discussion}

The counted quadratic factor is the common starting point of the paper.
Its symmetric coordinates remain regular at a merger, and its rational
traces remove the need to follow individual exterior roots. Divided
differences then give the complete integer hierarchy. Split-pole moments
are algebraically equivalent coordinates, while the complex exponent
organizes a holomorphic family and complete damped contour formulas.
Neither change of representation alone supplies the signed information
needed by gap transport.

The recent consequences sharpen that distinction in two directions.
The finite-head residual is explicitly small in moment order, yet its
independent positive comparison can grow factorially. The appropriate
question is therefore whether compatibility or a finite-divisor
reconstruction yields a useful error theorem. Separately, the joint
forcing, weighted gap energy and deflated derivative expansion retain
more coupling than separate absolute bounds. Their exactness identifies
which additional estimate is needed, but does not supply it. The scalar
energy example concerns the two-equation consumer only, and the static
divisor example concerns the stated enlarged-constant slice class.
Neither excludes a better estimate from the full canonical dynamics.

In the all-real regime there is an additional elementary sign structure:
the complete exterior denominators of an adjacent pair are positive.
The backward comparison uses it to improve the local safety allowance
from initial neighboring gaps. Regional minima make extremal selection
well-defined without a global minimizing pair, and their weak evolution
handles switches of the minimizing label. The distribution nevertheless
couples to the exterior interaction. A joint law of consecutive gaps is
therefore a concrete statistical input, with a computable consequence,
rather than a substitute name for the desired collision exclusion.

Two analytical tasks remain distinct. First, a finite canonical product
could provide a compatible approximation on a common interval if its
complete error, local contour clearance and moment transfer are proved.
The polynomial identity in \eqref{fin:polynomial-heat} and the existing
slice-jet passage are useful ingredients, not the missing time-uniform
theorem. Second, a statistical descent requires an actual estimate on the
lower tail of the three-gap safety margin, together with an argument that
eliminates all exceptions over the interval and at unbounded heights.
A density-one conclusion or fixed positive-blur modulus alone does not
meet that requirement. For an isolated very narrow gap, the neighboring
comparison gives little relative improvement, making those configurations
a substantive part of the remaining problem.

The statistical use of gap evolution has an important precedent in
Rodgers and Tao \cite{RodgersTao2020Newman}. Their contradiction under a
negative Newman constant gives the lower bound $\Lambda\ge0$. The
proposed application here has the other direction: continue an established
all-real time downward to improve an upper bound. The local comparison
does not provide the complete initial statistics or uniform height
coverage needed for that conclusion. No improved value of $\Lambda$
follows from the present results.

The collision controls in the project background are consistent with
these scopes. Quadratic and quartic heat polynomials obey the local
algebra while lacking the canonical growing population. Source
perturbations with Euler translates, finite square weights or finite
prime constraints change the full reflected unit-weight theta source;
the documented endpoint-defective subclasses also fail the fixed strip.
Reflected super-Gaussian controls fail the order-one product premise.
Smooth positive and positive-primitive controls alter theta arithmetic
and do not establish the complete strip/count/safe-line inputs used
here. Atomic or finite-profile controls lack the full canonical
population or profile; changed observations and fractional profiles
change the physical readout or integer multiplicities. RH-conditional
endpoint controls retain RH as an input and change ordinary Dirichlet
arithmetic. The static surgery of Proposition~\ref{prop:static}
has no common canonical heat tube, and the scalar family of
Proposition~\ref{joint:scalar-screen} has no asserted divisor realization.
The detailed source and premise screen is retained in
\cite{BakerControls,BakerCMReview,BakerCM13,BakerCM17}.

All examples used for the conclusions of this draft are analytical.
The quartic controls are exact heat polynomials, the static construction
is an asymptotic entire-function family, and the truncation limitation
is an explicit path in a positive matrix exponential. Numerical work may
support a future specified estimate, but the proofs here do not depend
on samples or floating-point evidence. The unresolved questions concern
complete-divisor, common-time and uniform-height control.

\appendix
\section{Clearance and contour estimates}
\label{app:estimates}

We give explicit constants for the spatial and index contours in
Sections~\ref{sec:complex} and \ref{sec:cauchy}. The source inputs
are \eqref{eq:domain} and \eqref{eq:count-safe}
\cite{BakerBase,BakerGT41}. Write \(q_0=16/15\) and
\(\Lambda(x)=\log(3+2T+|x|)\).

\subsection{Complete value sums}\label{app:value-sums}

Let \(p=u+\ii v\), \(u>1/2\). The band \(|\Re w_j|<1\) has at
most \(4\Cu L\) occurrences. Each pair of unit shells
\(n\le|\Re w_j|<n+1\), \(n\ge1\), has at most
\(2\Cu\Lambda(n)\), by translating their unit intervals into
the physical variable. For \(\gamma>0\) and \(k=1,2\), put
\begin{equation}
 E_k(\gamma)=\sum_{j\ge0}2^{-j\gamma}(1+(j+1)\log 2)^k .
 \label{app:dyadic}
\end{equation}
For every integer \(M_0\ge1\) and real \(X\ge1\), dyadic blocks
give
\begin{align*}
 \sum_{n\ge M_0} n^{-2u}\Lambda(n)
    &\le M_0^{1-2u}\Lambda(M_0) E_1(2u-1),\\
 \int_X^\infty x^{-2u}\Lambda(x)^k\dd x
    &\le X^{1-2u}\Lambda(X)^k E_k(2u-1).
\end{align*}
For blocks starting at \(s=2^jM_0\) or \(2^jX\), bound the
number or length of the terms by \(s\), and use
\(\Lambda(2^{j+1}x)\le\Lambda(x)+(j+1)\log2\) for \(x\ge1\).

An occurrence of \(Z\) has modulus
\(|\mu_j|^u\exp(v\Im\Log Q_j)\).
The collar gives \(|\mu_j|\le q_0/a^2<1\), and on shell \(n\ge1\)
one also has \(|\mu_j|\le q_0\ell^2/n^2\). Thus for
\(1/2<b\le u\le r<\infty\),
\begin{equation}
\begin{split}
 |Z(p,t)|&\le A_b \exp(\Theta|v|),\\
 A_b&=4\Cu L(q_0/a^2)^b+
        2\Cu(\ell^2q_0)^b E_1(2b-1)\Lambda(1).
\end{split}
\label{app:value-envelope}
\end{equation}
The right side bounds the sum of moduli. The shell bound gives
normal convergence on compact subsets of \(u>1/2\) and uniform
tails over the real tube at fixed \(T\). An omitted occurrence
when \(N>4T+8\) satisfies \(|\Re w_j|>N-1-2T>N/2\) and has a
principal log. With \(M_0=\lfloor N/2\rfloor\ge N/3\), this gives
\begin{equation}
 |Z-Z_N|\le 2\Cu(\ell^2q_0)^u E_1(2b-1)3^{2r-1}
                N^{1-2u}\Lambda(N)\exp(\pi|v|/2).
 \label{app:divisor-tail}
\end{equation}
For an individual integer \(u=n\) use \(E_1(2n-1)3^{2n-1}\) in
this bound. The same constant without \(3^{2r-1}\) works for
a centered shell cutoff with all \(|\Re w_j|<N\) retained,
provided \(N\ge8\). Here \(\Lambda(N)=O(\log N)\) when \(N>4T+8\).

\subsection{Safe cuts and removal of the remote connectors}\label{app:clearance}

In a physical unit interval \([U,U+1]\), the zero projections
in \([U-1,U+2]\) number at most \(N_U=3\Cu\log(4+|U|)\).
Deleting the intervals of radius
\(\eta_U=1/[4(1+N_U)]\) about them removes length at most \(1/2\).
The remaining interval contains a cut \(v_0\) at distance at
least \(\eta_U\) from all projected zeros, since projections
outside the enlarged interval have distance at least one.
Translating the finite cut windows at the reference time gives
the cuts of Section~\ref{sec:complex}. Compactness of the
connectors and continuity of \(H(t,c(t)+w)\) preserve clearance
in a smaller time interval. For remote closures we need only
a sequence of cuts at each slice.

To estimate at such a cut put \(x=v_0-c(t)\) and
\(\lambda_0=\log(3+|v_0|)\). The once-subtracted product
logarithmic derivative gives
\[
 h(t,x+\ii y)-h(t,x+4\ii)
  =\sum_\rho\left(\frac1{v_0+\ii y-\rho}-\frac1{v_0+4\ii-\rho}\right).
\]
The regularizers cancel before summation. For \(|y|\le4\)
the near shell \(|\Re\rho-v_0|<1\) has
\(2\Cu\lambda_0\) as a count bound. Its denominator factors
have moduli at least \(\eta_U\) and \(3\), and the numerator is
at most \(8\). On shell \(n\ge1\) both factors are at least \(n\)
and the count is at most \(2\Cu(\lambda_0+\log(n+1))\).
The sums \(\sum n^{-2}\le2\) and
\(\sum n^{-2}\log(n+1)\le3\) bound the far part by
\(32\Cu\lambda_0+48\Cu\). Since
\(\eta_U^{-1}\le4+24\Cu\lambda_0\), the upper safe-line
bound and the near term \(16\Cu\lambda_0/(3\eta_U)\) give
\[
 |h(t,x+\ii y)|\le240\Cu^2\lambda_0^2 .
\]
For \(|x|\ge6,\ |y|\le4\) the principal log has phase at most
\(\pi/2\) and \(|Q|\ge x^2/2\), because \(|D|\le1/16\).
The absolute integral of \(f_ph\dd w\) on the remote connector,
before Cauchy normalization, is therefore at most
\begin{equation}
 1920 \Cu^2(2\ell^2)^u |x|^{-2u}
                      \Lambda(x)^2\exp(\pi|v|/2).
 \label{app:edge-bound}
\end{equation}
This tends to zero for the \(u\) in Theorem~\ref{thm:complex-time}.

\subsection{Time bounds for the spatial formula}\label{app:time-bounds}

On either horizontal safe line the reciprocal-square and twice
reciprocal-cube sums for \(h_w,h_{ww}\) are absolute. The divisor
distance is at least \(3\) in the near shell and at least \(n\)
in shell \(n\ge1\). Using the same count and elementary sums
as above, including the safe-line bound for \(h\), gives
\[
 |h|,\ |h_w|,\ |h_{ww}|\le A\Lambda(x),\qquad A=24\Cu.
\]
Take \(C_c=\sup_I|c'|\) and \(C_D=\sup_I|D'|\), finite by the
target contour. The inherited deflated-field and pair bounds
give \(C_c\le16 \Cu L^2/a+190 \Cu L\) and
\(C_D\le2+56d \Cu q_0 L/a^2\) \cite{BakerGT41,BakerCM03}.
Then \(|h_t|\le C_h\Lambda(x)^2\) with
\(C_h=A(1+C_c)+2A^2\). For \(J_p=1+|p|\) and \(|x|\ge6\),
\[
 \begin{aligned}
 |f_p|&\le (2\ell^2)^u|x|^{-2u}\exp(\pi|v|/2),\\
 |(f_p)_t|&\le 2C_D J_p|f_p|/x^2,\qquad
 |(f_p)_w|\le 8J_p|f_p|/|x|,\qquad
 |(f_p)_{ww}|\le 68J_p^2|f_p|/x^2.
 \end{aligned}
\]
In fact \((f_p)_w=-2pw f_p/Q\) and
\((f_p)_{ww}=(-2p/Q+4p(p+1)w^2/Q^2)f_p\).
The displayed bounds follow using \(|w|\le2|x|\) and
\(|Q|\ge x^2/2\). Either differentiated integrand of
\eqref{eq:complex-time} is bounded on these lines by
\(C_{\rm sl}J_p^2 |f_p|\Lambda(x)^2\), where
\[
 C_{\rm sl}=C_h+A(2C_D+68+8C_c)+8A^2 .
\]
For the value integrand use \(A|f_p|\Lambda(x)\).
The spatial primitive in the proof of
Theorem~\ref{thm:complex-time} tends to zero at the infinite
endpoints by these same estimates.

On a compact subchart let \(P_0\) be the total finite contour
length, \(W_0=\max|w|\) and \(\kappa=\min|Q|>0\) there.
Let \(A_0\ge1\) bound \(|h|,|h_w|,|h_{ww}|\) on these pieces
over the subchart. Clearance makes each constant finite. Set
\[
\begin{split}
 R_{b,r}&=\max_{q=b,r}(\ell^2/\kappa)^q,\\
 C_{\rm ch}&=A_0(1+C_c+C_D/\kappa+2/\kappa+
                     4W_0^2/\kappa^2+2C_cW_0/\kappa)
                        +(2+2W_0/\kappa)A_0^2 .
\end{split}
\]
The finite contribution to the value formula and each
differentiated formula is at most
\((P_0/(2\pi))R_{b,r} C_{\rm ch}J_p^2\exp(\Theta|v|)\),
by the explicit derivatives of \(f_p,h\).
Bounding the four full rays by the integrals over \(|x|\ge6\)
then gives
\begin{equation}
\begin{split}
 |Z_t| &\le C_Z J_p^2\exp(\Theta|v|),\\
 C_Z &=\frac{P_0 R_{b,r} C_{\rm ch}}{2\pi}
   +\frac{2 C_{\rm sl}}{\pi}\max_{q=b,r}((2\ell^2)^q6^{1-2q})
                          E_2(2b-1)\Lambda(6)^2 .
\end{split}
\label{app:chart-time}
\end{equation}
The maximum arises because the scale factor is exponential in
the real part \(u\). On the fixed finite pieces the integrands are
\(C^1\) in physical time by clearance. On each safe-line point
they are continuous in time and the displayed estimates
give an integrable majorant on compact \(p\)-sets and compact
time subcharts. This proves the asserted differentiation,
continuity and holomorphy in Theorem~\ref{thm:complex-time}.
The finite chart constants concern fixed physical height.

For \(X\ge7\) let \(Z^{[X]}=\mathcal S^{[X]}(f_ph)\) with the
finite pieces retained and the four rays stopped at \(|x|=X\).
The dyadic integral estimates and \(2\ell^2<1\) give, on
\(b\le u\le r\),
\begin{equation}
\begin{aligned}
 |Z-Z^{[X]}|&\le U_0(X)\exp(\pi|v|/2),\\
 |Z_t-Z_t^{[X]}|&\le U_1(X)J_p^2\exp(\pi|v|/2),\\
 U_0(X)&=\frac{2A}{\pi}(2\ell^2)^b E_1(2b-1)
                                X^{1-2b}\Lambda(X),\\
 U_1(X)&=\frac{2C_{\rm sl}}{\pi}(2\ell^2)^b E_2(2b-1)
                                X^{1-2b}\Lambda(X)^2 .
\end{aligned}
\label{app:spatial-tails}
\end{equation}
The derivative of the cut value uses the unintegrated
derivative in \eqref{eq:complex-time}. The factor \(2/\pi\)
accounts for the four rays. The bounds hold with \(u\) in
the scale and power of \(X\) in place of \(b\), since the
dyadic constants are decreasing in \(2u-1\). This gives the
fixed-height \(X^{1-2u}\log X\) and \(X^{1-2u}\log^2 X\) rates
for \(X>4T+8\).

\subsection{Errors on the damped index contours}\label{app:cauchy-errors}

For a fixed strip in Section~\ref{sec:cauchy}, set
\[
 B_N=2\Cu(\ell^2q_0)^b E_1(2b-1)3^{2r-1}
                                           N^{1-2b}\Lambda(N).
\]
Then the divisor error is bounded by \(B_N\exp(\pi|v|/2)\)
throughout the strip. On the complete vertical difference,
\begin{equation}
 |\mathcal V_{b,r}(K(Z-Z_N))|\le
   \frac{2 C_\epsilon B_N}{\pi\lambda}
                   \sqrt{\frac{\pi}{\epsilon}}\exp(\pi^2/(16\epsilon)).
 \label{app:weighted-divisor-tail}
\end{equation}
Indeed \(\sqrt{\lambda^2+v^2}\ge\lambda\), and translating the
Gaussian integral after completing the square pays the factor
\(\exp(\pi^2/(16\epsilon))\).
This rate concerns the whole strip. Evaluating the signed residue
values gives the sharper \eqref{cau:residue-error}.

Write \(U=1+r\) so that \(J_p\le U+|v|\), and set
\(C_{\epsilon,t}=16C_D\ell^{-2}\exp(4\epsilon)\) for \(K_\epsilon\)
and \(C_{\epsilon,t}=0\) for \(G_\epsilon\). Formula
\eqref{cau:kernel-time} gives the same kernel bounds for \(K_t\)
with this constant in place of \(C_\epsilon\). The total absolute
bound for the differentiated horizontal integrals is
\[
 \frac{r-b}{\pi Y}\big[C_\epsilon C_Z(U+Y)^2+
                         C_{\epsilon,t}A_b\big]\Phi_\Theta(Y).
\]
For \(Y>\Theta/(2\epsilon)\) the differentiated vertical tail is
bounded by
\[
 \frac{2\Phi_\Theta(Y)}{\pi}\left[
 C_\epsilon C_Z(1+U/Y)^2
     \left(\frac{Y}{\kappa_\Theta(Y)}+\frac{1}{\kappa_\Theta(Y)^2}\right)
 +\frac{C_{\epsilon,t}A_b}{Y\kappa_\Theta(Y)}
 \right].
\]
Use \((U+v)^2/v\le(1+U/Y)^2v\) for \(v\ge Y\) and integrate
against \(\exp(-\kappa_\Theta(Y)(v-Y))\). This pays both the
complete derivative and the derivative of the finite rectangle
as its height grows.

Finally \eqref{app:spatial-tails} bounds the complete vertical
derivative error due to cutting the independent spatial rays by
\begin{equation}
 \frac{2}{\pi\lambda}\int_0^\infty
 \big[C_\epsilon U_1(X)(U+v)^2+C_{\epsilon,t}U_0(X)\big]
               \exp(-\epsilon v^2+\pi v/2)\dd v .
 \label{app:weighted-spatial-tail}
\end{equation}
For the value error replace the square bracket by
\(C_\epsilon U_0(X)\). All Gaussian moments are finite at the
stated fixed damping. A finite maximum of the compact chart
constants suffices for the established real tube at each fixed \(T\).
The estimates use real physical time and fixed residue indices.


% Supporting inherited certificate only; not an input to the main theory.
\section{An inherited compact certificate as a supporting example}
\label{app:compact-example}

This appendix applies the analytical screen of
Proposition~\ref{cmp:small-slice} to an existing certificate. The source
values are numerical supporting evidence, already obtained with complete
source, interpolation, quadrature and arithmetic allowances in
\cite{BakerGT15,BakerGT14}; no new source evaluation is performed here.
The main contour, lattice and comparison theorems do not depend on them.
The resulting arithmetic and scope were checked in \cite{BakerCM27}.

The certificate follows the Lehmer pair near $a=14010.1$ on
$0\le t\le0.1$. The gauged field $G=\e^{\pi a/8}H$ counts two zeros
inside $|z-a|<3/4$ and has a cleared annulus out to radius $21/20$.
Its centered annular logarithm has modulus less than $1.737$ on both
boundaries and hence throughout the annulus. Its directly enclosed
incoming values at $t_*=0.1$ imply the safe rounded bounds
\begin{equation}
 \begin{gathered}
 0.1787368\le D(t_*)\le0.1787369,\qquad
 0.6037742\le M(t_*)\le0.6037744,\\
 V(0)\le1-2(0.1787368)(0.6037742)
       =0.78416666313888.
 \end{gathered}
 \label{supp:incoming}
\end{equation}
The older certificate's drift symbol is $K_{\mathrm{GT}}=-M$; it is
not the forcing $K=2D^2\mathcal F$ used in the present paper.
Here $V(0)$ is at backward time zero, namely physical time $t_*$.

Consider the proposed local descent to $t=-0.0003$, of duration
$\Delta=0.1003$. With $d_{\mathrm{stop}}=0.0001$, a hypothetical
$v\le0.89$ would leave the strict margin
\[
 0.1787368-2(0.89)(0.1003)-0.0001=0.0001028>0,
\]
and would require the forcing floor $\kappa\ge1-0.89^2=0.2079$.
The existing certificate ends at physical time zero: neither the
negative-time source bounds nor the proposed floor are available from it.
Even the favorable initial-$M$ scalar allowance is at most
\[
 \frac{0.1787369}{2(1-0.1787369\cdot0.6037744)}<0.100180,
\]
which does not reach the proposed duration.

The already certified physical slice $t=0$ is enough to test the stronger
constant-$v$ proposal. At that slice,
\[
 |c-a|<0.1,\qquad 0<D<0.002,\qquad
 |\mathcal L-\mathcal L_{\mathrm{center}}|<7/4
 \quad\text{on }|z-a|=0.9.
\]
There $|Q|>0.638>3/5$ and $|Q'|\le2$. The complete log kernel
\eqref{cmp:log-kernel} gives
\begin{align*}
 |S_2|&\le 2(0.9)(7/4)2/(3/5)^3=175/6,\\
 |S_3|&\le 3(0.9)(7/4)2/(3/5)^4=875/12,\\
 |\mathcal F|&\le12|S_2|+16(0.002)|S_3|\le1057/3.
\end{align*}
Consequently the actual complete forcing satisfies
\begin{equation}
 |K(t=0)|<2(0.002)^2(1057/3)
        =1057/375000<0.003.
 \label{supp:small-K}
\end{equation}
This upper bound is not a failed lower enclosure. Any constant floor
covering this slice has $\kappa<0.003$, and the selected constant
parameter must obey $v>\sqrt{0.997}>0.998$. Even exact initial data and
zero stop threshold therefore give only
\[
 \Delta<\frac{0.1787369}{2(0.998)}<0.0896<0.1003.
\]
Refining the numerical floor cannot make this particular consumer pass.
The slice is inside the proposed stopping region, since its certified
$D(t=0)>0.00142117>0.0001$. It is used to reject a new allowance, not to
present already known positivity as new coverage.

The conclusion concerns the one constant-$v$ barrier from this fixed
anchor over an interval containing physical time zero. It proves no
collision or failure of actual descent, and excludes neither a sharper
transient comparison nor time-dependent or regular coupled evolution.
The same-source direct endpoint estimate is stronger on the old interval;
no method is certified here on the proposed negative-time extension.
No global all-real start, practical all-height cutoff or improved bound
on $\Lambda$ follows from this local supporting example.

\subsection{An unsigned force envelope at the same anchor}
\label{supp:force-envelope}

The subsequent comparison in \cite{BakerCM32,BakerCMReview29} uses
the same anchor to test a specific compression of the common node-force
equation. For two nodes put $r_U=|PU|$, where $P$ subtracts the mean
and $U_i=H_{zz}(x_i)/(2H_z(x_i))$. The exact common-label identity gives
\[
 g(u)=g(0)-2\int_0^u(U_1-U_0)\dd s
       \ge g(0)-2\sqrt2\int_0^u r_U\dd s.
\]
In backward time the internal graph term is antidissipative. Replacing
its gap denominator by the fixed stop
$s=2\sqrt{d_{\rm stop}}=1/50$ and the exterior forcing by an
independent nonnegative norm allowance gives the particular scalar
envelope
\[
 \mathcal R'=10000\mathcal R+J_{\rm ext},\qquad
 J_{\rm ext}\ge0,\qquad \mathcal R(0)\ge r_U(0).
\]
Even with zero exterior allowance and zero added errors,
$\mathcal R(u)\ge r_U(0)\e^{10000u}$.
This is a property of the envelope, not a lower growth estimate for
the actual force norm.

The safe incoming bounds in \eqref{supp:incoming} imply
$g(0)<423/500$ and $M(0)<151/250$. The edge force identity therefore
gives
\[
 U_1-U_0=\frac2{g(0)}-g(0)M(0)
  >\frac{1000}{423}-\frac{423}{500}\frac{151}{250}
  >\frac{37}{20}.
\]
Since $2\sqrt2r_U(0)=2|U_1-U_0|>37/10$, a positive margin from this
unsigned envelope requires
\begin{equation}
 \Delta<10^{-4}\log\left(1+
       \frac{10000(0.846-0.02)}{3.7}\right)
       =10^{-4}\log\frac{82637}{37}<0.0008.
 \label{supp:force-duration}
\end{equation}
The last strict comparison follows exactly from
\[
 \e^8>\sum_{k=0}^{10}\frac{8^k}{k!}
       =\frac{3830591}{1575}>\frac{82637}{37}.
\]
This rejects the fixed-stop-denominator, unsigned-exterior scalar
compression for the proposed $0.1003$ descent. It does not reject
state-dependent denominators, signed exterior correlations, other blocks
or the complete force dynamics.

\subsection{The retained constant-\texorpdfstring{$K$}{K} transient still misses the endpoint}
\label{supp:constant-K-transient}

Theorem~\ref{rf:constant-K} retains more than the earlier constant-$v$
barrier. It can nevertheless be screened using the same actual included
slice \eqref{supp:small-K}. Give its scalar model the favorable data
\[
 d_+=0.1787369,\qquad b_+=0.6037744,\qquad
 v_-=1-2d_+b_+,\qquad \kappa_*=1057/375000.
\]
The initial gap $2\sqrt D$ and initial backward velocity
$-2/\sqrt D+4M\sqrt D$ are both increasing on the positive anchor
rectangle. The stated upper endpoints therefore maximize both initial
data. Increasing $\kappa$ likewise favors the cooperative comparison.
Any constant floor covering the certified $t=0$ slice is less than
$\kappa_*$, so these choices are favorable even to an exact-data use.

Substitution of the rational values into
\eqref{rf:constant-K-polynomial} gives the exact inequality
\begin{equation}
 d_+-2v_-\Delta+
       \frac{v_-^2+\kappa_*-1}{d_+}\Delta^2
       <-\frac{8239}{10^8}<0,
 \qquad \Delta=1003/10000.
 \label{supp:transient-failure}
\end{equation}
The model starts positive and has reached its first zero before this
endpoint. Its negative polynomial value is not asserted as a
positive-gap comparison after that zero. Equivalently, even the
favorable endpoint requirement $d_\kappa(\Delta)>d_{\rm stop}$ would
need
\[
 \kappa>
 1-v_-^2+\frac{d_+}{\Delta^2}
       (d_{\rm stop}-d_++2v_-\Delta)
 >0.00605917>\kappa_*.
\]
For a stopped-trial use, successful passage would put the actual included
slice inside the trial region. If a lower stop occurred earlier, the
proposed target would already have failed. The necessary screen therefore
does not assume the desired whole-interval positivity.

This strengthens the earlier allowance rejection to the sharp scalar
constant-$K$ transient comparison from this anchor. It still excludes
neither time-dependent $K$ nor a constant physical $\mathcal F$ floor.

\subsection{A positive conditional margin for retained physical forcing}
\label{supp:retained-F}

Set
\[
 d=0.1787368,\qquad b=0.6037742,\qquad f=3/10,\qquad
 \Delta=1003/10000.
\]
These are the incoming lower data and a proposed complete forcing floor,
not a measured lower value of $\mathcal F$. The model of
Theorem~\ref{rf:forcing-comparison} has a rigorously positive margin.
Here is the exact rational calculation, as recorded in
\cite{BakerCM32,BakerCMReview29}.

Put $\lambda^2=2f=3/5$, $w=4\lambda^2\Delta^2$, $v=1-2db>0$ and
\[
 J=\frac d2+\frac{2db^2-2b}{\lambda^2}<0.
\]
Expanding \eqref{rf:physical-forcing-model} by double-angle identities
gives
\begin{equation}
 d_f(\Delta)=d+J[\cosh(2\lambda\Delta)-1]
                    -v\frac{\sinh(2\lambda\Delta)}{\lambda}.
 \label{supp:F-expansion}
\end{equation}
The positive Taylor series have the rational upper bounds
\begin{align*}
 C_+&=\sum_{j=0}^4\frac{w^j}{(2j)!}
       +\frac{w^5/10!}{1-w/(12\cdot11)},\\
 S_+&=2\Delta\left[
        \sum_{j=0}^4\frac{w^j}{(2j+1)!}
        +\frac{w^5/11!}{1-w/(13\cdot12)}\right].
\end{align*}
Indeed, after the first omitted terms the successive ratios are at
most $w/(12\cdot11)$ and $w/(13\cdot12)$, respectively; both are
strictly below one. Thus
$\cosh(2\lambda\Delta)\le C_+$ and
$\sinh(2\lambda\Delta)/\lambda\le S_+$.
The two coefficients $J$ and $-v$ in \eqref{supp:F-expansion} are
negative, so exact rational substitution proves
\begin{equation}
 d_f(\Delta)\ge d+J(C_+-1)-vS_+
       >\frac{16246}{10^8},\qquad
 d_f(\Delta)-d_{\rm stop}>\frac{6246}{10^8}.
 \label{supp:F-margin}
\end{equation}
All quantities in this comparison are rational; no sampling of the
canonical source or fit to a time profile enters it. The inequalities
$b>\sqrt{f/2}$ and $4bd<2$ hold for these data. Corollary
\ref{rf:forcing-stop} therefore makes the endpoint the model's minimum
on the full interval, so the second inequality is the strict lower-stop
margin, conditional on the actual source hypotheses.

For the same trial geometry $r=9/10$, $C=1/10$, $D_{\max}=1/5$ and
$q=11/25$, \eqref{rf:source-errors} gives the exact amplification factors
\[
 E_{\mathcal F,B}\le\frac{119250}{1331}e_B,
 \qquad
 E_{\mathcal F,\mathcal L}\le\frac{14175000}{14641}e_{\mathcal L}.
\]
The complete source/time error, the signed lower floor
$\mathcal F_\Gamma\ge3/10$, the negative-time contour clearance and
the unconditional coordinate bounds are unacquired. The annular
logarithm bound used earlier provides an absolute forcing bound, not
this signed lower floor. Small $K=2D^2\mathcal F$ at $t=0$ is compatible
with this proposed physical floor, but does not prove it. Direct regular
$D/M$ evolution with identical data and errors gives precisely the
same model and margin. Thus \eqref{supp:F-margin} is a conditional
consumer calculation, not a new canonical certificate or a lattice
advantage.

\subsection{A distinct time-zero census and its forward constants}
\label{fwd:census-example}

The finite census used in \cite{BakerCM34,BakerCMReview30} lies at a
different height from the Lehmer pair. Its constants are not transferred
to that pair. The accepted source record \cite{BakerOCCensus22} encloses
every zero in
\[
 \mathcal R=(988,2012)+\ii(-3/2,3/2)
\]
in 388 occupied cells, each of count one including multiplicity.
Their horizontal projections are disjoint, their maximum horizontal
width is $5/32768$, and successive projections have separation at
least $5085/8192$. The common recorded imaginary interval is
$[-3/32768,3/131072]$, which is not symmetric about zero.

Reality follows from the complete census and conjugation, rather than
from replacing these imaginary intervals by zero. The conjugate of a
root in $\mathcal R$ is another root in the same symmetric outer
rectangle, with the same real coordinate. Completeness puts it in an
occupied cell; disjoint horizontal projections force that cell to be
the original one. Its count one forces equality of the two roots and
simple multiplicity. There are consequently 388 real simple nodes,
387 internal consecutive gaps and 385 edges with both neighboring gaps
inside the census. This deduction retains the inherited winding
certificate's scope; it is not an independent replay of its source
arithmetic.

The initial squared-half-gap minimum and the finite-block coefficient are
\begin{equation}
 q_0=\frac{5085^2}{4\cdot8192^2}
       =\frac{25857225}{268435456},\qquad
 a_{388}=\frac{388}{194^2}=\frac1{97}.
 \label{supp:census-constants}
\end{equation}
If a future common tube keeps the block diameter below $1024$, the
force-dispersion contraction coefficient $2N/L_{\max}^2$ is
$\gamma=97/131072$. Even with zero exterior norm forcing, its
homogeneous envelope over the illustrative time $s=1/10$ decreases by
a relative amount
\[
 1-\e^{-\gamma s}<\gamma s=\frac{97}{1310720},
\]
about $0.0074$ percent. This is the reduction in that conservative
envelope, not an upper bound on actual contraction. Continued
confinement and zero exterior forcing have not been acquired.

For a paid uniform upper exterior interaction $C\ge0$, the forward
minimum envelope is
\[
 q(s)\ge\e^{-4Cs}q_0+
          \frac{1-\e^{-4Cs}}{388C},
\]
with limit $q_0+s/97$ when $C=0$. Its lower bound exceeds $q_0$ exactly
when
\begin{equation}
 C<\frac{a_{388}}{4q_0}
   =\frac{268435456}{10032603300}.
 \label{supp:census-exterior-target}
\end{equation}
No such whole-time exterior bound is supplied by the census. Both block
ends must be included in its acquisition: the 12-unit buffers of the
smaller observation window $[1000,2000]$ cannot be assigned to the
outermost retained nodes. Likewise, selecting only interior eligible
edges does not establish the cumulative-distance inequalities for the
minimum over the complete block.

This finite artifact supplies actual time-zero labels and gap geometry,
not a positive-time source tube, an exterior-force bound, or an all-height
gap law. The analytical forward comparisons remain conditional on their
independent source and boundary payments, and Theorem~\ref{rf:horizon}
still applies to any subsequent use of its specific backward consumer.

\bibliographystyle{plainurl}
\bibliography{references,induction}
\end{document}
